\hypertarget{section}{%
\section{}\label{section}}

\learningobjectives{

1. Web
2. HTML5CSS3JavaScript
3. Vue.js
4. 
5. 
}

\hypertarget{section-1}{%
\subsection{}\label{section-1}}

Web

WebHTMLJavaScriptjQueryAngularReactVueVue.js

HTMLCSSJavaScriptGIS

\hypertarget{section-2}{%
\subsection{}\label{section-2}}

!!! warning ``\,''

\begin{verbatim}
2024
\end{verbatim}

!!! info ``\,''

\begin{verbatim}
- - 
- PC
- - 
\end{verbatim}

\hypertarget{section-3}{%
\subsection{}\label{section-3}}

!!! info ``\,''

\begin{verbatim}
### [ ](section04-01.md)
- Web
- W3CWeb
- JavaScript
- ### [ HTML](section04-02.md)
- HTML
- HTML5
- - 
- SEO

### [ CSS](section04-03.md)
- CSS
- CSS3
- FlexboxGrid
- - UI

### [ JavaScript](section04-04.md)
- JavaScript 
- ES6+
- Promise
- DOM
- ### [ Vue](section04-05.md)
- - Vue.jsMVVM
- Vue 
- - 

### [ ](section04-06.md)
- Vue CLI
- WebpackVite
- - 

### [ ](section04-07.md)
- - CDN
- - 
\end{verbatim}

\keypoints{
||  ||
|------|------|-------------------|
|| UI ||
||  ||
|(SPA)|  ||
|DOM| UI ||
||  ||
}

\hypertarget{section-4}{%
\subsection{}\label{section-4}}

!!! note ``\,''

\begin{verbatim}
\important{}
- HTML5: 
- CSS3: 
- JavaScript ES6+: 

\important{}
- Vue.js 3.x: API
- Vue Router: 
- Vuex/Pinia: 

\important{}
- Vue CLI / Vite: 
- ESLint: 
- Prettier: 

\important{UI}
- Element Plus: Vue 3
- Ant Design Vue: UI
- Vant: Vue
\end{verbatim}

\hypertarget{section-5}{%
\subsection{}\label{section-5}}

!!! example ``\,''

\begin{verbatim}
\important{}
- - 
- - 

\important{}
- GIS
- - 
- \important{}
- WebSocket
- - 
- \important{}
- - 
- - 
\end{verbatim}

\hypertarget{section-6}{%
\subsection{}\label{section-6}}

!!! tip ``\,''

\begin{verbatim}
\important{1-2}
1. Web
2. HTML5
3. CSS3
4. JavaScript ES6+

\important{2-3}
1. Vue.js
2. 
3. 
4. 

\important{1-2}
1. 
2. 
3. 
4. 

\important{2-3}
1. 
2. 
3. 
4. 
\end{verbatim}

\hypertarget{section-7}{%
\subsection{}\label{section-7}}

``\,''

\important{}

\begin{itemize}
\item
  HTML/CSS
\item
  \important{}
\item
  JavaScript
\item
  \important{Vue}
\item
  Vue
\item
  \important{}
\item
  \begin{itemize}
  \tightlist
  \item
  \end{itemize}
\end{itemize}

\hypertarget{section-8}{%
\subsection{}\label{section-8}}

{[}1{]} World Wide Web Consortium. Web Content Accessibility Guidelines
(WCAG) 2.1{[}S{]}. 2018.

{[}2{]} . GB/T 25000.51-2016 SQuaRECOTS{[}S{]}. : , 2016.

{[}3{]} Mozilla Developer Network. HTML Living Standard{[}EB/OL{]}.
{[}2024-01-15{]}. https://developer.mozilla.org/en-US/docs/Web/HTML.

{[}4{]} Vue.js Team. Vue.js Guide{[}EB/OL{]}. {[}2024-01-15{]}.
https://vuejs.org/guide/.

{[}5{]} . {[}R{]}. : , 2023.

\hypertarget{section-9}{%
\subsection{}\label{section-9}}

\begin{enumerate}
\def\labelenumi{\arabic{enumi}.}
\item
  \important{}WebVue.js
\item
  \important{}
\item
  \important{}
\item
  \important{}
\end{enumerate}

\hypertarget{web}{%
\section{4.1.1 Web}\label{web}}

WebWebWeb

\hypertarget{web-1}{%
\subsection{Web}\label{web-1}}

WebFront-endWebWebHTMLCSSJavaScript

WebGIS

\important{} Web

WebWindowsmacOSLinuxWeb

\hypertarget{web-2}{%
\subsection{Web}\label{web-2}}

World Wide WebWebHTTPWeb

WebURLHTTPHTMLCSSJavaScriptWeb API

WebHTTPHTMLCSSJavaScript

\important{} WebW3CWebHTML5CSS3DOM

\hypertarget{section-10}{%
\subsection{}\label{section-10}}

WebHTMLCSSJavaScript

User InterfaceBrowser EngineRendering EngineNetworkingJavaScriptUIUI
BackendData StorageHTMLCSSHTTPJavaScriptUICookielocalStorage

DOMCSSJavaScript

\hypertarget{section-11}{%
\subsection{}\label{section-11}}

Rendering EngineHTMLCSSChromeEdgeBlinkSafariWebKitFirefoxGecko

HTMLDOMDocument Object ModelHTMLCSSCSSOMCSS Object
ModelDOMDOMCSSOMRender TreeLayoutPaint

DOMDOMDOMDOM

\important{} GPUGIS

\hypertarget{javascript}{%
\subsection{JavaScript}\label{javascript}}

JavaScriptJavaScriptWebJavaScriptChromeV8FirefoxSpiderMonkeySafariJavaScriptCoreECMAScriptJavaScript

JavaScriptJavaScriptJITJavaScript

JavaScriptJavaScriptJavaScriptWebSocketJavaScript

JavaScriptJavaScriptPromiseasync/awaitWeb Workers

\hypertarget{section-12}{%
\subsection{}\label{section-12}}

WebWeb

CSSCSS3JavaScript APIWeb APIHTMLCSS

CSSJavaScript polyfill

\important{} WebSocket API

\hypertarget{section-13}{%
\subsection{}\label{section-13}}

WebXMLHttpRequestFetch APIWebSocketServer-Sent Events

HTTPRESTful APIWebSocketServer-Sent Events

JSONWebGISGeoJSON

HTTP/2

WebHTML5CSS3JavaScriptW3CWeb

\hypertarget{html}{%
\section{4.2 HTML}\label{html}}

HTMLHyperText Markup LanguageWebWebHTMLCSSJavaScriptHTML

HTML5HTMLWebHTML5HTML5

!!! info ``HTML''

\begin{verbatim}
HTML5HTMLHTML
\end{verbatim}

\hypertarget{html-1}{%
\subsection{HTML}\label{html-1}}

\hypertarget{html-2}{%
\subsubsection{HTML}\label{html-2}}

HTMLHyperText Markup Language\important{}\important{elements}

HTML - \important{HyperText} - \important{Markup} - \important{Language}

\hypertarget{html-3}{%
\subsubsection{HTML}\label{html-3}}

HTML

\begin{lstlisting}[style=xml]
<!DOCTYPE html>
<html lang="zh-CN">
<head>
    <meta charset="UTF-8">
    <meta name="viewport" content="width=device-width, initial-scale=1.0">
    <title></title>
</head>
<body>
    <h1></h1>
    <p></p>
</body>
</html>
\end{lstlisting}

\important{}

\begin{itemize}
\item
  \code{<!DOCTYPE html>}HTML5
\item
  \code{<html>}
\item
  \code{<head>}
\item
  \code{<body>}
\end{itemize}

\hypertarget{html-4}{%
\subsubsection{HTML}\label{html-4}}

HTML\important{}\important{}\important{}\important{}

\begin{lstlisting}
<tagname></tagname>
\end{lstlisting}

\begin{lstlisting}[style=xml]
<h1></h1>
<p></p>
\end{lstlisting}

\important{}

\begin{enumerate}
\def\labelenumi{\arabic{enumi}.}
\item
  \important{}

  \textbackslash begin\{lstlisting}{[}style=xml{]}

  ``\textbackslash code\{
\item
  \important{}

  }\texttt{\textbackslash{}code\{html\ \ \ \ \textless{}img\ src="logo.jpg"\ alt=""\textgreater{}\ \ \ \ \textless{}br\textgreater{}\ \ \ \ \textless{}hr\textgreater{}\ \ \ \ }}\textbackslash code\{
\item
  \important{}

  \begin{itemize}
  \tightlist
  \item
    }

    \code{, }

    \code{, }

    -

    \code{, }

    \code{, }

    \code{, }

    \textbackslash code\{
  \end{itemize}
\item
  \important{}

  \begin{itemize}
  \tightlist
  \item
    }\code{, }\code{, }\code{, }\code{, }\textbackslash code\{
  \end{itemize}
\end{enumerate}

\hypertarget{html-5}{%
\subsubsection{HTML}\label{html-5}}

HTML\important{attributes}

\textbackslash end\{lstlisting}html

\begin{lstlisting}
\important{}
- }id\code{
- }class\code{CSSJavaScript
- }title\code{
- }lang\code{
- }style\code{CSS

### HTML

1. \important{}
   }`\code{html
   <P></P>  <!--  -->
   <p></p>    <!--  -->
   }`\code{

2. \important{}
   }`\code{html
   <img src="image.jpg" alt="">
   }`\code{

3. \important{}
   }`\code{html
   <!--  -->
   <p><strong></strong></p>
   
   <!--  -->
   <p><strong></p></strong>
   }`\code{

4. \important{}
   }`\code{html
   <br />
   <img src="image.jpg" alt="" />
   }`\code{

HTMLHTML5

## 4.2.1 HTML5

HTML5WebHTMLHTML5HTML

SemanticHTML}<header>\code{}<nav>\code{}<article>\code{}<section>\code{HTMLSEO

}<header>\code{}<nav>\code{}<main>\code{}<section>\code{}<article>\code{}<aside>\code{

\important{} HTML5

||  ||  |
|--------|----------|----------|------------------|
|}<header>\code{|  ||  |
|}<nav>\code{|  ||  |
|}<main>\code{|  ||  |
|}<section>\code{|  ||  |
|}<article>\code{|  ||  |
|}<aside>\code{|  ||  |
|}<footer>\code{|  ||  |
|}<figure>\code{|  ||  |
|}<figcaption>\code{|  ||  |
|}<time>\code{|  ||  |
|}<mark>\code{|  ||  |

### 

#### 1. }<header>\code{  - 

}<header>\code{HTML5}<div>\code{}<header>\code{

}<header>\code{}<header>\code{}<header>\code{}<address>\code{}<footer>\code{}<header>\code{

}<header>\code{}<header>\code{}<header>\code{

\end{lstlisting}

html

• •

41001500 \textbar{} 14:30

\begin{lstlisting}
#### 2. }<nav>\code{  - 

}<nav>\code{HTML5

}<nav>\code{}<nav>\code{

}<nav>\code{}aria-label\code{}aria-labelledby\code{

\end{lstlisting}

html

\textgreater{} \textgreater{} {}

\begin{lstlisting}
#### 3. }<main>\code{  - 

}<main>\code{HTML5HTML}<main>\code{}<article>\code{}<aside>\code{}<footer>\code{}<header>\code{}<nav>\code{

}<main>\code{}<main>\code{}<main>\code{

Web}<main>\code{}<section>\code{}<article>\code{

\end{lstlisting}

html

24

\begin{lstlisting}
#### 4. }<section>\code{  - 

}<section>\code{}<div>\code{}<section>\code{

}<section>\code{sectionh1-h6}<div>\code{

}<section>\code{sectionsection

\end{lstlisting}

html

85.23

86.00

88.50

2150/

1980/

\begin{lstlisting}
#### 5. }<article>\code{  - 

}<article>\code{

}<article>\code{}<section>\code{}<article>\code{}<section>\code{RSS}<article>\code{

}<article>\code{

\end{lstlisting}

html

2024-03-15 12:00

650-80

400-1234-5678

\begin{lstlisting}
#### 6. }<aside>\code{  - 

}<aside>\code{

}<aside>\code{

}<aside>\code{

\end{lstlisting}

html

±0.01

5

\begin{lstlisting}
#### 7. }<footer>\code{  - 

}<footer>\code{}<header>\code{}<footer>\code{

}<footer>\code{}<footer>\code{

}<footer>\code{

\end{lstlisting}

html

© 2024

15\%\ldots{}

\begin{lstlisting}
#### 8. }<figure>\code{  }<figcaption>\code{  - 

}<figure>\code{}<figcaption>\code{}<figure>\code{

}<figure>\code{}<figcaption>\code{}<figure>\code{

\end{lstlisting}

html

124 \textbar{} 2024-03-15 14:30

2024315 \textbar{}

\begin{lstlisting}
#### 9. }<time>\code{  - 

}<time>\code{HTML5}datetime\code{

}<time>\code{}datetime\code{ISO 8601"YYYY-MM-DD""YYYY-MM-DDTHH:MM:SS"}datetime\code{

}<time>\code{

\end{lstlisting}

html

2024315 14:30

2024315

2

\begin{lstlisting}
#### 10. }<mark>\code{  - 

}<mark>\code{}<strong>\code{}<em>\code{}<mark>\code{

}<mark>\code{CSS}<mark>\code{

}<mark>\code{

\end{lstlisting}

html

85.2386.00

``''10

\begin{lstlisting}

}<section>\code{}<header>\code{}<figure>\code{}<figcaption>\code{}<table>\code{}<footer>\code{

## 4.2.2 HTML5

WebHTML5

HTMLURLJavaScriptHTML5Web

HTML5

||  ||  |
|----------|----------|----------|-------------|
|}email\code{|  |}required\code{, }placeholder\code{|  |
|}url\code{|  |}required\code{, }placeholder\code{|  |
|}number\code{|  |}min\code{, }max\code{, }step\code{|  |
|}range\code{|  |}min\code{, }max\code{, }step\code{, }value\code{|  |
|}date\code{|  |}min\code{, }max\code{, }value\code{|  |
|}time\code{|  |}min\code{, }max\code{, }step\code{|  |
|}datetime-local\code{|  |}min\code{, }max\code{, }step\code{|  |
|}color\code{|  |}value\code{|  |
|}search\code{|  |}placeholder\code{, }results\code{|  |

### HTML5

#### 1. }email\code{  - 

}email\code{@

@

\end{lstlisting}

html

\begin{lstlisting}
#### 2. }url\code{  - 

}url\code{URLURLURL}email\code{URL

URLhttp://https://ftp://URL

URL

\end{lstlisting}

html

\begin{lstlisting}
#### 3. }number\code{  - 

}number\code{

}min\code{}max\code{}step\code{}value\code{

\end{lstlisting}

html ()

(m³/s)

\begin{lstlisting}
#### 4. }range\code{  - 

}range\code{

}range\code{}number\code{}min\code{}max\code{}step\code{}value\code{}number\code{}range\code{}<output>\code{JavaScript

\end{lstlisting}

html

90

\begin{lstlisting}
#### 5. }date\code{  - 

}date\code{

}min\code{}max\code{}value\code{ISO 8601YYYY-MM-DD

\end{lstlisting}

html

\begin{lstlisting}
#### 6. }time\code{  - 

}time\code{24HH:MMHH:MM:SS

}min\code{}max\code{}step\code{}step\code{}step="300"\code{5

\end{lstlisting}

html

\begin{lstlisting}
#### 7. }datetime-local\code{  - 

}datetime-local\code{}date\code{}time\code{

ISO 8601YYYY-MM-DDTHH:MM""

\end{lstlisting}

html

\begin{lstlisting}
#### 8. }color\code{  - 

}color\code{#4a90e2

CSSJavaScript

\end{lstlisting}

html

\begin{lstlisting}
#### 9. }search\code{  - 

}search\code{×""""

}search\code{}results\code{}placeholder\code{

\end{lstlisting}

html

\begin{lstlisting}
### 

####  - }fieldset\code{  }legend\code{

\end{lstlisting}

html

\begin{lstlisting}
#### 

\end{lstlisting}

html

\begin{lstlisting}
#### 

\end{lstlisting}

html

\begin{lstlisting}
}type="number"\code{}min\code{}max\code{}type="datetime-local"\code{}type="range"\code{

\important{} HTML5HTML5}required\code{}pattern\code{}minlength\code{}maxlength\code{

HTML5CSS}:valid\code{}:invalid\code{}:required\code{JavaScriptsetCustomValidity()

}step="0.01"\code{}min="0"\code{}pattern\code{

}<label>\code{}<fieldset>\code{}<legend>\code{}aria-describedby\code{

## 4.2.3 HTML5

HTML5WebHTML5

HTML5}<audio>\code{}<video>\code{Flash

### HTML5

||  ||  |
|--------|------|----------|-------------|
|}<video>\code{|  |}src\code{, }controls\code{, }autoplay\code{, }loop\code{, }muted\code{, }poster\code{|  |
|}<audio>\code{|  |}src\code{, }controls\code{, }autoplay\code{, }loop\code{, }muted\code{|  |
|}<source>\code{|  |}src\code{, }type\code{, }media\code{|  |
|}<track>\code{|  |}src\code{, }kind\code{, }srclang\code{, }label\code{|  |
|}<canvas>\code{|  |}width\code{, }height\code{|  |
|}<svg>\code{|  |}width\code{, }height\code{, }viewBox\code{|  |

### HTML5

#### 1. }<video>\code{  - 

}<video>\code{HTML5FlashMP4WebMOgg}<source>\code{

}<video>\code{}controls\code{}autoplay\code{}loop\code{}muted\code{}poster\code{

\end{lstlisting}

html

\begin{lstlisting}
#### 2. }<audio>\code{  - 

}<audio>\code{}<video>\code{MP3WAVOgg}<source>\code{

}controls\code{}autoplay\code{}loop\code{}preload\code{

\end{lstlisting}

html

\begin{lstlisting}
#### 3. }<source>\code{  - 

}<source>\code{HTML5}<video>\code{}<audio>\code{

}<source>\code{""

}<source>\code{}media\code{

\end{lstlisting}

html

\begin{lstlisting}
#### 4. }<track>\code{  - 

}<track>\code{}<video>\code{

}<track>\code{WebVTTWeb Video Text TracksWeb}kind\code{}subtitles\code{}captions\code{}descriptions\code{}chapters\code{}metadata\code{

}<track>\code{

\end{lstlisting}

html

\begin{lstlisting}
#### 5. }<canvas>\code{  - 

}<canvas>\code{JavaScriptHTML52D

Canvas APIAPICanvas

Canvas

\end{lstlisting}

html

Canvas

\begin{lstlisting}
#### 6. }<svg>\code{  - 

}<svg>\code{Scalable Vector GraphicsXMLCanvasSVGSVGDOMCSSJavaScript

SVGCSSJavaScriptSEOSVG

SVGSVG

\end{lstlisting}

html

\begin{verbatim}
<!--  -->
<rect id="waterFill" x="32" y="120" width="56" height="58" 
      fill="url(#waterGradient)" opacity="0.8">
    <animate attributeName="height" 
             values="20;80;60;100;70" 
             dur="10s" 
             repeatCount="indefinite"/>
    <animate attributeName="y" 
             values="160;100;120;80;110" 
             dur="10s" 
             repeatCount="indefinite"/>
</rect>

<!--  -->
<defs>
    <linearGradient id="waterGradient" x1="0%" y1="0%" x2="0%" y2="100%">
        <stop offset="0%" style="stop-color:#87CEEB;stop-opacity:1" />
        <stop offset="100%" style="stop-color:#4169E1;stop-opacity:1" />
    </linearGradient>
</defs>

<!--  -->
<g stroke="#666" stroke-width="1">
    <line x1="90" y1="60" x2="100" y2="60"/>
    <text x="105" y="65" font-size="12" fill="#666">90m</text>
    
    <line x1="90" y1="100" x2="100" y2="100"/>
    <text x="105" y="105" font-size="12" fill="#666">80m</text>
    
    <line x1="90" y1="140" x2="100" y2="140"/>
    <text x="105" y="145" font-size="12" fill="#666">70m</text>
    
    <line x1="90" y1="180" x2="100" y2="180"/>
    <text x="105" y="185" font-size="12" fill="#666">60m</text>
</g>

<!--  -->
<text x="60" y="15" text-anchor="middle" font-size="14" font-weight="bold" fill="#333">
    
</text>
\end{verbatim}

\begin{lstlisting}
### 

HTML5JavaScript

#### 

JavaScriptAPI

\end{lstlisting}

html

\begin{verbatim}
<button onclick="playVideo()"></button>
<button onclick="pauseVideo()"></button>
<button onclick="changeVolume(0.5)">50%</button>
<button onclick="seekTo(30)">30</button>
<button onclick="toggleFullscreen()"></button>
\end{verbatim}

\begin{lstlisting}
#### 

\end{lstlisting}

html

\begin{verbatim}
<button onclick="playAlert('warning')"></button>
<button onclick="playAlert('danger')"></button>
<button onclick="playAlert('emergency')"></button>
<button onclick="stopAllAlerts()"></button>
\end{verbatim}

}`\code{

HTML5CanvasSVGCanvasJavaScript API2DSVGCanvasSVG

\important{} CanvasCanvas API

CanvasSVGCanvasSVGDOM

WebGLWeb Audio APIWebRTC

## 4.2.4 HTML

HTMLHTML

HTML}

\code{}

\code{}

\code{}

\code{}

\code{}\code{

HTMLHTML}

\code{}

\code{}

\code{}

\code{}

\code{}

\code{}

\code{}

\code{}\code{

\important{} HTML

SEOSEOHTMLSEOh1h2h3}

\code{}\code{

HTMLWebHTML}

`HTML

HTMLHTMLDOMCSSJavaScript

HTML5HTML5CSSHTML

\hypertarget{css}{%
\section{4.3 CSS}\label{css}}

CSSCascading Style SheetsWebHTMLJavaScriptWebCSSCSS

CSS3CSSWebCSS3

!!! info ``CSS''

\begin{verbatim}
CSS3CSSCSS
\end{verbatim}

\hypertarget{css-1}{%
\subsection{CSS}\label{css-1}}

\hypertarget{css-2}{%
\subsubsection{CSS}\label{css-2}}

CSSCascading Style SheetsHTMLXMLCSS\important{}

CSS - \important{Cascading} - \important{Inheritance} -
\important{Separation}HTMLCSS

\hypertarget{css-3}{%
\subsubsection{CSS}\label{css-3}}

CSS\important{}HTML\important{}

\begin{enumerate}
\def\labelenumi{\arabic{enumi}.}
\item
  \important{}CSS
\item
  \important{}HTML
\item
  \important{}
\end{enumerate}

\hypertarget{css-4}{%
\subsubsection{CSS}\label{css-4}}

CSS\important{}\important{}

\begin{lstlisting}[style=css]
 {
    : ;
    : ;
}
\end{lstlisting}

\begin{lstlisting}[style=css]
h1 {
    color: blue;
    font-size: 24px;
    margin: 10px 0;
}
\end{lstlisting}

\important{}

\begin{itemize}
\item
  \important{Selector}HTML
\item
  \important{} \code{{}}
\item
  \important{Property}
\item
  \important{Value}
\item
  \important{;}
\end{itemize}

\hypertarget{css-5}{%
\subsubsection{CSS}\label{css-5}}

CSSHTML

\begin{enumerate}
\def\labelenumi{\arabic{enumi}.}
\item
  \important{}HTMLstyle
  \textbackslash begin\{lstlisting}{[}style=xml{]}

  ``\textbackslash code\{
\item
  \important{}HTML}

  \code{}

  \code{
  }\texttt{\textbackslash{}code\{html\ \ \ \ \textless{}head\textgreater{}\ \ \ \ \ \ \ \ \textless{}style\textgreater{}\ \ \ \ \ \ \ \ \ \ \ \ p\ \{\ \ \ \ \ \ \ \ \ \ \ \ \ \ \ \ color:\ blue;\ \ \ \ \ \ \ \ \ \ \ \ \ \ \ \ font-size:\ 14px;\ \ \ \ \ \ \ \ \ \ \ \ }\ \ \ \ \ \ \ \ \textless{}/style\textgreater{}\ \ \ \ \textless{}/head\textgreater{}\ \ \ \ }}\textbackslash code\{
\item
  \important{}CSS.css}\code{
  }\texttt{\textbackslash{}code\{html\ \ \ \ \textless{}head\textgreater{}\ \ \ \ \ \ \ \ \textless{}link\ rel="stylesheet"\ href="styles.css"\textgreater{}\ \ \ \ \textless{}/head\textgreater{}\ \ \ \ }}\textbackslash code\{
\end{enumerate}

\hypertarget{css-6}{%
\subsubsection{CSS}\label{css-6}}

\important{Cascading}CSS

\important{}

\begin{enumerate}
\def\labelenumi{\arabic{enumi}.}
\tightlist
\item
  style
\item
  ID
\item
\item
\end{enumerate}

\important{Inheritance}

\textbackslash end\{lstlisting}css body \{ font-family: ``Microsoft
YaHei''; color: \#333; } /\highlight{ body }/

\begin{lstlisting}
\important{}
- }font-family\code{, }font-size\code{, }color\code{, }line-height\code{
- }text-align\code{, }text-indent\code{, }text-transform\code{
- }list-style\code{

\important{}
- }width\code{, }height\code{, }margin\code{, }padding\code{, }border\code{
- }position\code{, }top\code{, }left\code{
- }display\code{, }float\code{

### CSS

CSS }/\highlight{ }/\code{ 

\end{lstlisting}

css /\highlight{  }/ p \{ color: blue; /\highlight{  }/ }

/\highlight{

}/

\begin{lstlisting}
CSSCSS3

## 4.3.1 CSS3

CSSCSSHTMLCSS3

CSSID

### CSS

#### 1.  - HTML

CSSHTML

HTML

\end{lstlisting}

css /\highlight{  }/ table \{ border-collapse: collapse;
/\highlight{  }/ width: 100\%; margin: 20px 0; /\highlight{ ...  ... }/
border-radius: 4px; border-left: 4px solid \#2196f3; /\highlight{  }/ }

\begin{lstlisting}
#### 2.  - 

CSSHTMLclass

HTMLclass

\end{lstlisting}

css /\highlight{  }/ .water-level \{ font-size: 18px; font-weight: bold;
color: \#2196f3; /\highlight{ ...  ... }/ box-shadow: 0 4px 12px
rgba(0,0,0,0.15); transform: translateY(-2px); }

\begin{lstlisting}
#### 3. ID - 

IDIDCSSID

ID"ID"IDIDJavaScript

\end{lstlisting}

css /\highlight{  }/ \#main-nav \{ background: \#1565c0; color: white;
padding: 0 20px; /\highlight{ ...  ... }/ gap: 20px; margin-bottom:
30px; }

\begin{lstlisting}
#### 4.  - HTML

CSS3HTML

\end{lstlisting}

css /\highlight{  }/ input{[}required{]} \{ border-left: 4px solid
\#2196f3; background: rgba(33, 150, 243, 0.05); }
/\highlight{ ...  ... }/ {[}data-device=``flow-meter''{]}::before \{
content: '`; } {[}data-device=``rain-gauge''{]}::before \{ content:'`;
} {[}data-device=``temperature''{]}::before \{ content:''; }

\begin{lstlisting}
### CSS3

CSS3

#### 1.  - 

HTML

\end{lstlisting}

css /\highlight{  -  }/ .monitor-panel \textgreater{} .data-item \{
border-bottom: 1px solid \#eee; padding: 10px 0; display: flex;
/\highlight{ ...  ... }/ border-color: \#f44336; /\highlight{  }/
background: rgba(244, 67, 54, 0.05); }

\begin{lstlisting}
#### 2.  - 

JavaScript

\end{lstlisting}

css /\highlight{  -  }/ .monitor-card:hover \{ box-shadow: 0 8px 16px
rgba(0,0,0,0.15); transform: translateY(-3px) scale(1.02); transition:
all 0.3s cubic-bezier(0.4, 0.0, 0.2, 1); /\highlight{ ...  ... }/
border-radius: 4px; } .nav-link:active \{ color: \#01579b; }

\begin{lstlisting}
#### 3.  - 

HTMLHTML

}::before\code{}::after\code{}content\code{

\end{lstlisting}

css /\highlight{  }/ .water-level::before \{ content: ``\,``;
margin-right: 8px; font-size: 1.2em; /\highlight{ ...  ... }/ color:
\#1565c0; font-size: 1.1em; }

\begin{lstlisting}
### CSS3

CSS3

#### 1. 

CSS3

\end{lstlisting}

css /\highlight{  }/ .rounded-card \{ border-radius: 12px;
/\highlight{  }/ background: white; border: 1px solid \#e0e0e0;
/\highlight{ ...  ... }/ linear-gradient(white, white) padding-box,
linear-gradient(45deg, \#2196f3, \#4caf50, \#ff9800) border-box; }

\begin{lstlisting}
#### 2. 

CSS3

\end{lstlisting}

css /\highlight{  }/ .title-shadow \{ text-shadow: 2px 2px 4px
rgba(0,0,0,0.3), /\highlight{  }/ 1px 1px 2px rgba(0,0,0,0.5);
/\highlight{  }/ /\highlight{ ...  ... }/ padding: 2px 6px;
border-radius: 4px; }

\begin{lstlisting}
## 4.3.2 CSS3

WebCSS3FlexboxGrid

CSSfloatpositiondisplayCSS3FlexboxGridWebFlexboxGrid

### Flexbox

FlexboxCSS3

Flexboxmain axiscross axis

#### 1. Flex

}display: flex\code{}display: inline-flex\code{

\end{lstlisting}

css /\highlight{  }/ .flex-container \{ display: flex; justify-content:
space-between; /\highlight{  }/ align-items: center; /\highlight{  }/
/\highlight{ ...  ... }/ flex-direction: column; /\highlight{  }/ } }

\begin{lstlisting}
#### 2. Flex

CSS

\end{lstlisting}

css /\highlight{  }/ .monitor-dashboard \{ display: flex; gap: 20px;
flex-wrap: wrap; /\highlight{  }/ /\highlight{ ...  ... }/ text-align:
left; } }

\begin{lstlisting}
#### 3. Flex

Flexbox

Flex

\end{lstlisting}

css /\highlight{  }/ .platform-header \{ display: flex; align-items:
center; padding: 0 20px; /\highlight{ ...  ... }/ border-radius: 8px;
padding: 20px; }

\begin{lstlisting}
### Grid

CSS GridCSS3FlexboxGridGridExcel

#### 1. Grid

Grid}display: grid\code{rowcolumngrid linegrid area

\end{lstlisting}

css /\highlight{  }/ .main-dashboard \{ display: grid;
/\highlight{ () () () }/ grid-template-columns: 260px 1fr 320px;
/\highlight{ ...  ... }/ justify-content: center; gap: 10px; }

\begin{lstlisting}
#### 2. Grid

Grid

\end{lstlisting}

css /\highlight{  }/ .water-monitoring-platform \{ display: grid;
grid-template-areas: ``logo navigation user-info'' ``sidebar
main-dashboard alerts'' ``sidebar data-tables alerts'' ``footer footer
footer''; grid-template-columns: 280px 1fr 350px; grid-template-rows:
70px 400px 1fr 60px; gap: 20px; padding: 20px; min-height: 100vh;
background: \#f8fafc; }

/\highlight{  }/ .platform-logo \{ grid-area: logo; display: flex;
align-items: center; justify-content: center; background:
linear-gradient(45deg, \#2196F3, \#21CBF3); color: white; font-weight:
bold; border-radius: 8px; }

.main-navigation \{ grid-area: navigation; display: grid;
grid-template-columns: repeat(auto-fit, minmax(120px, 1fr)); gap: 10px;
align-items: center; background: white; border-radius: 8px; padding: 0
20px; box-shadow: 0 2px 10px rgba(0,0,0,0.05); }

.user-info-panel \{ grid-area: user-info; background: white;
border-radius: 8px; padding: 15px; display: flex; align-items: center;
justify-content: space-between; box-shadow: 0 2px 10px rgba(0,0,0,0.05);
}

/\highlight{  }/ .critical-alert-banner \{ /\highlight{ 12 }/ grid-row:
1 / 2; grid-column: 1 / -1; /\highlight{ -1 }/ background:
linear-gradient(90deg, \#ff6b6b, \#ff8e8e); color: white; padding: 10px
20px; border-radius: 6px; text-align: center; font-weight: 600;
animation: pulse 2s infinite; z-index: 10; }

/\highlight{  }/ .data-export-section \{ /\highlight{  }/ grid-column:
sidebar-end / alerts-start; grid-row: data-tables-start / footer-start;
background: white; border-radius: 8px; padding: 25px; display: grid;
grid-template-rows: auto 1fr auto; gap: 15px; }

/\highlight{  }/ .full-screen-report \{ /\highlight{  }/ grid-row: 2 /
-1; grid-column: 1 / -1; background: rgba(255,255,255,0.98);
backdrop-filter: blur(10px); border-radius: 12px; padding: 30px;
box-shadow: 0 15px 35px rgba(0,0,0,0.1); z-index: 100; display: none;
/\highlight{  }/ }

.full-screen-report.active \{ display: grid; grid-template-rows: auto
1fr auto; gap: 20px; }

/\highlight{  }/ .dynamic-sensor-card \{ /\highlight{ JavaScript }/
background: white; border-radius: 8px; padding: 20px; border-left: 4px
solid \#3498db; box-shadow: 0 3px 12px rgba(0,0,0,0.08); transition: all
0.3s ease;

\begin{verbatim}
/\highlight{  }/
display: grid;
grid-template-areas: 
    "icon title status"
    "icon value trend"
    "controls controls controls";
grid-template-columns: 60px 1fr auto;
grid-template-rows: auto auto auto;
gap: 10px 15px;
align-items: center;
\end{verbatim}

}

.sensor-icon \{ grid-area: icon; width: 50px; height: 50px;
border-radius: 50\%; background: linear-gradient(135deg, \#667eea 0\%,
\#764ba2 100\%); display: flex; align-items: center; justify-content:
center; color: white; font-size: 1.5em; }

.sensor-title \{ grid-area: title; font-weight: 600; color: \#2c3e50;
margin: 0; }

.sensor-status \{ grid-area: status; padding: 4px 12px; border-radius:
20px; font-size: 0.8em; font-weight: 500; text-align: center; }

.sensor-value \{ grid-area: value; font-size: 1.4em; font-weight: bold;
color: \#27ae60; font-family: `Consolas', monospace; }

.sensor-trend \{ grid-area: trend; font-size: 0.9em; color: \#7f8c8d; }

.sensor-controls \{ grid-area: controls; display: flex; gap: 8px;
margin-top: 10px; }

\begin{lstlisting}
#### 3. Grid

Grid

\end{lstlisting}

css /\highlight{  }/ .monitoring-cards-responsive \{ display: grid;
/\highlight{ auto-fitminmax }/ grid-template-columns: repeat(auto-fit,
minmax(320px, 1fr)); /\highlight{ ...  ... }/ font-size: 0.9em; } }

\begin{lstlisting}
### 

WebCSS

#### 1. 

CSS

\end{lstlisting}

css /\highlight{  }/

/\highlight{  (≥1400px) }/ @media screen and (min-width: 1400px) \{
.water-monitoring-dashboard \{ max-width: 1360px; margin: 0 auto;
grid-template-columns: 320px 1fr 400px; /\highlight{  }/ }

\begin{verbatim}
.monitoring-cards-grid {
    grid-template-columns: repeat(4, 1fr);    /\highlight{ 4 }/
    gap: 30px 25px;
}

.data-charts-section {
    grid-template-columns: repeat(3, 1fr);    /\highlight{ 3 }/
}

.detailed-table {
    font-size: 1rem;                         /\highlight{  }/
}

/\highlight{  }/
.large-screen-details {
    display: block;
}
\end{verbatim}

}

/\highlight{  (1200px - 1399px) }/ @media screen and (min-width: 1200px)
and (max-width: 1399px) \{ .water-monitoring-dashboard \{
grid-template-columns: 280px 1fr 350px; padding: 20px; }

\begin{verbatim}
.monitoring-cards-grid {
    grid-template-columns: repeat(3, 1fr);    /\highlight{ 3 }/
    gap: 25px 20px;
}

.data-charts-section {
    grid-template-columns: repeat(2, 1fr);    /\highlight{ 2 }/
}

/\highlight{  }/
.card-title {
    font-size: 1.1em;
}

.data-value {
    font-size: 1.8em;
}
\end{verbatim}

}

/\highlight{  (992px - 1199px) }/ @media screen and (min-width: 992px)
and (max-width: 1199px) \{ .water-monitoring-dashboard \{
grid-template-areas: ``header header header'' ``sidebar content info''
``footer footer footer''; grid-template-columns: 250px 1fr 300px;
grid-template-rows: 70px 1fr 50px; gap: 15px; }

\begin{verbatim}
.monitoring-cards-grid {
    grid-template-columns: repeat(2, 1fr);    /\highlight{ 2 }/
    gap: 20px 15px;
}

.data-charts-section {
    grid-template-columns: 1fr 1fr;
}

/\highlight{  }/
.main-navigation {
    font-size: 0.9em;
    gap: 8px;
}

.nav-item {
    padding: 8px 12px;
}
\end{verbatim}

}

/\highlight{  (768px - 991px) }/ @media screen and (min-width: 768px)
and (max-width: 991px) \{ .water-monitoring-dashboard \{
grid-template-areas: ``header header'' ``sidebar content'' ``info info''
``footer footer''; grid-template-columns: 200px 1fr; grid-template-rows:
auto 1fr auto auto; gap: 15px; padding: 15px; }

\begin{verbatim}
.monitoring-cards-grid {
    grid-template-columns: repeat(2, 1fr);
    gap: 18px 12px;
}

.data-charts-section {
    grid-template-columns: 1fr;              /\highlight{  }/
    gap: 20px;
}

/\highlight{  }/
.data-table-responsive {
    display: grid;
    grid-template-columns: 1fr;
    gap: 10px;
}

.table-row-tablet {
    display: grid;
    grid-template-columns: repeat(3, 1fr);
    gap: 10px;
    padding: 12px;
    background: white;
    border-radius: 6px;
    box-shadow: 0 2px 4px rgba(0,0,0,0.1);
}

/\highlight{  }/
.secondary-info {
    display: none;
}

.primary-data {
    font-size: 1.2em;
    font-weight: bold;
}
\end{verbatim}

}

/\highlight{  (576px - 767px) }/ @media screen and (min-width: 576px)
and (max-width: 767px) \{ .water-monitoring-dashboard \{
grid-template-areas: ``header'' ``content'' ``sidebar'' ``info''
``footer''; grid-template-columns: 1fr; gap: 12px; padding: 12px; }

\begin{verbatim}
.monitoring-cards-grid {
    grid-template-columns: 1fr 1fr;          /\highlight{ 2 }/
    gap: 15px 10px;
}

.monitoring-card-compact {
    padding: 15px;
    display: grid;
    grid-template-areas: 
        "title value"
        "status trend";
    grid-template-columns: 1fr auto;
    gap: 8px;
    align-items: center;
}

.card-title {
    grid-area: title;
    font-size: 0.9em;
    font-weight: 600;
    color: #2c3e50;
}

.card-value {
    grid-area: value;
    font-size: 1.3em;
    font-weight: bold;
    color: #27ae60;
    text-align: right;
}

.card-status {
    grid-area: status;
    font-size: 0.8em;
    color: #7f8c8d;
}

.card-trend {
    grid-area: trend;
    font-size: 0.8em;
    text-align: right;
}

/\highlight{  }/
.main-navigation {
    display: grid;
    grid-template-columns: repeat(auto-fit, minmax(80px, 1fr));
    gap: 5px;
    font-size: 0.8em;
}
\end{verbatim}

}

/\highlight{  (≤575px) }/ @media screen and (max-width: 575px) \{
.water-monitoring-dashboard \{ grid-template-areas: ``header''
``content'' ``footer''; grid-template-columns: 1fr; gap: 10px; padding:
10px; margin: 0; }

\begin{verbatim}
.monitoring-cards-grid {
    grid-template-columns: 1fr;              /\highlight{  }/
    gap: 12px;
}

.monitoring-card-mobile {
    padding: 12px;
    display: grid;
    grid-template-areas: 
        "title"
        "value"
        "details"
        "actions";
    gap: 10px;
    border-left: 4px solid #3498db;
}

.card-title-mobile {
    grid-area: title;
    font-size: 1em;
    font-weight: 600;
    color: #2c3e50;
    margin: 0;
}

.card-value-mobile {
    grid-area: value;
    font-size: 2em;
    font-weight: bold;
    color: #27ae60;
    text-align: center;
    padding: 10px 0;
    background: #f8f9fa;
    border-radius: 4px;
}

.card-details-mobile {
    grid-area: details;
    display: grid;
    grid-template-columns: 1fr 1fr;
    gap: 8px;
    font-size: 0.8em;
    color: #666;
}

.card-actions-mobile {
    grid-area: actions;
    display: flex;
    gap: 8px;
    justify-content: center;
}

.btn-mobile {
    flex: 1;
    padding: 8px 12px;
    border: none;
    border-radius: 4px;
    font-size: 0.9em;
    cursor: pointer;
    transition: all 0.2s ease;
}

/\highlight{  }/
.mobile-sidebar {
    position: fixed;
    top: 0;
    left: -300px;                            /\highlight{  }/
    width: 300px;
    height: 100vh;
    background: white;
    box-shadow: 2px 0 10px rgba(0,0,0,0.1);
    transition: left 0.3s ease;
    z-index: 1000;
    overflow-y: auto;
}

.mobile-sidebar.active {
    left: 0;                                 /\highlight{  }/
}

.mobile-overlay {
    position: fixed;
    top: 0;
    left: 0;
    right: 0;
    bottom: 0;
    background: rgba(0,0,0,0.5);
    z-index: 999;
    display: none;
}

.mobile-overlay.active {
    display: block;
}

/\highlight{  }/
.table-mobile {
    display: block;
}

.table-row-mobile {
    display: block;
    background: white;
    margin-bottom: 10px;
    border-radius: 6px;
    padding: 12px;
    box-shadow: 0 2px 6px rgba(0,0,0,0.1);
}

.table-cell-mobile {
    display: flex;
    justify-content: space-between;
    align-items: center;
    padding: 6px 0;
    border-bottom: 1px dotted #ddd;
}

.table-cell-mobile:last-child {
    border-bottom: none;
}

.cell-label {
    font-weight: 600;
    color: #666;
    font-size: 0.9em;
}

.cell-value {
    font-weight: bold;
    color: #2c3e50;
}
\end{verbatim}

}

/\highlight{  }/ @media screen and (-webkit-min-device-pixel-ratio: 2),
screen and (min-resolution: 192dpi) \{ .icon, .chart-canvas,
.data-visualization \{ image-rendering: -webkit-optimize-contrast;
image-rendering: optimizeQuality; }

\begin{verbatim}
/\highlight{  }/
body {
    -webkit-font-smoothing: antialiased;
    -moz-osx-font-smoothing: grayscale;
}
\end{verbatim}

}

/\highlight{  }/ @media print \{ .water-monitoring-dashboard \{ display:
block !important; grid-template-areas: none !important;
grid-template-columns: none !important; }

\begin{verbatim}
.no-print {
    display: none !important;
}

.monitoring-card {
    break-inside: avoid;
    margin-bottom: 20px;
}

.chart-container {
    break-inside: avoid;
}

body {
    font-size: 12pt;
    line-height: 1.4;
}
\end{verbatim}

}

\begin{lstlisting}
#### 2. 

\end{lstlisting}

css /\highlight{  }/ .touch-friendly-button \{ min-height: 44px;
/\highlight{ Apple }/ min-width: 44px; padding: 12px 20px;
/\highlight{ ...  ... }/ zoom: 1; } }

\begin{lstlisting}
## 4.3.4 CSS3

CSS3Web

### CSS3

(Transition)CSS3

#### 1. 

CSS3CSS

}transition-property\code{}transition-duration\code{}transition-timing-function\code{}transition-delay\code{

\end{lstlisting}

css /\highlight{  }/ .water-data-card \{ background: white;
border-radius: 12px; padding: 25px; margin: 15px; border-left: 5px solid
\#3498db; box-shadow: 0 4px 6px rgba(0,0,0,0.1); cursor: pointer;
position: relative; overflow: hidden;

\begin{verbatim}
/\highlight{  }/
transition: 
    transform 0.3s cubic-bezier(0.25, 0.8, 0.25, 1),    /\highlight{  }/
    box-shadow 0.3s cubic-bezier(0.25, 0.8, 0.25, 1),  /\highlight{  }/
    border-left-width 0.2s ease-out,                   /\highlight{  }/
    background-color 0.2s ease-in-out;                 /\highlight{  }/
\end{verbatim}

}

.water-data-card:hover \{ transform: translateY(-8px) scale(1.02);
/\highlight{  }/ box-shadow: 0 10px 25px rgba(0,0,0,0.15),
/\highlight{  }/ 0 5px 10px rgba(52, 152, 219, 0.2); /\highlight{  }/
border-left-width: 8px; /\highlight{  }/ background-color: \#fafbfc;
/\highlight{  }/ }

.water-data-card:active \{ transform: translateY(-2px) scale(0.98);
/\highlight{  }/ transition-duration: 0.1s; /\highlight{  }/ }

/\highlight{  }/ .critical-value \{ font-size: 2.2em; font-weight: bold;
color: \#2c3e50; display: inline-block; position: relative;

\begin{verbatim}
/\highlight{  }/
transition: 
    color 0.5s ease,
    transform 0.3s cubic-bezier(0.68, -0.55, 0.265, 1.55); /\highlight{  }/
\end{verbatim}

}

/\highlight{  }/ .critical-value.warning \{ color: \#f39c12; animation:
valueWarning 2s ease-in-out infinite; }

.critical-value.danger \{ color: \#e74c3c; animation: valueDanger 1.5s
ease-in-out infinite; }

.critical-value.updated \{ transform: scale(1.15); color: \#27ae60; }

@keyframes valueWarning \{ 0\%, 100\% \{ opacity: 1; transform:
scale(1); } 50\% \{ opacity: 0.7; transform: scale(1.05); } }

@keyframes valueDanger \{ 0\%, 100\% \{ opacity: 1; transform: scale(1);
box-shadow: 0 0 0 0 rgba(231, 76, 60, 0.4); } 25\% \{ opacity: 0.8;
transform: scale(1.08); box-shadow: 0 0 0 8px rgba(231, 76, 60, 0.2); }
50\% \{ opacity: 1; transform: scale(1.12); box-shadow: 0 0 0 12px
rgba(231, 76, 60, 0.1); } 75\% \{ opacity: 0.9; transform: scale(1.05);
box-shadow: 0 0 0 6px rgba(231, 76, 60, 0.15); } }

/\highlight{  }/ .progress-indicator \{ width: 100\%; height: 8px;
background: \#ecf0f1; border-radius: 4px; overflow: hidden; position:
relative; }

.progress-bar \{ height: 100\%; background: linear-gradient(90deg,
\#3498db, \#2ecc71); border-radius: 4px;

\begin{verbatim}
/\highlight{  }/
transition: width 0.8s cubic-bezier(0.25, 0.46, 0.45, 0.94);

/\highlight{  }/
position: relative;
overflow: hidden;
\end{verbatim}

}

.progress-bar::before \{ content: '\,'; position: absolute; top: 0;
left: -100\%; width: 100\%; height: 100\%; background: linear-gradient(
90deg, transparent, rgba(255,255,255,0.4), transparent ); animation:
progressShine 2s infinite; }

@keyframes progressShine \{ 0\% \{ left: -100\%; } 50\% \{ left: 100\%;
} 100\% \{ left: 100\%; } }

/\highlight{  }/ .status-indicator \{ width: 12px; height: 12px;
border-radius: 50\%; display: inline-block; position: relative;

\begin{verbatim}
/\highlight{  }/
transition: 
    background-color 0.4s ease,
    box-shadow 0.4s ease,
    transform 0.2s ease;
\end{verbatim}

}

.status-indicator.online \{ background-color: \#27ae60; box-shadow: 0 0
0 3px rgba(39, 174, 96, 0.2), inset 0 1px 1px rgba(255,255,255,0.3); }

.status-indicator.offline \{ background-color: \#e74c3c; box-shadow: 0 0
0 3px rgba(231, 76, 60, 0.2), inset 0 1px 1px rgba(0,0,0,0.1); }

.status-indicator.maintenance \{ background-color: \#f39c12; box-shadow:
0 0 0 3px rgba(243, 156, 18, 0.2), inset 0 1px 1px
rgba(255,255,255,0.3); animation: maintenanceBlink 2s ease-in-out
infinite; }

@keyframes maintenanceBlink \{ 0\%, 100\% \{ opacity: 1; } 50\% \{
opacity: 0.5; } }

/\highlight{  }/ .form-field \{ position: relative; margin-bottom: 20px;
}

.form-input \{ width: 100\%; padding: 12px 16px; border: 2px solid
\#e1e8ed; border-radius: 6px; font-size: 16px; background: white;

\begin{verbatim}
/\highlight{  }/
transition: 
    border-color 0.3s ease,
    box-shadow 0.3s ease,
    background-color 0.2s ease;
\end{verbatim}

}

.form-input:focus \{ outline: none; border-color: \#3498db; box-shadow:
0 0 0 3px rgba(52, 152, 219, 0.1); background-color: \#fafbfc; }

.form-input.valid \{ border-color: \#27ae60; background-color: rgba(39,
174, 96, 0.05); }

.form-input.invalid \{ border-color: \#e74c3c; background-color:
rgba(231, 76, 60, 0.05); animation: shakeError 0.5s ease-in-out; }

@keyframes shakeError \{ 0\%, 100\% \{ transform: translateX(0); }
10\%, 30\%, 50\%, 70\%, 90\% \{ transform: translateX(-5px); } 20\%,
40\%, 60\%, 80\% \{ transform: translateX(5px); } }

/\highlight{  }/ .error-message \{ color: \#e74c3c; font-size: 14px;
margin-top: 8px; opacity: 0; transform: translateY(-10px);

\begin{verbatim}
transition: 
    opacity 0.3s ease,
    transform 0.3s ease;
\end{verbatim}

}

.error-message.show \{ opacity: 1; transform: translateY(0); }

/\highlight{  }/ .action-button \{ background: \#3498db; color: white;
border: none; padding: 12px 24px; border-radius: 6px; font-size: 16px;
font-weight: 600; cursor: pointer; position: relative; overflow: hidden;

\begin{verbatim}
/\highlight{  }/
transition: 
    background-color 0.2s ease,
    transform 0.1s ease,
    box-shadow 0.2s ease;
\end{verbatim}

}

.action-button:hover \{ background: \#2980b9; transform:
translateY(-2px); box-shadow: 0 4px 12px rgba(52, 152, 219, 0.3); }

.action-button:active \{ transform: translateY(0); box-shadow: 0 2px 4px
rgba(52, 152, 219, 0.2); }

.action-button:disabled \{ background: \#bdc3c7; cursor: not-allowed;
transform: none; box-shadow: none; }

/\highlight{  }/ .action-button.loading \{ color: transparent; cursor:
wait; }

.action-button.loading::after \{ content: '\,'; position: absolute; top:
50\%; left: 50\%; width: 16px; height: 16px; margin: -8px 0 0 -8px;
border: 2px solid rgba(255,255,255,0.3); border-top: 2px solid white;
border-radius: 50\%; animation: buttonSpin 1s linear infinite; }

@keyframes buttonSpin \{ 0\% \{ transform: rotate(0deg); } 100\% \{
transform: rotate(360deg); } }

/\highlight{  }/ .monitor-card \{ background: white; border-radius: 8px;
\# \ldots{} \ldots{} transform: scale(0.98); transition: transform 0.1s
ease; }

\begin{lstlisting}
#### 2. 

\end{lstlisting}

css /\highlight{  }/ .form-input \{ border: 2px solid \#e0e0e0; padding:
12px; border-radius: 4px; /\highlight{ ...  ... }/ transform:
translateY(-15px) scale(0.8); color: \#2196f3; }

\begin{lstlisting}
### CSS3

#### 1. 

\end{lstlisting}

css /\highlight{  }/ @keyframes spin \{ from \{ transform: rotate(0deg);
} to \{ transform: rotate(360deg); } } /\highlight{ ...  ... }/
background: rgba(33, 150, 243, 0.3); animation: ripple 2s infinite; }

\begin{lstlisting}
#### 2. 

\end{lstlisting}

css /\highlight{  }/ @keyframes dataUpdate \{ 0\%, 100\% \{ background:
transparent; } 50\% \{ background: rgba(76, 175, 80, 0.3); } }
/\highlight{ ...  ... }/ .warning-indicator \{ animation: pulse 2s
infinite; }

\begin{lstlisting}
#### 3. 

\end{lstlisting}

css /\highlight{  }/ @keyframes progressFill \{ from \{ width: 0\%; }
to \{ width: 100\%; } } /\highlight{ ...  ... }/ background:
linear-gradient(to top, \#1976d2, \#42a5f5); animation: waterRise 2s
ease-out; }

\begin{lstlisting}
## 4.3.4 UI

UI

### 

#### 1. 

CSS

\end{lstlisting}

css /\highlight{  }/ :root \{ /\highlight{  -  }/ --primary-color:
\#1565c0; --primary-light: \#42a5f5; /\highlight{ ...  ... }/
.status-normal \{ color: var(--success-color); } .status-warning \{
color: var(--warning-color); } .status-error \{ color:
var(--error-color); }

\begin{lstlisting}
#### 2. 

\end{lstlisting}

css /\highlight{  }/ .water-gradient-1 \{ background:
linear-gradient(135deg, \#1565c0, \#42a5f5); }

\begin{verbatim}
/\highlight{ ...  ... }/
    var(--secondary-color),
    var(--success-color));
\end{verbatim}

}

\begin{lstlisting}
### 

#### 1. 

\end{lstlisting}

css /\highlight{  }/ .btn \{ display: inline-flex; align-items: center;
justify-content: center; /\highlight{ ...  ... }/ background:
var(--primary-color); color: white; }

\begin{lstlisting}
#### 2. 

Web

\end{lstlisting}

css /\highlight{  }/ .card \{ background: white; border-radius: 8px;
box-shadow: 0 2px 8px rgba(0, 0, 0, 0.1); /\highlight{ ...  ... }/
.monitor-card.error::before \{ background: var(--error-color); }

\begin{lstlisting}
#### 3. 

\end{lstlisting}

css /\highlight{  }/ .data-table \{ width: 100\%; border-collapse:
collapse; background: white; /\highlight{ ...  ... }/
.status-indicator.normal \{ background: var(--success-color); }
.status-indicator.warning \{ background: var(--warning-color); }
.status-indicator.error \{ background: var(--error-color); }

\begin{lstlisting}
### 

#### 1. 

Web

\end{lstlisting}

css /\highlight{  }/ @media screen and (max-width: 768px) \{ .container
\{ padding: 15px; } /\highlight{ ...  ... }/ color: var(--gray-600); }
} }``

CSS3FlexboxGridUICSS3JavaScriptHTMLCSS

\hypertarget{javascript-1}{%
\section{4.4 JavaScript}\label{javascript-1}}

JavaScriptWebJavaScriptAPIJavaScript

JavaScriptES6+PromiseJavaScript

JavaScriptDOMJavaScriptVue.js

!!! info ``JavaScript''

\begin{verbatim}
JavaScriptJavaScriptJavaScript
\end{verbatim}

\hypertarget{javascript-2}{%
\subsection{JavaScript}\label{javascript-2}}

\hypertarget{javascript-3}{%
\subsubsection{JavaScript}\label{javascript-3}}

JavaScript\important{}JavaScriptNode.js

JavaScript - \important{} - \important{} - \important{} - \important{} -
\important{}

\hypertarget{javascriptweb}{%
\subsubsection{JavaScriptWeb}\label{javascriptweb}}

JavaScriptWeb

\begin{enumerate}
\def\labelenumi{\arabic{enumi}.}
\item
  \important{DOM}HTML
\item
  \important{}
\item
  \important{}
\end{enumerate}

JavaScript - - - -

\hypertarget{javascript-4}{%
\subsubsection{JavaScript}\label{javascript-4}}

\hypertarget{section-14}{%
\paragraph{1.}\label{section-14}}

JavaScript

\begin{lstlisting}[style=javascript]
// varES5
var oldWay = "";

// letES6+
let waterLevel = 85.5;
let stationName = "";

// constES6+
const MAX_WATER_LEVEL = 100;
const API_URL = "https://api.water-monitor.com";
\end{lstlisting}

\hypertarget{section-15}{%
\paragraph{2.}\label{section-15}}

JavaScript

\important{}
\begin{lstlisting}[style=javascript]
// 
let temperature = 25.5;
let stationCount = 10;

// 
let message = "";
let description = \code{${temperature}°C}; // 

// 
let isOnline = true;
let hasAlert = false;

// undefined
let undefinedValue;
console.log(undefinedValue); // undefined

// null
let emptyValue = null;

// SymbolES6+
let id = Symbol('id');

// BigIntES2020+
let bigNumber = 123456789012345678901234567890n;
\end{lstlisting}

\important{}
\begin{lstlisting}[style=javascript]
// 
let station = {
    id: 1,
    name: "",
    location: { lat: 34.5, lng: 112.5 },
    status: ""
};

// 
let waterLevels = [85.5, 86.2, 84.8, 87.1];
let stations = ["A", "B", "C"];

// 
function calculateAverage(values) {
    let sum = values.reduce((total, val) => total + val, 0);
    return sum / values.length;
}
\end{lstlisting}

\hypertarget{section-16}{%
\paragraph{3.}\label{section-16}}

\begin{lstlisting}[style=javascript]
// 
let a = 10, b = 3;
console.log(a + b); // 13 
console.log(a - b); // 7  
console.log(a \highlight{ b); // 30 
console.log(a / b); // 3.333... 
console.log(a % b); // 1 

// 
console.log(a > b);  // true
console.log(a === b); // false 
console.log(a == "10"); // true 

// 
let isOnline = true;
let hasData = false;
console.log(isOnline && hasData); // false 
console.log(isOnline || hasData); // true  
console.log(!isOnline); // false 
\end{lstlisting}

\hypertarget{section-17}{%
\paragraph{4.}\label{section-17}}

\begin{lstlisting}[style=javascript]
// 
let waterLevel = 85.5;

if (waterLevel > 90) {
    console.log("");
} else if (waterLevel < 20) {
    console.log("");
} else {
    console.log("");
}

// 
// for
for (let i = 0; i < stations.length; i++) {
    console.log(\code{${stations[i]}});
}

// while
let attempts = 0;
while (attempts < 3) {
    console.log(\code{${attempts + 1}});
    attempts++;
}

// for...of
for (let level of waterLevels) {
    console.log(\code{${level}});
}
\end{lstlisting}

\hypertarget{section-18}{%
\paragraph{5.}\label{section-18}}

\begin{lstlisting}[style=javascript]
// 
function checkWaterLevel(level) {
    if (level > 90) {
        return "";
    } else if (level > 70) {
        return "";
    } else {
        return "";
    }
}

// 
const calculateFlow = function(velocity, area) {
    return velocity } area;
};

// ES6+
const getTemperatureStatus = (temp) => {
    return temp > 30 ? "" : temp < 0 ? "" : "";
};

// 
const double = x => x \highlight{ 2;
const greet = () => "";
\end{lstlisting}

\hypertarget{javascript-5}{%
\subsubsection{JavaScript}\label{javascript-5}}

JavaScript\important{}

\begin{enumerate}
\def\labelenumi{\arabic{enumi}.}
\item
  \important{}
\item
  \important{}
\item
  \important{}
\end{enumerate}

\begin{lstlisting}[style=javascript]
// 
const SYSTEM_NAME = "";

function processData() {
    // 
    const data = "";
    
    function analyzeData() {
        // 
        console.log(\code{${SYSTEM_NAME}${data}});
    }
    
    analyzeData();
}
\end{lstlisting}

JavaScriptJavaScriptES6+JavaScript

\hypertarget{javascript-6}{%
\subsection{JavaScript}\label{javascript-6}}

\hypertarget{section-19}{%
\subsubsection{}\label{section-19}}

\hypertarget{number}{%
\paragraph{1. Number}\label{number}}

JavaScriptIEEE 754

\begin{lstlisting}[style=javascript]
// 
let decimal = 42;          // 
let binary = 0b101010;     // ES6+
let octal = 0o52;          // ES6+
let hex = 0x2A;            // 

// 
let infinity = Infinity;    // 
let negInfinity = -Infinity; // 
let notANumber = NaN;      // 

// 
let waterLevel = 85.67;
let temperature = -5.2;
let pressure = 1.01325e5;  // 101325

// 
console.log(0.1 + 0.2);              // 0.30000000000000004
console.log((0.1 + 0.2).toFixed(2)); // "0.30"
console.log(Math.round((0.1 + 0.2) } 100) / 100); // 0.3

// 
console.log(Number.isInteger(42));     // true
console.log(Number.isNaN(NaN));        // true
console.log(Number.isFinite(100));     // true
console.log(Number.isFinite(Infinity)); // false
\end{lstlisting}

\important{}
\begin{lstlisting}[style=javascript]
// 
function processWaterLevel(rawLevel) {
    // 
    if (!Number.isFinite(rawLevel)) {
        console.error("");
        return null;
    }
    
    // 
    const level = Math.round(rawLevel \highlight{ 100) / 100;
    
    // 
    if (level > 90) return { level, status: "" };
    if (level > 70) return { level, status: "" };
    return { level, status: "" };
}
\end{lstlisting}

\hypertarget{string}{%
\paragraph{2. String}\label{string}}

JavaScriptJavaScript

\hypertarget{section-20}{%
\subparagraph{}\label{section-20}}

JavaScript

\begin{lstlisting}[style=javascript]
// 
let str1 = "";  // 
let str2 = '';  // 
let str3 = \code{};    // ES6
\end{lstlisting}

\hypertarget{section-21}{%
\subparagraph{}\label{section-21}}

ES6(\textbackslash code\{)

\begin{lstlisting}[style=javascript]
let stationName = "";
let currentLevel = 85.5;

// 
let report = }

========
${stationName}
${currentLevel}
${currentLevel > 80 ? '' : ''}
${new Date().toLocaleString()}
\code{;

console.log(report);
\end{lstlisting}

\hypertarget{section-22}{%
\subparagraph{}\label{section-22}}

\begin{lstlisting}[style=javascript]
let text = "Smart Water Management System";

// 
console.log(text.length);        // 29

// 
console.log(text[0]);           // "S"
console.log(text[text.length - 1]); // "m" ()

// charAt
console.log(text.charAt(6));    // "W"
console.log(text.charAt(100));  // "" ()
\end{lstlisting}

\hypertarget{section-23}{%
\subparagraph{}\label{section-23}}

\begin{lstlisting}[style=javascript]
let text = "Smart Water Management System";

// 
console.log(text.indexOf("Water"));      // 6 ()
console.log(text.lastIndexOf("a"));      // 24 ()
console.log(text.indexOf("River"));      // -1 (-1)

// 
console.log(text.includes("Smart"));     // true ()
console.log(text.startsWith("Smart"));   // true ()
console.log(text.endsWith("System"));    // true ()
\end{lstlisting}

\hypertarget{section-24}{%
\subparagraph{}\label{section-24}}

\begin{lstlisting}[style=javascript]
let text = "Smart Water Management System";

// slice
console.log(text.slice(6, 11));         // "Water"
console.log(text.slice(-6));            // "System" (6)
console.log(text.slice(6, -7));         // "Water Management"

// substring
console.log(text.substring(6, 11));     // "Water"
console.log(text.substring(11, 6));     // "Water" ()

// substr
console.log(text.substr(6, 5));         // "Water"
\end{lstlisting}

\hypertarget{section-25}{%
\subparagraph{}\label{section-25}}

\begin{lstlisting}[style=javascript]
let text = "Smart Water Management System";

// 
console.log(text.toLowerCase());        // "smart water management system"
console.log(text.toUpperCase());        // "SMART WATER MANAGEMENT SYSTEM"

// 
let spacedText = "    ";
console.log(spacedText.trim());         // ""
console.log(spacedText.trimStart());    // "  "
console.log(spacedText.trimEnd());      // "  "
\end{lstlisting}

\hypertarget{section-26}{%
\subparagraph{}\label{section-26}}

CSV

\begin{lstlisting}[style=javascript]
let text = "Smart Water Management System";

// 
let parts = text.split(" ");            // ["Smart", "Water", "Management", "System"]
let limited = text.split(" ", 2);       // ["Smart", "Water"] ()

// 
let joined = parts.join("-");           // "Smart-Water-Management-System"
let withComma = parts.join(", ");       // "Smart, Water, Management, System"
\end{lstlisting}

\hypertarget{section-27}{%
\subparagraph{}\label{section-27}}

\begin{lstlisting}[style=javascript]
let text = "Smart Water Management System";

// 
let updated = text.replace("Smart", "Intelligent");
console.log(updated); // "Intelligent Water Management System"

// 
let globalReplace = text.replace(/a/g, "@");  // 'a'
console.log(globalReplace); // "Sm@rt W@ter M@n@gement System"

// 
let capitalized = text.replace(/\b\w+\b/g, function(word) {
    return word.toUpperCase();
});
console.log(capitalized); // "SMART WATER MANAGEMENT SYSTEM"
\end{lstlisting}

\hypertarget{section-28}{%
\subparagraph{}\label{section-28}}

\important{}

\begin{lstlisting}[style=javascript]
// 
function parseStationData(dataString) {
    // : "ST001||85.5||2023-12-01 10:30:00"
    const parts = dataString.split("|");
    
    // 
    if (parts.length !== 5) {
        throw new Error(}5${parts.length}\code{);
    }
    
    return {
        id: parts[0].trim(),
        name: parts[1].trim(),
        waterLevel: parseFloat(parts[2]),
        status: parts[3].trim(),
        timestamp: new Date(parts[4].trim())
    };
}

// 
try {
    const rawData = "ST001||85.5||2023-12-01 10:30:00";
    const stationData = parseStationData(rawData);
    console.log(":", stationData);
} catch (error) {
    console.error(":", error.message);
}
\end{lstlisting}

\important{}

\begin{lstlisting}[style=javascript]
// 
function generateReportTitle(stationName, date, reportType = "") {
    // 
    const formattedDate = date.toISOString().slice(0, 10);
    
    // 
    const cleanStationName = stationName.replace(/[^\w\u4e00-\u9fa5]/g, "_");
    
    // 
    const fileName = }${cleanStationName}_${reportType}_${formattedDate}\code{.replace(/\s+/g, "_");
    
    return {
        title: }${stationName} ${reportType}\code{,
        fileName: fileName + ".pdf",
        displayName: }${stationName} - ${formattedDate}\code{
    };
}

// 
const reportInfo = generateReportTitle("", new Date());
console.log(":", reportInfo);
// : { 
//   title: " ",
//   fileName: "__2023-12-01.pdf",
//   displayName: " - 2023-12-01"
// }
\end{lstlisting}

\important{}

\begin{lstlisting}[style=javascript]
// 
function validateStationName(name) {
    // 
    const trimmed = name.trim();
    
    // 
    if (!trimmed) {
        return { isValid: false, error: "" };
    }
    
    // 
    if (trimmed.length < 2 || trimmed.length > 50) {
        return { isValid: false, error: "2-50" };
    }
    
    // 
    const invalidChars = /[<>:"/\\|?}]/;
    if (invalidChars.test(trimmed)) {
        return { isValid: false, error: "" };
    }
    
    return { isValid: true, value: trimmed };
}

// 
console.log(validateStationName("  01  ")); 
// : { isValid: true, value: "01" }

console.log(validateStationName("<>"));
// : { isValid: false, error: "" }
\end{lstlisting}

\hypertarget{boolean}{%
\paragraph{3. Boolean}\label{boolean}}

JavaScript}true\code{}false\code{JavaScript

##### 

\begin{lstlisting}[style=javascript]
// 
let isOnline = true;          // 
let hasError = false;         // 
let alertEnabled = true;      // 
let maintenanceMode = false;  // 
\end{lstlisting}

##### AND (&&) - 

AND

\begin{lstlisting}[style=javascript]
let a = true, b = false;

// 
console.log(a && b);        // false ()
console.log(true && true);  // true

// 
console.log(true && "hello");   // "hello" ()
console.log(false && "hello");  // false ()

// 
let user = { isAdmin: true, name: "" };
user.isAdmin && console.log(}\$\{user.name}\code{); // 
\end{lstlisting}

##### OR (||) -   

OR

\begin{lstlisting}[style=javascript]
let a = true, b = false;

// 
console.log(a || b);        // true ()
console.log(false || false); // false

// 
console.log(false || "default"); // "default" ()
console.log(true || "backup");   // true ()

// 
function createConfig(options) {
    return {
        host: options.host || "localhost",
        port: options.port || 8080,
        timeout: options.timeout || 30000
    };
}
\end{lstlisting}

##### NOT (!) - 

NOT

\begin{lstlisting}[style=javascript]
let isOnline = true;

// 
console.log(!isOnline);     // false
console.log(!false);        // true

// 
console.log(!!"hello");     // true (true)
console.log(!!0);          // false (0false)
console.log(!!null);       // false (nullfalse)

// 
function toggleMaintenanceMode(currentMode) {
    return !currentMode;  // 
}
\end{lstlisting}

#####  (??) - ES2020

}null\code{}undefined\code{

\begin{lstlisting}[style=javascript]
let userInput = null;
let defaultValue = "";

// 
console.log(userInput ?? defaultValue); // ""

// 
console.log(0 || "default");    // "default" (0)
console.log(0 ?? "default");    // 0 (0nullundefined)

console.log("" || "default");   // "default" ()
console.log("" ?? "default");   // "" (nullundefined)
\end{lstlisting}

##### 

JavaScript

\begin{lstlisting}[style=javascript]
// Falsy- false
console.log(Boolean(false));      // false
console.log(Boolean(0));          // false
console.log(Boolean(-0));         // false
console.log(Boolean(0n));         // false (BigInt)
console.log(Boolean(""));         // false ()
console.log(Boolean(null));       // false
console.log(Boolean(undefined));  // false
console.log(Boolean(NaN));        // false

// Truthy- true
console.log(Boolean(1));          // true
console.log(Boolean(-1));         // true
console.log(Boolean("hello"));    // true
console.log(Boolean(" "));        // true ()
console.log(Boolean([]));         // true ()
console.log(Boolean({}));         // true ()
console.log(Boolean(function(){})); // true ()
\end{lstlisting}

##### 

\important{}

\begin{lstlisting}[style=javascript]
// 
function checkStationStatus(station) {
    const isOnline = station.status === "";
    const hasRecentData = station.lastUpdate > Date.now() - 300000; // 5
    const isLevelNormal = station.waterLevel >= 20 && station.waterLevel <= 90;
    const isTempNormal = station.temperature >= -10 && station.temperature <= 50;
    
    // 
    const isHealthy = isOnline && hasRecentData && isLevelNormal && isTempNormal;
    
    // 
    const issues = [
        !isOnline && "",
        !hasRecentData && "", 
        !isLevelNormal && "",
        !isTempNormal && ""
    ].filter(Boolean); // false
    
    return {
        isHealthy,
        issues,
        statusLevel: issues.length === 0 ? "" : 
                    issues.length <= 2 ? "" : ""
    };
}

// 
const station = {
    id: "ST001",
    name: "",
    status: "",
    waterLevel: 95,  // 
    temperature: 25,
    lastUpdate: Date.now() - 100000 // 140
};

const statusResult = checkStationStatus(station);
console.log(":", statusResult);
// : {
//   isHealthy: false,
//   issues: [""],
//   statusLevel: ""
// }
\end{lstlisting}

\important{}

\begin{lstlisting}[style=javascript]
// 
function createMonitoringConfig(userConfig = {}) {
    // 
    const config = {
        // 
        refreshInterval: userConfig.refreshInterval ?? 30000,    // ()
        alertEnabled: userConfig.alertEnabled ?? true,          // 
        maxRetries: userConfig.maxRetries ?? 3,                 // 
        
        //  - 
        apiUrl: userConfig.apiUrl || "https://api.water-monitor.com",
        timeout: userConfig.timeout || 10000,
        
        // 
        waterLevelThresholds: {
            low: userConfig.waterLevelThresholds?.low ?? 20,
            high: userConfig.waterLevelThresholds?.high ?? 90,
            danger: userConfig.waterLevelThresholds?.danger ?? 95
        },
        
        //  - 
        features: {
            realTimeChart: !!(userConfig.features?.realTimeChart ?? true),
            dataExport: !!(userConfig.features?.dataExport ?? false),
            alertHistory: !!(userConfig.features?.alertHistory ?? true)
        }
    };
    
    return config;
}

// 
const customConfig = {
    refreshInterval: 60000,
    waterLevelThresholds: { high: 85 },
    features: { dataExport: true }
};

const finalConfig = createMonitoringConfig(customConfig);
console.log(":", finalConfig);
\end{lstlisting}

\important{}

\begin{lstlisting}[style=javascript]
// 
function checkUserPermission(user, action, resource) {
    // 
    const isLoggedIn = user && user.isAuthenticated;
    const hasRole = user?.role && user.role !== "guest";
    const isActive = user?.status === "active";
    
    // false
    if (!isLoggedIn || !hasRole || !isActive) {
        return {
            allowed: false,
            reason: ""
        };
    }
    
    // 
    const isAdmin = user.role === "admin";
    const isSupervisor = user.role === "supervisor";
    const isOperator = user.role === "operator";
    
    // 
    let hasPermission = false;
    
    if (action === "read") {
        // 
        hasPermission = isLoggedIn;
    } else if (action === "write") {
        // 
        hasPermission = isAdmin || isOperator;
    } else if (action === "delete") {
        // 
        hasPermission = isAdmin;
    } else if (action === "configure") {
        // 
        hasPermission = isAdmin || isSupervisor;
    }
    
    return {
        allowed: hasPermission,
        reason: hasPermission ? "" : }${user.role}${action}\code{
    };
}

// 
const user = {
    id: "user001",
    name: "",
    role: "operator",
    isAuthenticated: true,
    status: "active"
};

const readPermission = checkUserPermission(user, "read", "monitoring-data");
const deletePermission = checkUserPermission(user, "delete", "station-config");

console.log(":", readPermission);  // { allowed: true, reason: "" }
console.log(":", deletePermission); // { allowed: false, reason: "operatordelete" }
\end{lstlisting}

#### 4. 

JavaScriptJavaScript

##### 

\important{}

\begin{lstlisting}[style=javascript]
let num = 42;
let bool = true;
let obj = { name: "" };

// String()
console.log(String(num));        // "42"
console.log(String(bool));       // "true" 
console.log(String(obj));        // "[object Object]"

// toString()
console.log(num.toString());     // "42"
console.log(bool.toString());    // "true"
// nullundefinedtoString()

// 
console.log(}\$\{num}\code{);      // "42"
console.log(num + "");          // "42"
\end{lstlisting}

\important{}

API

\begin{lstlisting}[style=javascript]
let str = "42";
let floatStr = "42.5";
let invalidStr = "42px";

// Number()
console.log(Number(str));        // 42
console.log(Number(floatStr));   // 42.5
console.log(Number(invalidStr)); // NaN
console.log(Number(""));         // 0 (0)
console.log(Number("  "));       // 0 (0)

// parseInt()
console.log(parseInt(str));      // 42
console.log(parseInt(floatStr)); // 42 ()
console.log(parseInt(invalidStr)); // 42 ()
console.log(parseInt("42.8px")); // 42

// parseFloat()  
console.log(parseFloat(floatStr)); // 42.5
console.log(parseFloat(invalidStr)); // 42 ()

// 
console.log(+str);               // 42 ()
console.log(+"42.5");            // 42.5
\end{lstlisting}

\important{}

\begin{lstlisting}[style=javascript]
// Boolean()
console.log(Boolean(1));         // true
console.log(Boolean(0));         // false
console.log(Boolean(""));        // false
console.log(Boolean(" "));       // true ()
console.log(Boolean([]));        // true (true)
console.log(Boolean({}));        // true (true)
console.log(Boolean(null));      // false
console.log(Boolean(undefined)); // false

// 
console.log(!!"hello");          // true
console.log(!!0);               // false
console.log(!!"");              // false
\end{lstlisting}

##### 

JavaScript

\important{}

\begin{lstlisting}[style=javascript]
// 
console.log(5 + "3");            // "53" ()
console.log("5" + 3);            // "53" ()
console.log(5 + 3 + "2");        // "82" (5+3=8"2")
console.log("2" + 5 + 3);        // "253" ()

// 
let level = 85.5;
console.log(}\$\{level}\code{);  // "85.5" ()
\end{lstlisting}

\important{}  

\begin{lstlisting}[style=javascript]
console.log("5" \highlight{ 3);            // 15 ("5"5)
console.log("5" - 3);            // 2 ("5"5)
console.log("5" / "2");          // 2.5 ()
console.log(true + 1);           // 2 (true1)
console.log(false } 5);          // 0 (false0)

// 
console.log("10" > 5);           // true ("10"10)
console.log("10" > "5");         // false (Unicode)
\end{lstlisting}

\important{}

JavaScript

\begin{lstlisting}[style=javascript]
// if
if ("") {
    console.log(""); // false
}

if ("hello") {
    console.log("");   // true
}

// 
console.log("" && "hello");      // "" ()
console.log("hi" && "hello");    // "hello" ()
console.log("" || "hello");      // "hello" ()
\end{lstlisting}

##### 

\important{}

\begin{lstlisting}[style=javascript]
// 
function safeToNumber(value, defaultValue = 0) {
    // nullundefined
    if (value === null || value === undefined) {
        return defaultValue;
    }
    
    // 
    if (typeof value === 'number') {
        return Number.isNaN(value) ? defaultValue : value;
    }
    
    // 
    const num = Number(value);
    
    // 
    return Number.isNaN(num) ? defaultValue : num;
}

// 
function safeToString(value, defaultValue = "") {
    if (value === null || value === undefined) {
        return defaultValue;
    }
    
    return String(value);
}

// 
console.log(safeToNumber("42.5"));      // 42.5
console.log(safeToNumber("invalid"));   // 0
console.log(safeToNumber(null, -1));    // -1
console.log(safeToString(null, "N/A")); // "N/A"
\end{lstlisting}

\important{}

\begin{lstlisting}[style=javascript]
// API
function processMonitoringData(rawData) {
    try {
        // 
        const data = typeof rawData === 'string' ? JSON.parse(rawData) : rawData;
        
        if (!data || typeof data !== 'object') {
            throw new Error("");
        }
        
        // 
        const processedData = {
            // ID
            stationId: safeToString(data.stationId || data.id || data.station_id).trim(),
            
            // 
            waterLevel: safeToNumber(data.waterLevel || data.water_level || data.level),
            temperature: safeToNumber(data.temperature || data.temp, null),
            pressure: safeToNumber(data.pressure, null),
            flow: safeToNumber(data.flow || data.flowRate, null),
            
            // 
            isActive: !!(data.isActive || data.active || data.status === "active"),
            status: safeToString(data.status, "unknown").toLowerCase(),
            
            // Date
            timestamp: data.timestamp ? new Date(data.timestamp) : new Date(),
            
            // 
            location: data.location ? {
                lat: safeToNumber(data.location.lat || data.location.latitude),
                lng: safeToNumber(data.location.lng || data.location.longitude)
            } : null
        };
        
        return {
            success: true,
            data: processedData,
            errors: []
        };
        
    } catch (error) {
        return {
            success: false,
            data: null,
            errors: [}: ${error.message}\code{]
        };
    }
}

// 
const rawApiData = {
    station_id: "ST001",
    water_level: "85.5",
    temp: "25.2",
    active: "true",
    timestamp: "2023-12-01T10:30:00Z"
};

const result = processMonitoringData(rawApiData);
console.log(":", result);
\end{lstlisting}

\important{}

\begin{lstlisting}[style=javascript]
// 
function validateWaterLevel(input) {
    // 
    const trimmed = String(input).trim();
    
    // 
    if (!trimmed) {
        return { 
            isValid: false, 
            error: "",
            value: null 
        };
    }
    
    // 
    const level = parseFloat(trimmed);
    
    // 
    if (Number.isNaN(level)) {
        return { 
            isValid: false, 
            error: "",
            value: null 
        };
    }
    
    // 
    if (level < -50 || level > 300) {
        return { 
            isValid: false, 
            error: "-50300",
            value: null 
        };
    }
    
    // 
    const roundedLevel = Math.round(level \highlight{ 100) / 100;
    
    return { 
        isValid: true, 
        value: roundedLevel,
        displayValue: }${roundedLevel}\code{
    };
}

// 
function validateStationConfig(configInput) {
    const errors = [];
    const config = {};
    
    // 
    const nameResult = validateStationName(configInput.name);
    if (!nameResult.isValid) {
        errors.push(}: ${nameResult.error}\code{);
    } else {
        config.name = nameResult.value;
    }
    
    // 
    const thresholds = {};
    ['low', 'high', 'danger'].forEach(key => {
        const input = configInput[}${key}Threshold\code{];
        if (input !== undefined && input !== '') {
            const result = validateWaterLevel(input);
            if (!result.isValid) {
                errors.push(}${key}: ${result.error}\code{);
            } else {
                thresholds[key] = result.value;
            }
        }
    });
    
    // 
    const intervalInput = configInput.refreshInterval;
    if (intervalInput !== undefined) {
        const interval = safeToNumber(intervalInput, 0);
        if (interval < 1000 || interval > 300000) {
            errors.push("1-300");
        } else {
            config.refreshInterval = interval;
        }
    }
    
    return {
        isValid: errors.length === 0,
        errors,
        config: errors.length === 0 ? { ...config, thresholds } : null
    };
}

// 
console.log(validateWaterLevel("85.67"));  
// { isValid: true, value: 85.67, displayValue: "85.67" }

console.log(validateWaterLevel("invalid")); 
// { isValid: false, error: "", value: null }
\end{lstlisting}

### 

JavaScript

#### 1. 

##### 

JavaScript

\begin{lstlisting}[style=javascript]
let a = 10, b = 3;

// 
console.log(a + b);    // 13 
console.log(a - b);    // 7  
console.log(a } b);    // 30 
console.log(a / b);    // 3.333... 

// 
console.log(a % b);    // 1 (103)
console.log(10 % 4);   // 2
console.log(15 % 5);   // 0 (0)

// ES2016+
console.log(a ** b);   // 1000 (103)
console.log(2 ** 8);   // 256 (28)
console.log(9 ** 0.5); // 3 ()
\end{lstlisting}

##### 

\begin{lstlisting}[style=javascript]
let x = 5;

// ++variable
console.log(++x);      // 6 (x66)
console.log(x);        // 6

// variable++  
let y = 5;
console.log(y++);      // 5 (5y6)
console.log(y);        // 6

// --variable
let m = 5;
console.log(--m);      // 4 (m44)

// variable--
let n = 5;
console.log(n--);      // 5 (5n4)
console.log(n);        // 4
\end{lstlisting}

##### 

\begin{lstlisting}[style=javascript]
let str = "42";
let bool = true;

// 
console.log(+str);     // 42 ()
console.log(+bool);    // 1 (true1)
console.log(+false);   // 0 (false0)
console.log(+"");      // 0 (0)

// 
console.log(-str);     // -42 ()
console.log(-true);    // -1
console.log(-false);   // -0
\end{lstlisting}

##### 

\important{}

\begin{lstlisting}[style=javascript]
// Q = A × V ( =  × )
function calculateFlow(crossSectionArea, velocity) {
    // 
    if (crossSectionArea <= 0 || velocity < 0) {
        throw new Error("0");
    }
    
    const flow = crossSectionArea \highlight{ velocity;
    
    return {
        flow: Math.round(flow } 1000) / 1000, // 3
        unit: "m³/s",
        formula: }${crossSectionArea} × ${velocity} = ${flow.toFixed(3)}\code{
    };
}

// 
function calculateTrapezoidalFlow(bottomWidth, depth, sideSlope, velocity) {
    // A = (b + m\highlight{h) } h
    // b: , h: , m: 
    const area = (bottomWidth + sideSlope \highlight{ depth) } depth;
    
    return calculateFlow(area, velocity);
}

// 
const flowResult = calculateFlow(25.5, 1.8);
console.log(":", flowResult);
// : { flow: 45.9, unit: "m³/s", formula: "25.5 × 1.8 = 45.900" }
\end{lstlisting}

\important{}

\begin{lstlisting}[style=javascript]
// 
function calculateWaterLevelRate(currentLevel, previousLevel, timeInterval) {
    // timeInterval
    const levelChange = currentLevel - previousLevel;
    const timeHours = timeInterval / (1000 \highlight{ 60 } 60); // 
    
    if (timeHours === 0) {
        return { rate: 0, trend: "" };
    }
    
    const ratePerHour = levelChange / timeHours;
    
    // 
    let trend;
    if (Math.abs(ratePerHour) < 0.01) {
        trend = "";
    } else if (ratePerHour > 0) {
        trend = ratePerHour > 0.1 ? "" : "";
    } else {
        trend = ratePerHour < -0.1 ? "" : "";
    }
    
    return {
        rate: Math.round(ratePerHour \highlight{ 1000) / 1000, // 3
        unit: "/",
        trend,
        timeSpan: }${timeHours.toFixed(2)}\code{
    };
}

// 
const rateResult = calculateWaterLevelRate(85.5, 84.8, 2 } 60 \highlight{ 60 } 1000); // 2
console.log(":", rateResult);
// : { rate: 0.35, unit: "/", trend: "", timeSpan: "2.00" }
\end{lstlisting}

\important{}

\begin{lstlisting}[style=javascript]
// 
function calculateAverage(values, weights = null) {
    if (!values || values.length === 0) {
        return { average: 0, count: 0 };
    }
    
    // 
    const validValues = values.filter(val => typeof val === 'number' && !isNaN(val));
    
    if (validValues.length === 0) {
        return { average: 0, count: 0 };
    }
    
    let sum, totalWeight;
    
    if (weights && weights.length === validValues.length) {
        // 
        sum = validValues.reduce((acc, val, index) => acc + val * weights[index], 0);
        totalWeight = weights.reduce((acc, weight) => acc + weight, 0);
        
        return {
            average: sum / totalWeight,
            count: validValues.length,
            type: "weighted"
        };
    } else {
        // 
        sum = validValues.reduce((acc, val) => acc + val, 0);
        
        return {
            average: sum / validValues.length,
            count: validValues.length,
            type: "arithmetic"
        };
    }
}

// 
function calculateVarianceAndStdDev(values) {
    const avgResult = calculateAverage(values);
    if (avgResult.count < 2) {
        return { variance: 0, standardDeviation: 0, count: avgResult.count };
    }
    
    const mean = avgResult.average;
    const validValues = values.filter(val => typeof val === 'number' && !isNaN(val));
    
    // 
    const squaredDiffs = validValues.map(val => (val - mean) ** 2);
    const variance = squaredDiffs.reduce((acc, val) => acc + val, 0) / (validValues.length - 1);
    
    return {
        variance: Math.round(variance \highlight{ 1000) / 1000,
        standardDeviation: Math.round(Math.sqrt(variance) } 1000) / 1000,
        count: validValues.length,
        mean: Math.round(mean \highlight{ 1000) / 1000
    };
}

// 
const waterLevels = [85.2, 84.8, 86.1, 85.5, 84.9, 85.8, 86.2];
const avgResult = calculateAverage(waterLevels);
const statsResult = calculateVarianceAndStdDev(waterLevels);

console.log(":", avgResult);
console.log(":", statsResult);
\end{lstlisting}

\important{}

\begin{lstlisting}[style=javascript]
// 
function roundToPrecision(number, precision = 2) {
    const factor = Math.pow(10, precision);
    return Math.round(number } factor) / factor;
}

// 
function processCalculationResult(result, precision = 3) {
    if (typeof result !== 'number' || isNaN(result)) {
        return { value: 0, displayValue: "" };
    }
    
    const rounded = roundToPrecision(result, precision);
    
    return {
        value: rounded,
        displayValue: rounded.toFixed(precision),
        scientific: result.toExponential(precision)
    };
}

// 
const calculation = 85.23456789 \highlight{ 1.41421356;
const processed = processCalculationResult(calculation);
console.log(":", processed);
// : { value: 120.563, displayValue: "120.563", scientific: "1.206e+2" }
\end{lstlisting}

#### 2. 

##### 

\begin{lstlisting}[style=javascript]
let waterLevel1 = 85.5;
let waterLevel2 = 90.0;
let threshold = 80;

// 
console.log(waterLevel1 < waterLevel2);   // true (85.5 < 90.0)
console.log(waterLevel1 > threshold);     // true (85.5 > 80)
console.log(waterLevel2 <= 90);          // true (90.0 <= 90)
console.log(waterLevel1 >= threshold);    // true (85.5 >= 80)

// Unicode
console.log("apple" < "banana");          // true
console.log("Apple" < "banana");          // true (Unicode)
console.log("10" < "9");                  // true ()

// 
console.log("85" > 80);                   // true ("85"85)
console.log(true > false);                // true (true1false0)
\end{lstlisting}

##### 

JavaScript=====

\begin{lstlisting}[style=javascript]
let x = 5, y = 10, z = "5";

// ===- 
console.log(x === z);     // false (5"5")
console.log(x === 5);     // true ()
console.log(null === undefined); // false ()

// !==
console.log(x !== z);     // true ()
console.log(x !== 5);     // false ()

// ==- 
console.log(x == z);      // true ("5"5)
console.log(true == 1);   // true (true1)
console.log(false == 0);  // true (false0)
console.log(null == undefined); // true ()

// !=
console.log(x != z);      // false ()
\end{lstlisting}

##### 

\begin{lstlisting}[style=javascript]
// NaN
console.log(NaN === NaN);          // false (NaN)
console.log(NaN == NaN);           // false
console.log(Number.isNaN(NaN));    // true (NaN)

// ES6
console.log(Object.is(NaN, NaN));  // true
console.log(Object.is(+0, -0));    // false ()
console.log(+0 === -0);            // true ()

// nullundefined
console.log(null == undefined);    // true ()
console.log(null === undefined);   // false ()
console.log(null == 0);            // false
console.log(undefined == 0);       // false
\end{lstlisting}

##### 

\begin{lstlisting}[style=javascript]
// 
let station1 = { name: "A", level: 85 };
let station2 = { name: "A", level: 85 };
let station3 = station1;

console.log(station1 === station2); // false ()
console.log(station1 === station3); // true ()

// 
let arr1 = [1, 2, 3];
let arr2 = [1, 2, 3];
console.log(arr1 === arr2);         // false ()

// 
let str1 = "hello";
let str2 = "hello";
console.log(str1 === str2);         // true ()
\end{lstlisting}

##### 

\important{}

\begin{lstlisting}[style=javascript]
// 
const WATER_LEVEL_THRESHOLDS = {
    DANGER: 95,      // 
    WARNING: 80,     // 
    NORMAL_HIGH: 70, // 
    NORMAL_LOW: 20,  // 
    LOW: 10          // 
};

// 
function classifyWaterLevel(level) {
    // 
    if (typeof level !== 'number' || isNaN(level)) {
        return { 
            level: "invalid", 
            message: "",
            color: "gray",
            priority: 0 
        };
    }
    
    // 
    if (level >= WATER_LEVEL_THRESHOLDS.DANGER) {
        return { 
            level: "danger", 
            message: }${level}\code{,
            color: "red",
            priority: 5 
        };
    } else if (level >= WATER_LEVEL_THRESHOLDS.WARNING) {
        return { 
            level: "warning", 
            message: }${level}\code{,
            color: "orange", 
            priority: 4
        };
    } else if (level >= WATER_LEVEL_THRESHOLDS.NORMAL_HIGH) {
        return { 
            level: "normal-high", 
            message: }${level}\code{,
            color: "yellow",
            priority: 2 
        };
    } else if (level >= WATER_LEVEL_THRESHOLDS.NORMAL_LOW) {
        return { 
            level: "normal", 
            message: }${level}\code{,
            color: "green",
            priority: 1 
        };
    } else if (level >= WATER_LEVEL_THRESHOLDS.LOW) {
        return { 
            level: "low", 
            message: }${level}\code{,
            color: "blue",
            priority: 2 
        };
    } else {
        return { 
            level: "very-low", 
            message: }${level}\code{,
            color: "purple",
            priority: 3 
        };
    }
}

// 
function analyzeMultipleStations(stations) {
    return stations.map(station => {
        const classification = classifyWaterLevel(station.waterLevel);
        return {
            ...station,
            ...classification,
            needsAttention: classification.priority >= 3
        };
    }).sort((a, b) => b.priority - a.priority); // 
}

// 
const stations = [
    { id: "ST001", name: "", waterLevel: 96.5 },
    { id: "ST002", name: "", waterLevel: 75.2 },
    { id: "ST003", name: "", waterLevel: 15.8 },
];

const analysisResult = analyzeMultipleStations(stations);
console.log(":", analysisResult);
\end{lstlisting}

\important{}

\begin{lstlisting}[style=javascript]
// 
function checkDataConsistency(currentData, historicalData) {
    const issues = [];
    
    // ID
    if (currentData.stationId !== historicalData.stationId) {
        issues.push({
            type: "id_mismatch",
            message: "ID",
            severity: "high"
        });
    }
    
    // 
    const timeDiff = Math.abs(currentData.timestamp - historicalData.timestamp);
    const maxReasonableGap = 24 } 60 \highlight{ 60 } 1000; // 24
    
    if (timeDiff > maxReasonableGap) {
        issues.push({
            type: "time_gap",
            message: }${(timeDiff / (60 \highlight{ 60 } 1000)).toFixed(1)}\code{,
            severity: "medium"
        });
    }
    
    // 
    const levelDiff = Math.abs(currentData.waterLevel - historicalData.waterLevel);
    const maxReasonableChange = 10; // 10
    
    if (levelDiff > maxReasonableChange) {
        issues.push({
            type: "level_jump",
            message: }${levelDiff.toFixed(2)}\code{,
            severity: "high"
        });
    }
    
    // 
    const minReasonableLevel = -50;
    const maxReasonableLevel = 300;
    
    if (currentData.waterLevel < minReasonableLevel || 
        currentData.waterLevel > maxReasonableLevel) {
        issues.push({
            type: "value_range",
            message: }${currentData.waterLevel}\code{,
            severity: "high"
        });
    }
    
    return {
        isConsistent: issues.length === 0,
        issues,
        riskLevel: issues.some(issue => issue.severity === "high") ? "high" : 
                  issues.some(issue => issue.severity === "medium") ? "medium" : "low"
    };
}

// 
const currentData = {
    stationId: "ST001",
    waterLevel: 85.5,
    timestamp: Date.now()
};

const historicalData = {
    stationId: "ST001", 
    waterLevel: 84.2,
    timestamp: Date.now() - 30 \highlight{ 60 } 1000 // 30
};

const consistencyResult = checkDataConsistency(currentData, historicalData);
console.log(":", consistencyResult);
\end{lstlisting}

\important{}

\begin{lstlisting}[style=javascript]
// 
function calculateStationScore(station) {
    let score = 100; // 100
    
    // 
    const levelClass = classifyWaterLevel(station.waterLevel);
    switch(levelClass.level) {
        case "danger":
            score -= 40;
            break;
        case "warning": 
            score -= 20;
            break;
        case "very-low":
            score -= 15;
            break;
        case "low":
            score -= 5;
            break;
    }
    
    // 
    if (station.status !== "online") {
        score -= 30;
    }
    
    // 
    const updateAge = Date.now() - new Date(station.lastUpdate).getTime();
    const ageHours = updateAge / (60 \highlight{ 60 } 1000);
    
    if (ageHours > 24) {
        score -= 25;
    } else if (ageHours > 2) {
        score -= 10;
    } else if (ageHours > 0.5) {
        score -= 5;
    }
    
    return Math.max(0, score); // 0
}

// 
function sortStationsByPriority(stations, criteria = "comprehensive") {
    return [...stations].sort((a, b) => {
        switch(criteria) {
            case "waterLevel":
                return b.waterLevel - a.waterLevel; // 
                
            case "priority":
                const priorityA = classifyWaterLevel(a.waterLevel).priority;
                const priorityB = classifyWaterLevel(b.waterLevel).priority;
                return priorityB - priorityA; // 
                
            case "score":
                const scoreA = calculateStationScore(a);
                const scoreB = calculateStationScore(b);
                return scoreB - scoreA; // 
                
            case "comprehensive":
            default:
                //  >  > 
                const priorityDiff = classifyWaterLevel(b.waterLevel).priority - 
                                   classifyWaterLevel(a.waterLevel).priority;
                if (priorityDiff !== 0) return priorityDiff;
                
                const scoreDiff = calculateStationScore(b) - calculateStationScore(a);
                if (scoreDiff !== 0) return scoreDiff;
                
                return b.waterLevel - a.waterLevel;
        }
    });
}

// 
const monitoringStations = [
    { 
        id: "ST001", 
        name: "", 
        waterLevel: 96.5, 
        status: "online",
        lastUpdate: new Date(Date.now() - 10 \highlight{ 60 } 1000) // 10
    },
    { 
        id: "ST002", 
        name: "", 
        waterLevel: 75.2, 
        status: "offline",
        lastUpdate: new Date(Date.now() - 3 \highlight{ 60 } 60 \highlight{ 1000) // 3
    }
];

const sortedStations = sortStationsByPriority(monitoringStations, "comprehensive");
console.log(":", sortedStations);
\end{lstlisting}

### 

JavaScript

#### 1. 

##### if...else 

\begin{lstlisting}[style=javascript]
// if...else
function checkWaterLevelAlert(waterLevel) {
    if (waterLevel > 95) {
        console.log(" ");
        return "emergency";
    } else if (waterLevel > 85) {
        console.log("\faExclamationTriangle ");
        return "warning"; 
    } else if (waterLevel > 70) {
        console.log("ℹ ");
        return "notice";
    } else if (waterLevel < 20) {
        console.log("\faExclamationTriangle ");
        return "low";
    } else {
        console.log("\faCheck ");
        return "normal";
    }
}

// 
const alertLevel = checkWaterLevelAlert(87.5);
console.log(":", alertLevel); // warning
\end{lstlisting}

##### 

\begin{lstlisting}[style=javascript]
// 
const getStatusColor = (level) => level > 80 ? "red" : level > 50 ? "yellow" : "green";
const getStatusIcon = (isOnline) => isOnline ? "🟢" : "";

// 
const temperature = 25;
const weatherDescription = temperature > 30 ? "" : 
                          temperature > 20 ? "" : 
                          temperature > 10 ? "" : "";

// 
function displayStationStatus(station) {
    const statusText = station.isOnline ? }${station.name}\code{ : }${station.name}\code{;
    const levelStatus = station.waterLevel > 80 ? "" : "";
    
    return }${statusText} - : ${levelStatus}\code{;
}
\end{lstlisting}

##### switch

switchif...else

\begin{lstlisting}[style=javascript]
// switch
function getSeasonByMonth(month) {
    switch (month) {
        case 12:
        case 1:
        case 2:
            return "";
        case 3:
        case 4:
        case 5:
            return "";
        case 6:
        case 7:
        case 8:
            return "";
        case 9:
        case 10:
        case 11:
            return "";
        default:
            return "";
    }
}

// switch
function getDeviceActionByStatus(status) {
    switch (status.toLowerCase()) {
        case "online":
            return { action: "monitor", message: "", color: "green" };
        case "offline":
            return { action: "reconnect", message: "", color: "red" };
        case "maintenance":
            return { action: "wait", message: "", color: "orange" };
        case "error":
            return { action: "diagnose", message: "", color: "red" };
        default:
            return { action: "unknown", message: "", color: "gray" };
    }
}
\end{lstlisting}

#####  - 

switch

\begin{lstlisting}[style=javascript]
// switch
const alertLevelMap = {
    high: "danger",
    medium: "warning", 
    low: "normal",
    unknown: "undefined"
};

const getAlertLevel = (value) => alertLevelMap[value] || "unknown";

// 
const stationTypeConfig = {
    river: { 
        icon: "", 
        maxLevel: 100, 
        checkInterval: 300000  // 5
    },
    reservoir: { 
        icon: "", 
        maxLevel: 200, 
        checkInterval: 600000  // 10
    },
    lake: { 
        icon: "", 
        maxLevel: 150, 
        checkInterval: 900000  // 15
    }
};

function getStationConfig(type) {
    return stationTypeConfig[type] || {
        icon: "",
        maxLevel: 50,
        checkInterval: 300000
    };
}
\end{lstlisting}

##### 

\important{}

\begin{lstlisting}[style=javascript]
// 
function evaluateStationHealth(station) {
    const { waterLevel, temperature, pressure, lastUpdate, deviceStatus } = station;
    const now = Date.now();
    const updateAge = now - new Date(lastUpdate).getTime();
    
    // 
    if (deviceStatus !== "online") {
        return { 
            status: "offline", 
            message: "", 
            priority: "high",
            actions: ["", "", ""]
        };
    }
    
    // 
    if (updateAge > 5 } 60 \highlight{ 1000) { // 5
        return { 
            status: "stale", 
            message: } ${Math.round(updateAge / 60000)} \code{, 
            priority: "medium",
            actions: ["", ""]
        };
    }
    
    // 
    if (waterLevel > 90 || waterLevel < 10) {
        const condition = waterLevel > 90 ? "" : "";
        return { 
            status: "critical", 
            message: }${condition}: ${waterLevel}\code{, 
            priority: "high",
            actions: ["", "", ""]
        };
    }
    
    // 
    if (temperature > 35 || temperature < -10) {
        return { 
            status: "warning", 
            message: }: ${temperature}°C\code{, 
            priority: "medium",
            actions: ["", ""]
        };
    }
    
    // 
    if (pressure && (pressure < 0.8 || pressure > 1.2)) {
        return { 
            status: "warning", 
            message: }: ${pressure}bar\code{, 
            priority: "medium",
            actions: ["", ""]
        };
    }
    
    // 
    return { 
        status: "healthy", 
        message: "", 
        priority: "low",
        actions: [""]
    };
}

// 
const stationData = {
    id: "ST001",
    name: "",
    waterLevel: 95.5,
    temperature: 28,
    pressure: 1.01,
    lastUpdate: new Date(Date.now() - 2 } 60 \highlight{ 1000), // 2
    deviceStatus: "online"
};

const healthStatus = evaluateStationHealth(stationData);
console.log(":", healthStatus);
\end{lstlisting}

\important{}

\begin{lstlisting}[style=javascript]
// 
function determineWarningLevel(monitoringData) {
    const { 
        currentLevel, 
        trend, 
        rateOfChange, 
        weather, 
        season,
        historicalData 
    } = monitoringData;
    
    let warningLevel = "normal";
    let factors = [];
    let score = 0;
    
    // 
    if (currentLevel > 95) {
        score += 50;
        factors.push("");
    } else if (currentLevel > 85) {
        score += 30;
        factors.push("");
    } else if (currentLevel > 75) {
        score += 15;
        factors.push("");
    }
    
    // 
    if (trend === "rising" && rateOfChange > 2) {
        score += 25;
        factors.push("");
    } else if (trend === "rising" && rateOfChange > 0.5) {
        score += 15;
        factors.push("");
    } else if (trend === "falling" && rateOfChange < -2) {
        score += 10;
        factors.push("");
    }
    
    // 
    if (weather === "heavy_rain") {
        score += 30;
        factors.push("");
    } else if (weather === "rain") {
        score += 15;
        factors.push("");
    } else if (weather === "drought") {
        score += 10;
        factors.push("");
    }
    
    // 
    if (season === "flood_season") {
        score += 10;
        factors.push("");
    }
    
    // 
    if (historicalData) {
        const avgLevel = historicalData.averageLevel;
        const deviation = Math.abs(currentLevel - avgLevel);
        
        if (deviation > 20) {
            score += 20;
            factors.push("");
        } else if (deviation > 10) {
            score += 10;
            factors.push("");
        }
    }
    
    // 
    if (score >= 70) {
        warningLevel = "emergency";
    } else if (score >= 50) {
        warningLevel = "high";
    } else if (score >= 30) {
        warningLevel = "medium";
    } else if (score >= 15) {
        warningLevel = "low";
    }
    
    return {
        level: warningLevel,
        score,
        factors,
        recommendation: getRecommendationByLevel(warningLevel),
        timestamp: new Date().toISOString()
    };
}

// 
function getRecommendationByLevel(level) {
    const recommendations = {
        emergency: [
            "",
            "", 
            "",
            ""
        ],
        high: [
            "",
            "",
            "",
            ""
        ],
        medium: [
            "",
            "",
            "",
            ""
        ],
        low: [
            "",
            "",
            ""
        ],
        normal: [
            "",
            ""
        ]
    };
    
    return recommendations[level] || recommendations.normal;
}

// 
const monitoringData = {
    currentLevel: 88.5,
    trend: "rising",
    rateOfChange: 1.5,
    weather: "heavy_rain",
    season: "flood_season",
    historicalData: { averageLevel: 65 }
};

const warningResult = determineWarningLevel(monitoringData);
console.log(":", warningResult);
\end{lstlisting}

#### 2. 

JavaScript

##### for - 

for

\begin{lstlisting}[style=javascript]
// for
function calculateHourlyAverage(hourlyData) {
    let totalSum = 0;
    let validCount = 0;
    
    // for
    for (let i = 0; i < hourlyData.length; i++) {
        // 
        if (hourlyData[i] !== null && hourlyData[i] !== undefined && !isNaN(hourlyData[i])) {
            totalSum += hourlyData[i];
            validCount++;
        }
    }
    
    return {
        average: validCount > 0 ? totalSum / validCount : 0,
        validCount: validCount,
        totalCount: hourlyData.length,
        dataIntegrity: validCount / hourlyData.length
    };
}

// 
const hourlyLevels = [85.2, 84.8, null, 86.1, 85.5, undefined, 84.9, 85.8];
const avgResult = calculateHourlyAverage(hourlyLevels);
console.log(":", avgResult);
\end{lstlisting}

##### while - 

while

\begin{lstlisting}[style=javascript]
// 
function attemptConnection(maxAttempts = 5) {
    let attempts = 0;
    let connected = false;
    let lastError = null;
    
    while (attempts < maxAttempts && !connected) {
        attempts++;
        console.log(}${attempts}...\code{);
        
        try {
            // 
            const success = Math.random() > 0.6; // 40%
            
            if (success) {
                connected = true;
                console.log("");
            } else {
                throw new Error("");
            }
        } catch (error) {
            lastError = error.message;
            console.log(}${error.message}\code{);
            
            if (attempts < maxAttempts) {
                console.log("3...");
                // 
            }
        }
    }
    
    return {
        success: connected,
        attempts: attempts,
        error: connected ? null : lastError
    };
}

// 
function collectValidData(validator) {
    let data;
    let isValid = false;
    let attempts = 0;
    const maxAttempts = 10;
    
    while (!isValid && attempts < maxAttempts) {
        attempts++;
        
        // 
        data = {
            waterLevel: Math.random() } 120 - 10, // -10  110 
            temperature: Math.random() \highlight{ 60 - 20, // -20  40 
            timestamp: Date.now()
        };
        
        // 
        isValid = validator(data);
        
        if (!isValid) {
            console.log(}${attempts}...\code{);
        }
    }
    
    if (isValid) {
        console.log(}${attempts}\code{);
        return { success: true, data, attempts };
    } else {
        console.log(}\code{);
        return { success: false, data: null, attempts };
    }
}
\end{lstlisting}

##### do...while - 

do...while

\begin{lstlisting}[style=javascript]
// 
function getUserInput(promptMessage, validator) {
    let userInput;
    let isValid = false;
    
    do {
        // 
        userInput = prompt(promptMessage);
        
        if (userInput === null) {
            // 
            return { success: false, value: null, message: "" };
        }
        
        // 
        const validationResult = validator(userInput);
        isValid = validationResult.isValid;
        
        if (!isValid) {
            alert(}: ${validationResult.error}\code{);
        }
        
    } while (!isValid);
    
    return { success: true, value: userInput, message: "" };
}

// 
function collectMinimumDataSet() {
    let dataSet = [];
    let collectionRound = 0;
    
    do {
        collectionRound++;
        console.log(}${collectionRound}...\code{);
        
        // 
        const newData = {
            round: collectionRound,
            timestamp: Date.now(),
            waterLevel: Math.random() } 100,
            flow: Math.random() \highlight{ 50,
            temperature: Math.random() } 40
        };
        
        dataSet.push(newData);
        console.log(}${collectionRound}\code{);
        
        // 3
    } while (dataSet.length < 3 || !isDataSetComplete(dataSet));
    
    return {
        dataSet,
        rounds: collectionRound,
        totalPoints: dataSet.length
    };
}

// 
function isDataSetComplete(dataSet) {
    // 
    return dataSet.every(data => 
        data.waterLevel !== undefined && 
        data.flow !== undefined && 
        data.temperature !== undefined
    );
}
\end{lstlisting}

##### for...in - 

for...in

\begin{lstlisting}[style=javascript]
// 
function displayStationInfo(station) {
    console.log("===  ===");
    
    // 
    for (let property in station) {
        // 
        if (station.hasOwnProperty(property)) {
            const value = station[property];
            const formattedValue = formatPropertyValue(property, value);
            console.log(}${getPropertyDisplayName(property)}: ${formattedValue}\code{);
        }
    }
    
    console.log("===  ===");
}

// 
function formatPropertyValue(property, value) {
    switch (property) {
        case 'waterLevel':
            return }${value}\code{;
        case 'temperature':
            return }${value}°C\code{;
        case 'timestamp':
            return new Date(value).toLocaleString();
        case 'isOnline':
            return value ? '' : '';
        case 'location':
            return }: ${value.lat}, : ${value.lng}\code{;
        default:
            return value;
    }
}

// 
function getPropertyDisplayName(property) {
    const displayNames = {
        id: 'ID',
        name: '', 
        waterLevel: '',
        temperature: '',
        isOnline: '',
        location: '',
        timestamp: ''
    };
    
    return displayNames[property] || property;
}

// 
const station = {
    id: 'ST001',
    name: '',
    waterLevel: 85.5,
    temperature: 25.2,
    isOnline: true,
    location: { lat: 30.5, lng: 114.3 },
    timestamp: Date.now()
};

displayStationInfo(station);
\end{lstlisting}

##### for...of - 

for...ofSetMap

\begin{lstlisting}[style=javascript]
// 
function processStationList(stations) {
    const results = [];
    const summary = {
        total: 0,
        online: 0,
        warning: 0,
        normal: 0
    };
    
    // for...of
    for (let station of stations) {
        // 
        const processed = {
            id: station.id,
            name: station.name,
            waterLevel: station.waterLevel,
            status: determineStationStatus(station),
            lastCheck: new Date().toISOString()
        };
        
        results.push(processed);
        
        // 
        summary.total++;
        if (station.isOnline) summary.online++;
        if (processed.status === 'warning') summary.warning++;
        if (processed.status === 'normal') summary.normal++;
    }
    
    return { results, summary };
}

// 
function determineStationStatus(station) {
    if (!station.isOnline) return 'offline';
    if (station.waterLevel > 80) return 'warning';
    if (station.waterLevel < 20) return 'low';
    return 'normal';
}

// forEach - 
function displayStationData(stations) {
    console.log("===  ===");
    
    // forEach
    stations.forEach((station, index) => {
        const statusIcon = getStatusIcon(station.isOnline, station.waterLevel);
        console.log(}${index + 1}. ${statusIcon} ${station.name}: ${station.waterLevel}m\code{);
    });
}

// 
function getStatusIcon(isOnline, waterLevel) {
    if (!isOnline) return '';
    if (waterLevel > 80) return '🟡';
    if (waterLevel < 20) return '';
    return '🟢';
}

// 
const stations = [
    { id: 'ST001', name: '', waterLevel: 85.5, isOnline: true },
    { id: 'ST002', name: '', waterLevel: 65.2, isOnline: true },
    { id: 'ST003', name: '', waterLevel: 15.8, isOnline: false }
];

const processResult = processStationList(stations);
console.log(":", processResult);

displayStationData(stations);
\end{lstlisting}

##### 

JavaScriptbreakcontinue

\important{break - }

break

\begin{lstlisting}[style=javascript]
// 
function findFirstCriticalStation(stations) {
    let criticalStation = null;
    
    for (let i = 0; i < stations.length; i++) {
        const station = stations[i];
        
        // 
        if (station.status !== "online") {
            console.log(}${station.name}\code{);
            continue; // 
        }
        
        // 
        if (station.waterLevel > 90) {
            console.log(}${station.name}${station.waterLevel}\code{);
            criticalStation = station;
            break; // 
        }
        
        console.log(}${station.name}${station.waterLevel}\code{);
    }
    
    return {
        found: criticalStation !== null,
        station: criticalStation,
        message: criticalStation ? 
            }${criticalStation.name}\code{ : 
            ""
    };
}

// 
function validateDataQuality(dataPoints) {
    const issues = [];
    let fatalError = false;
    
    for (let i = 0; i < dataPoints.length; i++) {
        const point = dataPoints[i];
        
        // 
        if (!point || typeof point !== 'object') {
            issues.push(}${i + 1}: \code{);
            fatalError = true;
            break; // 
        }
        
        // 
        if (!point.hasOwnProperty('timestamp')) {
            issues.push(}${i + 1}: \code{);
            continue; // 
        }
        
        if (!point.hasOwnProperty('waterLevel')) {
            issues.push(}${i + 1}: \code{);
            continue;
        }
        
        // 
        if (point.waterLevel < -100 || point.waterLevel > 500) {
            issues.push(}${i + 1}:  (${point.waterLevel})\code{);
            fatalError = true;
            break; // 
        }
    }
    
    return {
        isValid: !fatalError && issues.length === 0,
        issues,
        fatalError,
        checkedCount: fatalError ? 
            dataPoints.findIndex(p => !p || typeof p !== 'object' || 
                                (p.waterLevel && (p.waterLevel < -100 || p.waterLevel > 500))) + 1 :
            dataPoints.length
    };
}
\end{lstlisting}

\important{continue - }

continue

\begin{lstlisting}[style=javascript]
// 
function analyzeValidStations(stations) {
    const analysis = {
        total: stations.length,
        processed: 0,
        skipped: 0,
        online: 0,
        offline: 0,
        warning: 0,
        normal: 0,
        skippedReasons: []
    };
    
    for (let station of stations) {
        // 
        if (!station || !station.id) {
            analysis.skipped++;
            analysis.skippedReasons.push("");
            continue;
        }
        
        // 
        if (station.id.startsWith('TEST')) {
            analysis.skipped++;
            analysis.skippedReasons.push(}: ${station.id}\code{);
            continue;
        }
        
        // 
        if (station.maintenance === true) {
            analysis.skipped++;
            analysis.skippedReasons.push(}: ${station.name}\code{);
            continue;
        }
        
        // 
        analysis.processed++;
        
        // 
        if (station.status === 'online') {
            analysis.online++;
            
            // 
            if (station.waterLevel > 80) {
                analysis.warning++;
            } else {
                analysis.normal++;
            }
        } else {
            analysis.offline++;
        }
    }
    
    return analysis;
}

// 
const stationList = [
    { id: 'ST001', name: '', status: 'online', waterLevel: 85.2 },
    { id: 'TEST001', name: '', status: 'online', waterLevel: 50 },
    { id: 'ST002', name: '', status: 'offline', waterLevel: null },
    { id: 'ST003', name: '', maintenance: true, status: 'online' },
    null, // 
    { id: 'ST004', name: '', status: 'online', waterLevel: 65.8 }
];

const analysis = analyzeValidStations(stationList);
console.log(":", analysis);
\end{lstlisting}

\important{label- }

\begin{lstlisting}[style=javascript]
// 
function findEmergencyStationInRegions(regions) {
    const searchResults = {
        found: false,
        region: null,
        station: null,
        searchPath: []
    };
    
    // 
    regionLoop: for (let region of regions) {
        console.log(}${region.name}\code{);
        searchResults.searchPath.push(}: ${region.name}\code{);
        
        // 
        if (region.priority === 'low') {
            console.log(}${region.name}\code{);
            searchResults.searchPath.push(}: ${region.name}\code{);
            continue regionLoop;
        }
        
        // 
        for (let station of region.stations) {
            searchResults.searchPath.push(}: ${station.name}\code{);
            
            // 
            if (station.status !== "online") {
                console.log(}${station.name}\code{);
                continue; // 
            }
            
            // 
            if (station.waterLevel > 95) {
                console.log(}${region.name} - ${station.name}\code{);
                searchResults.found = true;
                searchResults.region = region;
                searchResults.station = station;
                searchResults.searchPath.push(}: ${station.name}\code{);
                
                // 
                break regionLoop;
            }
            
            console.log(}${station.name}...\code{);
        }
        
        console.log(}${region.name}\code{);
        searchResults.searchPath.push(}: ${region.name}\code{);
    }
    
    if (searchResults.found) {
        console.log("");
    } else {
        console.log("");
        searchResults.searchPath.push("");
    }
    
    return searchResults;
}

// 
function validateNestedData(dataStructure) {
    const validationResults = {
        isValid: true,
        errors: [],
        validatedItems: 0
    };
    
    // 
    outerValidation: for (let categoryName in dataStructure) {
        const category = dataStructure[categoryName];
        
        if (!Array.isArray(category)) {
            validationResults.errors.push(} ${categoryName} \code{);
            validationResults.isValid = false;
            break outerValidation; // 
        }
        
        // 
        for (let i = 0; i < category.length; i++) {
            const item = category[i];
            validationResults.validatedItems++;
            
            // 
            if (!item || typeof item !== 'object') {
                validationResults.errors.push(
                    }${categoryName}[${i}]: \code{
                );
                continue; // 
            }
            
            // 
            if (!item.id || !item.timestamp) {
                validationResults.errors.push(
                    }${categoryName}[${i}]: \code{
                );
                validationResults.isValid = false;
                
                // 
                if (validationResults.errors.length > 10) {
                    console.log("");
                    break outerValidation;
                }
            }
        }
    }
    
    return validationResults;
}

// 
const regions = [
    {
        name: "", 
        priority: "high",
        stations: [
            { name: "", status: "online", waterLevel: 88.5 },
            { name: "", status: "online", waterLevel: 96.2 }
        ]
    },
    {
        name: "", 
        priority: "medium",
        stations: [
            { name: "", status: "offline", waterLevel: 75.0 }
        ]
    }
];

const searchResult = findEmergencyStationInRegions(regions);
console.log(":", searchResult);
\end{lstlisting}

#### 3. 

##### try...catch...finally

JavaScript

\begin{lstlisting}[style=javascript]
// JSON
function safeDataParse(jsonString, defaultValue = null) {
    let parseResult = {
        success: false,
        data: null,
        error: null
    };
    
    try {
        // JSON
        const data = JSON.parse(jsonString);
        console.log("");
        
        parseResult.success = true;
        parseResult.data = data;
        
        return parseResult;
        
    } catch (error) {
        // 
        console.error(":", error.message);
        
        parseResult.error = {
            type: 'ParseError',
            message: error.message,
            timestamp: new Date()
        };
        
        // 
        parseResult.data = defaultValue;
        
        return parseResult;
        
    } finally {
        // 
        console.log("");
        
        // 
        logOperation('data_parse', {
            timestamp: new Date(),
            success: parseResult.success,
            inputLength: jsonString ? jsonString.length : 0
        });
    }
}

// 
function fetchMonitoringData(url, options = {}) {
    const requestInfo = {
        url,
        startTime: Date.now(),
        retryCount: 0
    };
    
    try {
        // URL
        if (!url || typeof url !== 'string') {
            throw new Error('URL');
        }
        
        // 
        const timeout = options.timeout || 10000;
        const controller = new AbortController();
        
        // 
        const timeoutId = setTimeout(() => {
            controller.abort();
        }, timeout);
        
        console.log(}: ${url}\code{);
        
        // fetch
        const response = simulateFetch(url, { 
            ...options, 
            signal: controller.signal 
        });
        
        clearTimeout(timeoutId);
        
        return {
            success: true,
            data: response,
            requestInfo
        };
        
    } catch (error) {
        console.error(} (${url}):\code{, error.message);
        
        return {
            success: false,
            error: {
                type: error.name || 'RequestError',
                message: error.message,
                url,
                timestamp: new Date()
            },
            requestInfo
        };
        
    } finally {
        requestInfo.endTime = Date.now();
        requestInfo.duration = requestInfo.endTime - requestInfo.startTime;
        
        console.log(}: ${requestInfo.duration}ms\code{);
    }
}

// fetch
function simulateFetch(url, options) {
    // 
    if (Math.random() < 0.2) {
        throw new Error('');
    }
    
    return {
        status: 'success',
        data: {
            stationId: 'ST001',
            waterLevel: 85.5,
            timestamp: new Date()
        }
    };
}

// 
function logOperation(operation, details) {
    const logEntry = {
        operation,
        timestamp: new Date().toISOString(),
        ...details
    };
    
    // 
    console.log(':', logEntry);
}
\end{lstlisting}

##### 

\begin{lstlisting}[style=javascript]
// 
class WaterLevelError extends Error {
    constructor(message, level, stationId) {
        super(message);
        this.name = 'WaterLevelError';
        this.level = level;
        this.stationId = stationId;
        this.timestamp = new Date();
    }
}

class DataValidationError extends Error {
    constructor(message, field, value) {
        super(message);
        this.name = 'DataValidationError';
        this.field = field;
        this.value = value;
        this.timestamp = new Date();
    }
}

class DeviceConnectionError extends Error {
    constructor(message, deviceId, lastContact) {
        super(message);
        this.name = 'DeviceConnectionError';
        this.deviceId = deviceId;
        this.lastContact = lastContact;
        this.timestamp = new Date();
    }
}

// 
function validateWaterLevel(level, stationId) {
    // 
    if (typeof level !== 'number') {
        throw new DataValidationError(
            "", 
            "waterLevel", 
            level
        );
    }
    
    // 
    if (isNaN(level)) {
        throw new DataValidationError(
            "NaN",
            "waterLevel",
            level
        );
    }
    
    // 
    if (level < -100) {
        throw new WaterLevelError(
            "",
            level,
            stationId
        );
    }
    
    if (level > 300) {
        throw new WaterLevelError(
            "",
            level,
            stationId
        );
    }
    
    // 
    if (level > 100) {
        throw new WaterLevelError(
            "",
            level,
            stationId
        );
    }
    
    return true;
}

// 
function checkDeviceConnection(device) {
    const now = Date.now();
    const lastContact = new Date(device.lastContact).getTime();
    const timeSinceContact = now - lastContact;
    
    // 
    if (timeSinceContact > 30 \highlight{ 60 } 1000) { // 30
        throw new DeviceConnectionError(
            }30\code{,
            device.id,
            device.lastContact
        );
    }
    
    // 
    if (device.status !== 'online') {
        throw new DeviceConnectionError(
            }: ${device.status}\code{,
            device.id,
            device.lastContact
        );
    }
    
    return true;
}

##### 

\end{lstlisting}javascript
// 
function processMonitoringData(rawData, options = {}) {
    const processingLog = {
        startTime: Date.now(),
        steps: [],
        errors: [],
        warnings: []
    };
    
    try {
        // 1
        processingLog.steps.push('');
        if (!rawData || typeof rawData !== 'object') {
            throw new DataValidationError("", "rawData", rawData);
        }
        
        // 2
        processingLog.steps.push('');
        const stationId = String(rawData.stationId || '').trim();
        if (!stationId) {
            throw new DataValidationError("ID", "stationId", rawData.stationId);
        }
        
        // 3
        processingLog.steps.push('');
        let waterLevel;
        try {
            waterLevel = parseFloat(rawData.waterLevel);
            validateWaterLevel(waterLevel, stationId);
        } catch (error) {
            if (error instanceof WaterLevelError) {
                // 
                processingLog.warnings.push({
                    type: 'water_level_anomaly',
                    message: error.message,
                    level: error.level,
                    stationId: error.stationId
                });
                waterLevel = error.level;
            } else {
                throw error;
            }
        }
        
        // 4
        if (rawData.deviceInfo) {
            try {
                processingLog.steps.push('');
                checkDeviceConnection(rawData.deviceInfo);
            } catch (error) {
                if (error instanceof DeviceConnectionError) {
                    processingLog.warnings.push({
                        type: 'device_connection',
                        message: error.message,
                        deviceId: error.deviceId
                    });
                }
            }
        }
        
        // 
        const processedData = {
            stationId,
            waterLevel,
            timestamp: new Date(),
            status: waterLevel > 80 ? "warning" : "normal",
            hasWarnings: processingLog.warnings.length > 0
        };
        
        return {
            success: true,
            data: processedData,
            processingLog
        };
        
    } catch (error) {
        processingLog.errors.push({
            type: error.name,
            message: error.message,
            timestamp: new Date()
        });
        
        console.error(":", error.message);
        
        return {
            success: false,
            error: error.message,
            errorType: error.name,
            data: null,
            processingLog
        };
    } finally {
        processingLog.endTime = Date.now();
        processingLog.duration = processingLog.endTime - processingLog.startTime;
        console.log(}: ${processingLog.duration}ms\code{);
    }
}

// 
const testData = {
    stationId: 'ST001',
    waterLevel: '95.5',
    deviceInfo: {
        id: 'DEV001',
        status: 'online',
        lastContact: new Date(Date.now() - 10 \highlight{ 60 } 1000) // 10
    }
};

const result = processMonitoringData(testData);
console.log(':', result);
\begin{lstlisting}
}`\code{

### 

JavaScriptJavaScript

#### 1. 

\end{lstlisting}javascript
// 
function calculateFlow(area, velocity) {
    return area \highlight{ velocity;
}

// 
const calculatePressure = function(force, area) {
    return force / area;
};

//  (IIFE)
const moduleResult = (function() {
    const privateVar = "";
    
    return {
        getPrivateVar: function() {
            return privateVar;
        }
    };
})();

// ES6+
const convertTemperature = (celsius) => celsius } 9/5 + 32;
const getStatusIcon = (status) => status === "online" ? "🟢" : "";

// 
const processStationData = (station) => {
    const level = station.waterLevel;
    const status = level > 80 ? "warning" : "normal";
    
    return {
        id: station.id,
        level,
        status,
        processed: true
    };
};
\begin{lstlisting}
#### 2. 

\end{lstlisting}javascript
// ES6+
function createAlert(message, type = "info", duration = 3000) {
    return {
        message,
        type,
        duration,
        timestamp: Date.now()
    };
}

// ...rest
function calculateAverage(...values) {
    if (values.length === 0) return 0;
    const sum = values.reduce((total, value) => total + value, 0);
    return sum / values.length;
}

// 
function createStation({id, name, location = {}, waterLevel = 0}) {
    return {
        id,
        name,
        latitude: location.lat || 0,
        longitude: location.lng || 0,
        waterLevel,
        status: "active"
    };
}

// 
function validateAndProcess(data) {
    // 
    if (arguments.length === 0) {
        throw new Error("");
    }
    
    // 
    if (typeof data !== 'object' || data === null) {
        throw new TypeError("");
    }
    
    // ...
    return processData(data);
}
\begin{lstlisting}
#### 3. 

\end{lstlisting}javascript
// 
const globalVar = "";

function outerFunction(outerParam) {
    const outerVar = "";
    
    function innerFunction(innerParam) {
        const innerVar = "";
        
        // 
        console.log(globalVar);  // 
        console.log(outerParam); // 
        console.log(outerVar);   // 
        console.log(innerParam); // 
        console.log(innerVar);   // 
    }
    
    return innerFunction;
}

// 
function createCounter(initialValue = 0) {
    let count = initialValue;
    
    return {
        increment: function() {
            count++;
            return count;
        },
        decrement: function() {
            count--;
            return count;
        },
        getValue: function() {
            return count;
        },
        reset: function() {
            count = initialValue;
            return count;
        }
    };
}

// 
const stationCounter = createCounter(0);
console.log(stationCounter.increment()); // 1
console.log(stationCounter.increment()); // 2
console.log(stationCounter.getValue());  // 2

// 
const WaterMonitoringModule = (function() {
    // 
    let stations = [];
    let alertThreshold = 80;
    
    function validateLevel(level) {
        return typeof level === 'number' && level >= 0 && level <= 200;
    }
    
    // 
    return {
        addStation: function(station) {
            if (validateLevel(station.waterLevel)) {
                stations.push(station);
                return true;
            }
            return false;
        },
        
        getAlerts: function() {
            return stations.filter(s => s.waterLevel > alertThreshold);
        },
        
        setThreshold: function(newThreshold) {
            if (validateLevel(newThreshold)) {
                alertThreshold = newThreshold;
            }
        }
    };
})();
\begin{lstlisting}
JavaScript

###  - 

JavaScript\important{}\important{}

#### 1.  - 

JavaScript

\important{}

JavaScriptpropertiesmethods

- \important{}
- \important{}
- \important{}
- \important{}

\important{}

 }{}\textbackslash code\{

\textbackslash end\{lstlisting}javascript // const waterStation = \{
id: ``WS001'', name: ``\,``, waterLevel: 15.2, isOnline: true };

\begin{lstlisting}
\important{}

\end{lstlisting}

javascript // const monitoringStation = \{ id: ``WS001'', name: ``\,``,

\begin{verbatim}
// 
location: {
    latitude: 30.7128,
    longitude: 111.3200,
    altitude: 45.5,
    description: ""
},

// 
currentData: {
    waterLevel: 15.2,
    temperature: 18.3,
    timestamp: new Date()
}
\end{verbatim}

};

\begin{lstlisting}
\important{}

\end{lstlisting}

javascript const waterMonitor = \{ currentLevel: 15.2, alertThreshold:
20.0,

\begin{verbatim}
// 
updateLevel: function(newLevel) {
    this.currentLevel = newLevel;
    console.log(}: ${newLevel}\code{);
},

// ES6
checkAlert() {
    return this.currentLevel > this.alertThreshold;
},

// 
getStatus() {
    if (this.checkAlert()) {
        return "";
    }
    return "";
}
\end{verbatim}

};

// waterMonitor.updateLevel(18.5);
console.log(waterMonitor.getStatus());

\begin{lstlisting}
#### 2. 

\important{}

JavaScript

\end{lstlisting}

javascript // const sensorData = \{ temperature: 18.5, pressure: 1013.25
};

// sensorData.humidity = 65; // sensorData{[}``windSpeed''{]} = 12.3; //

// if (``temperature'' in sensorData) \{ console.log(``\,``); }

\begin{lstlisting}
\important{}

JavaScript

\end{lstlisting}

javascript const stationData = \{ id: ``WS001'', waterLevel: 15.2,
temperature: 18.5, lastUpdate: new Date() };

// const propertyNames = Object.keys(stationData); console.log(``:'',
propertyNames);

//\\
const propertyValues = Object.values(stationData); console.log(``:'',
propertyValues);

// const entries = Object.entries(stationData); entries.forEach(({[}key,
value{]}) =\textgreater{} \{ console.log(}\$\{key}:
\$\{value}\code{);
});

\begin{lstlisting}
\important{}

\end{lstlisting}

javascript // const stationInfo = \{ id: ``WS001'', name: ``\,'' };

// const locationInfo = \{ latitude: 30.5928, longitude: 114.3055 };

// Object.assign const completeInfo = Object.assign(\{}, stationInfo,
locationInfo);

// ES6 const modernMerge = \{ \ldots stationInfo, \ldots locationInfo,
status: ``active'' // };

\begin{lstlisting}
\important{}

\end{lstlisting}

javascript const originalReading = \{ stationId: ``WS001'', data: \{
waterLevel: 15.2, temperature: 18.3 } };

// - const shallowCopy = \{ \ldots originalReading };
shallowCopy.data.waterLevel = 20.5; //

// - const deepCopy = JSON.parse(JSON.stringify(originalReading));
deepCopy.data.temperature = 25.0; //

\begin{lstlisting}
#### 3.  - 

JavaScript

\important{}

- -   
- - 

- \important{}0
- \important{}
- \important{}
- \important{}

\important{}

\end{lstlisting}

javascript // const waterLevels = {[}12.5, 13.2, 14.1, 15.8{]};

// const emptyReadings = new Array(); const fixedLength = new Array(7);
// 7

// const stationInfo = {[}``WS001'', ``\,``, 15.2, true, new Date(){]};

\begin{lstlisting}
\important{}

0length

\end{lstlisting}

javascript const dailyReadings = {[}15.2, 16.1, 17.3, 18.5{]};

// console.log(``:'', dailyReadings{[}0{]}); console.log(``:'',
dailyReadings{[}dailyReadings.length - 1{]});

// console.log(``:'', dailyReadings.length);

// dailyReadings{[}1{]} = 16.8; //

\begin{lstlisting}
\important{}

\end{lstlisting}

javascript const stationList = {[}``WS001'', ``WS002'', ``WS003''{]};

// stationList.push(``WS004'');

//\\
stationList.unshift(``WS000'');

// const removedStation = stationList.pop();

// const firstStation = stationList.shift();

console.log(``:'', stationList);

\begin{lstlisting}
\important{}

\end{lstlisting}

javascript const monitoringStations = {[} \{ id: ``WS001'', name: ````,
waterLevel: 15.2, status:''normal'' }, \{ id: ``WS002'', name: ````,
waterLevel: 12.8, status:''low'' }{]};

// console.log(``:'', monitoringStations{[}0{]}.name);
console.log(``:'', monitoringStations{[}1{]}.waterLevel);

\begin{lstlisting}
#### 4.  - 

JavaScript

\important{forEach - }

forEach

\end{lstlisting}

javascript const dailyReadings = {[}15.2, 16.1, 17.3, 18.5{]};

// dailyReadings.forEach((reading, index) =\textgreater{} \{
console.log(}\$\{index + 1}: \$\{reading}\code{);
    
    // 
    if (reading > 18.0) {
        console.log(}${reading}\code{);
    }
});

\begin{lstlisting}
\important{map - }

map

\end{lstlisting}

javascript const temperatures = {[}18.3, 19.1, 17.8{]};

// const fahrenheitTemps = temperatures.map(celsius =\textgreater{} (\{
celsius: celsius, fahrenheit: (celsius \highlight{ 9/5) + 32
}));

// const stationData = {[} \{ id: ``WS001'', waterLevel: 15.2 }, \{ id:
``WS002'', waterLevel: 18.7 }{]}; const waterLevels =
stationData.map(station =\textgreater{} station.waterLevel);

\begin{lstlisting}
\important{filter - }

filter

\end{lstlisting}

javascript const monitoringData = {[} \{ station: ``WS001'', level:
15.2, alert: false }, \{ station: ``WS002'', level: 22.1, alert: true
}, \{ station: ``WS003'', level: 18.7, alert: false }{]};

// const alertStations = monitoringData.filter(data =\textgreater{}
data.alert);

// 20 const highLevelStations = monitoringData.filter(data
=\textgreater{} data.level \textgreater{} 20);

\begin{lstlisting}
\important{findsome - }

findsome

\end{lstlisting}

javascript const stationList = {[} \{ id: ``WS001'', status: ``online''
}, \{ id: ``WS002'', status: ``offline'' }, \{ id: ``WS003'', status:
``online'' }{]};

// const offlineStation = stationList.find(station =\textgreater{}
station.status === ``offline'');

// const hasOfflineStation = stationList.some(station =\textgreater{}
station.status === ``offline'');

console.log(``:'', offlineStation?.id); console.log(``:'',
hasOfflineStation);

\begin{lstlisting}
\important{reduce - }

reducereducer

\end{lstlisting}

javascript const readings = {[}15.2, 16.1, 17.3, 18.5, 16.8{]};

// const averageLevel = readings.reduce((sum, reading, index, array)
=\textgreater{} \{ sum += reading; return index === array.length - 1 ?
sum / array.length : sum; }, 0);

// const maxLevel = readings.reduce((max, current) =\textgreater{}
current \textgreater{} max ? current : max, readings{[}0{]});

console.log(}: \$\{averageLevel.toFixed(2)}\code{);
console.log(}: \$\{maxLevel}\code{);
\begin{lstlisting}
#### 5. 

JavaScript

\important{}

- \important{}
- \important{}
- \important{}
- \important{}

\important{}

\end{lstlisting}javascript
// 
const waterMonitoringSystem = {
    systemName: "",
    stations: [
        {
            id: "WS001",
            name: "",
            location: { lat: 30.7128, lng: 111.3200 },
            readings: [
                { timestamp: new Date(), waterLevel: 15.2, temperature: 18.3 },
                { timestamp: new Date(), waterLevel: 15.8, temperature: 18.7 }
            ]
        },
        {
            id: "WS002", 
            name: "",
            location: { lat: 32.0042, lng: 112.1225 },
            readings: [
                { timestamp: new Date(), waterLevel: 12.1, temperature: 17.9 }
            ]
        }
    ],
    
    // 
    getCurrentLevels() {
        return this.stations.map(station => ({
            stationName: station.name,
            currentLevel: station.readings[station.readings.length - 1]?.waterLevel || 0
        }));
    },
    
    // 
    findHighWaterStations(threshold = 20) {
        return this.stations.filter(station => {
            const latestReading = station.readings[station.readings.length - 1];
            return latestReading && latestReading.waterLevel > threshold;
        });
    }
};

// 
const currentLevels = waterMonitoringSystem.getCurrentLevels();
console.log(":", currentLevels);

const alertStations = waterMonitoringSystem.findHighWaterStations(15);
console.log(":", alertStations);
\begin{lstlisting}
\important{}

\end{lstlisting}javascript
// 
const processingFunctions = {
    // 
    calculateAverageLevel(stationReadings) {
        const levels = stationReadings.map(reading => reading.waterLevel);
        return levels.reduce((sum, level) => sum + level, 0) / levels.length;
    },
    
    // 
    generateStationReport(station) {
        const readings = station.readings;
        if (readings.length === 0) return null;
        
        return {
            stationId: station.id,
            stationName: station.name,
            totalReadings: readings.length,
            averageLevel: this.calculateAverageLevel(readings),
            latestReading: readings[readings.length - 1]
        };
    }
};

// 
const stationReports = waterMonitoringSystem.stations
    .map(station => processingFunctions.generateStationReport(station))
    .filter(report => report !== null);

console.log(":", stationReports);
\begin{lstlisting}

###  - 

7×24\important{}

#### 1. JavaScript

JavaScript

\end{lstlisting}javascript
// 

// 1.  (SyntaxError) - 
// 
/} // function processWaterData() \{ console.log(``\,``); //
\highlight{/

// 2.  (Runtime Error) - 
function calculateFlowRate(volume, time) {
    // ReferenceError - 
    try {
        if (time === 0) {
            // InfinityNaN
            throw new Error("");
        }
        
        // TypeError - 
        const result = volume / time;
        return result;
        
    } catch (error) {
        console.error(":", error.message);
        return null;
    }
}

// 3.  - 
function validateWaterLevel(level, stationId) {
    // 
    if (level > 0 && level < 100) {  // 
        return true;
    }
    
    // 
    console.log(}${stationId}${level}\code{);
    return false;
}

// JavaScript
function demonstrateErrorTypes() {
    try {
        // SyntaxError - eval
        eval("function invalid syntax");
    } catch (e) {
        console.log(":", e instanceof SyntaxError);
    }
    
    try {
        // ReferenceError - 
        console.log(undefinedVariable);
    } catch (e) {
        console.log(":", e instanceof ReferenceError);
    }
    
    try {
        // TypeError - 
        const nullValue = null;
        nullValue.someMethod();
    } catch (e) {
        console.log(":", e instanceof TypeError);
    }
    
    try {
        // RangeError - 
        const arr = new Array(-1);
    } catch (e) {
        console.log(":", e instanceof RangeError);
    }
}
\begin{lstlisting}
#### 2. try-catch-finally

try-catch-finallyJavaScripttrycatchfinally

\important{try-catch}

try-catch

\end{lstlisting}javascript
// try-catch
function processWaterLevel(rawData) {
    try {
        // 
        const level = parseFloat(rawData);
        
        if (isNaN(level)) {
            throw new Error("");
        }
        
        if (level < 0 || level > 50) {
            throw new Error("");
        }
        
        console.log(}: ${level}\code{);
        return level;
        
    } catch (error) {
        // 
        console.error(":", error.message);
        return null;
    }
}

// 
processWaterLevel("15.2");    // 
processWaterLevel("invalid"); // 
processWaterLevel("-5");      // 
\begin{lstlisting}
\important{finally}

finally

\end{lstlisting}javascript
function readSensorData(sensorId) {
    let connection = null;
    
    try {
        // 
        connection = connectToSensor(sensorId);
        
        // 
        const data = connection.readData();
        
        if (!data) {
            throw new Error("");
        }
        
        return {
            success: true,
            data: data,
            timestamp: new Date()
        };
        
    } catch (error) {
        console.error(}${sensorId}:\code{, error.message);
        
        return {
            success: false,
            error: error.message,
            timestamp: new Date()
        };
        
    } finally {
        // 
        if (connection) {
            connection.close();
            console.log(}${sensorId}\code{);
        }
    }
}

// 
function connectToSensor(sensorId) {
    return {
        readData: () => ({ temperature: 18.5, humidity: 65 }),
        close: () => console.log("")
    };
}
\begin{lstlisting}
\important{}

\end{lstlisting}javascript
// 
class WaterSystemError extends Error {
    constructor(message, errorCode) {
        super(message);
        this.name = 'WaterSystemError';
        this.errorCode = errorCode;
    }
}

class SensorError extends WaterSystemError {
    constructor(message, sensorId) {
        super(message, 'SENSOR_ERROR');
        this.sensorId = sensorId;
    }
}

class DataValidationError extends WaterSystemError {
    constructor(message, fieldName, value) {
        super(message, 'DATA_VALIDATION_ERROR');
        this.fieldName = fieldName;
        this.invalidValue = value;
    }
}

// 
function validateAndProcessData(reading) {
    try {
        if (!reading.stationId) {
            throw new DataValidationError("ID", "stationId", reading.stationId);
        }
        
        if (reading.waterLevel < 0) {
            throw new DataValidationError("", "waterLevel", reading.waterLevel);
        }
        
        console.log("");
        return true;
        
    } catch (error) {
        if (error instanceof DataValidationError) {
            console.error(}: ${error.message}\code{);
            console.error(}: ${error.fieldName}, : ${error.invalidValue}\code{);
        } else {
            console.error(":", error.message);
        }
        return false;
    }
}
\begin{lstlisting}
#### 3. Promise

JavaScript

\important{Promise}

Promise}.catch()\code{async/await

\end{lstlisting}javascript
// async/await
async function fetchWaterLevelData(stationId) {
    try {
        const response = await fetch(}/api/stations/${stationId}/current\code{);
        
        // HTTP
        if (!response.ok) {
            throw new Error(}HTTP: ${response.status} ${response.statusText}\code{);
        }
        
        const data = await response.json();
        console.log(}${stationId}:\code{, data);
        return data;
        
    } catch (error) {
        if (error instanceof TypeError) {
            console.error(':', error.message);
        } else if (error instanceof SyntaxError) {
            console.error(':', error.message);
        } else {
            console.error(':', error.message);
        }
        return null;
    }
}

// 
fetchWaterLevelData("WS001").then(data => {
    if (data) {
        console.log(":", data);
    }
});
\begin{lstlisting}
\important{Promise}

Promise}.catch()\code{

\end{lstlisting}javascript
// Promise
function processWaterStationData(stationId) {
    return fetchWaterLevelData(stationId)
        .then(data => {
            if (!data) {
                throw new Error("");
            }
            return { ...data, processed: true };
        })
        .then(processedData => {
            console.log(":", processedData);
            return processedData;
        })
        .catch(error => {
            console.error(}${stationId}:\code{, error.message);
            return { error: error.message, stationId };
        })
        .finally(() => {
            console.log(}${stationId}\code{);
        });
}
\begin{lstlisting}
\important{}

}Promise.allSettled()\code{Promise

\end{lstlisting}javascript
async function fetchMultipleStationsData(stationIds) {
    // Promise.allSettled
    const results = await Promise.allSettled(
        stationIds.map(stationId => fetchWaterLevelData(stationId))
    );
    
    // 
    const successful = [];
    const failed = [];
    
    results.forEach((result, index) => {
        if (result.status === 'fulfilled' && result.value) {
            successful.push({
                stationId: stationIds[index],
                data: result.value
            });
        } else {
            failed.push({
                stationId: stationIds[index],
                error: result.reason?.message || ''
            });
        }
    });
    
    console.log(}${successful.length}\code{);
    console.log(}${failed.length}\code{);
    
    return { successful, failed };
}

// 
const stationIds = ["WS001", "WS002", "WS003"];
fetchMultipleStationsData(stationIds).then(result => {
    console.log(":", result);
});
\begin{lstlisting}
\important{WebSocket}

\end{lstlisting}javascript
function createRealtimeConnection(stationId) {
    return new Promise((resolve, reject) => {
        const ws = new WebSocket(}ws://api.example.com/realtime/${stationId}\code{);
        
        // 
        const timeout = setTimeout(() => {
            ws.close();
            reject(new Error(''));
        }, 10000);
        
        ws.onopen = () => {
            clearTimeout(timeout);
            console.log(}${stationId}\code{);
            resolve(ws);
        };
        
        ws.onmessage = (event) => {
            try {
                const data = JSON.parse(event.data);
                console.log(}${stationId}:\code{, data);
            } catch (error) {
                console.error(':', error.message);
            }
        };
        
        ws.onerror = (error) => {
            clearTimeout(timeout);
            console.error(}WebSocket:\code{, error);
            reject(error);
        };
        
        ws.onclose = (event) => {
            if (event.code !== 1000) {
                console.warn(}: ${event.code} ${event.reason}\code{);
            }
        };
    });
}

// 
async function establishConnection(stationId, maxRetries = 3) {
    for (let attempt = 1; attempt <= maxRetries; attempt++) {
        try {
            const connection = await createRealtimeConnection(stationId);
            return connection;
        } catch (error) {
            console.error(}${attempt}:\code{, error.message);
            
            if (attempt < maxRetries) {
                const delay = 1000 } attempt; // 
                console.log(}${delay}ms...\code{);
                await new Promise(resolve => setTimeout(resolve, delay));
            }
        }
    }
    
    throw new Error(}\$\{maxRetries}\code{);
}

\begin{lstlisting}
#### 4. 

JavaScriptconsole.log

\important{}

consoleJavaScript

\end{lstlisting}

javascript // class WaterSystemDebugger \{ constructor(debugMode =
false) \{ this.debugMode = debugMode; this.performanceMarkers = new
Map(); this.debugLogs = {[}{]}; this.errorHistory = {[}{]}; }

\begin{verbatim}
// 1. 
debugLog(message, data = null, level = 'info') {
    if (!this.debugMode) return;
    
    const timestamp = new Date().toISOString();
    const debugEntry = { timestamp, level, message, data };
    
    this.debugLogs.push(debugEntry);
    
    // 
    switch (level) {
        case 'error':
            console.error(} [${timestamp}] ${message}\code{, data);
            break;
        case 'warn':
            console.warn(}\faExclamationTriangle [${timestamp}] ${message}\code{, data);
            break;
        case 'info':
            console.info(}ℹ [${timestamp}] ${message}\code{, data);
            break;
        case 'debug':
            console.debug(} [${timestamp}] ${message}\code{, data);
            break;
        case 'table':
            console.table(data);
            break;
        case 'group':
            console.group(message);
            if (data) console.log(data);
            break;
        case 'groupEnd':
            console.groupEnd();
            break;
        default:
            console.log(}\faPencil [${timestamp}] ${message}\code{, data);
    }
}

// 2. 
startPerformanceMarker(name) {
    if (!this.debugMode) return;
    
    this.performanceMarkers.set(name, {
        startTime: performance.now(),
        startMemory: performance.memory ? performance.memory.usedJSHeapSize : null
    });
    
    console.time(name);
}

endPerformanceMarker(name) {
    if (!this.debugMode) return;
    
    const marker = this.performanceMarkers.get(name);
    if (!marker) {
        console.warn(} "${name}" \code{);
        return;
    }
    
    const endTime = performance.now();
    const duration = endTime - marker.startTime;
    const endMemory = performance.memory ? performance.memory.usedJSHeapSize : null;
    const memoryDiff = endMemory && marker.startMemory ? 
                      endMemory - marker.startMemory : null;
    
    console.timeEnd(name);
    
    const perfInfo = {
        duration: }${duration.toFixed(2)}ms\code{,
        memoryChange: memoryDiff ? }${(memoryDiff / 1024 / 1024).toFixed(2)}MB\code{ : 'N/A'
    };
    
    this.debugLog(} "${name}" \code{, perfInfo, 'info');
    
    this.performanceMarkers.delete(name);
    return { duration, memoryChange: memoryDiff };
}

// 3. 
conditionalBreakpoint(condition, message = '') {
    if (condition && this.debugMode) {
        console.log(} : ${message}\code{);
        debugger; // 
    }
}

// 4. 
traceFunction(fn, name) {
    if (!this.debugMode) return fn;
    
    return (...args) => {
        this.debugLog(}: ${name}\code{, { args }, 'group');
        
        try {
            const result = fn.apply(this, args);
            
            // 
            if (result && typeof result.then === 'function') {
                return result
                    .then(value => {
                        this.debugLog(}: ${name}\code{, { result: value });
                        this.debugLog('', null, 'groupEnd');
                        return value;
                    })
                    .catch(error => {
                        this.debugLog(}: ${name}\code{, { error: error.message }, 'error');
                        this.debugLog('', null, 'groupEnd');
                        throw error;
                    });
            } else {
                this.debugLog(}: ${name}\code{, { result });
                this.debugLog('', null, 'groupEnd');
                return result;
            }
        } catch (error) {
            this.debugLog(}: ${name}\code{, { error: error.message }, 'error');
            this.debugLog('', null, 'groupEnd');
            throw error;
        }
    };
}

// 5. 
createWatchedObject(obj, name = 'object') {
    if (!this.debugMode) return obj;
    
    return new Proxy(obj, {
        get: (target, property) => {
            const value = target[property];
            this.debugLog(}: ${name}.${String(property)}\code{, { value });
            return value;
        },
        
        set: (target, property, value) => {
            const oldValue = target[property];
            target[property] = value;
            this.debugLog(}: ${name}.${String(property)}\code{, 
                { oldValue, newValue: value });
            return true;
        }
    });
}

// 6. 
traceDataFlow(data, description) {
    if (!this.debugMode) return data;
    
    const traceId = }trace_${Date.now()}_${Math.random().toString(36).substr(2, 9)}\code{;
    
    this.debugLog(}: ${description}\code{, { traceId, data }, 'group');
    
    // 
    if (typeof data === 'object' && data !== null) {
        Object.defineProperty(data, '__traceId', {
            value: traceId,
            writable: false,
            enumerable: false
        });
    }
    
    this.debugLog('', null, 'groupEnd');
    return data;
}

// 7. 
captureErrorContext(error, context = {}) {
    const errorInfo = {
        message: error.message,
        stack: error.stack,
        type: error.constructor.name,
        timestamp: new Date().toISOString(),
        context: {
            userAgent: navigator.userAgent,
            url: window.location.href,
            ...context
        },
        systemState: this.captureSystemState()
    };
    
    this.errorHistory.push(errorInfo);
    
    // 
    if (this.errorHistory.length > 50) {
        this.errorHistory = this.errorHistory.slice(-25);
    }
    
    this.debugLog('', errorInfo, 'error');
    return errorInfo;
}

// 8. 
captureSystemState() {
    return {
        timestamp: new Date().toISOString(),
        memory: performance.memory ? {
            used: }${(performance.memory.usedJSHeapSize / 1024 / 1024).toFixed(2)}MB\code{,
            total: }${(performance.memory.totalJSHeapSize / 1024 / 1024).toFixed(2)}MB\code{,
            limit: }${(performance.memory.jsHeapSizeLimit / 1024 / 1024).toFixed(2)}MB\code{
        } : 'N/A',
        performanceMarkers: Array.from(this.performanceMarkers.keys()),
        debugLogsCount: this.debugLogs.length,
        errorHistoryCount: this.errorHistory.length
    };
}

// 9. 
generateDebugReport() {
    return {
        reportGeneratedAt: new Date().toISOString(),
        systemState: this.captureSystemState(),
        recentLogs: this.debugLogs.slice(-20),
        errorHistory: this.errorHistory,
        summary: {
            totalLogs: this.debugLogs.length,
            totalErrors: this.errorHistory.length,
            activePerformanceMarkers: this.performanceMarkers.size
        }
    };
}

// 10. 
createDebugPanel() {
    if (!this.debugMode || typeof document === 'undefined') return;
    
    const panel = document.createElement('div');
    panel.style.cssText = }
        position: fixed;
        top: 10px;
        right: 10px;
        width: 300px;
        max-height: 400px;
        background: rgba(0, 0, 0, 0.9);
        color: #00ff00;
        font-family: monospace;
        font-size: 12px;
        padding: 10px;
        border-radius: 5px;
        overflow-y: auto;
        z-index: 10000;
        border: 1px solid #333;
    \code{;
    
    panel.innerHTML = }
        <div style="font-weight: bold; margin-bottom: 10px;">
             Water System Debug Panel
            <button onclick="this.parentElement.parentElement.remove()" 
                    style="float: right; background: red; color: white; border: none; cursor: pointer;">×</button>
        </div>
        <div id="debugPanelContent">Loading...</div>
    \code{;
    
    document.body.appendChild(panel);
    
    // 
    const updatePanel = () => {
        const content = document.getElementById('debugPanelContent');
        if (!content) return;
        
        const state = this.captureSystemState();
        content.innerHTML = }
            <div>: ${state.memory.used || 'N/A'}</div>
            <div>: ${state.performanceMarkers.length}</div>
            <div>: ${state.debugLogsCount}</div>
            <div>: ${state.errorHistoryCount}</div>
            <div style="margin-top: 10px; font-size: 11px;">
                :<br>
                ${this.debugLogs.slice(-5).map(log => 
                    \code{${log.level}: ${log.message}}
                ).join('<br>')}
            </div>
        \code{;
    };
    
    updatePanel();
    const updateInterval = setInterval(updatePanel, 1000);
    
    // 
    panel.addEventListener('remove', () => {
        clearInterval(updateInterval);
    });
}
\end{verbatim}

}

// const debugger = new WaterSystemDebugger(true);

// debugger.createDebugPanel();

// function processWaterStationData(stationId, data) \{ const
trackedData = debugger.traceDataFlow(data, }\$\{stationId}\code{);
    
    debugger.startPerformanceMarker('dataProcessing');
    
    try {
        // 
        debugger.conditionalBreakpoint(
            data.waterLevel > 25, 
            }${stationId}: ${data.waterLevel}m\code{
        );
        
        // 
        const processedData = {
            stationId,
            ...data,
            processedAt: new Date(),
            alert: data.waterLevel > 20
        };
        
        debugger.debugLog('', { stationId, processedData });
        
        return processedData;
        
    } catch (error) {
        debugger.captureErrorContext(error, { stationId, data });
        throw error;
    } finally {
        debugger.endPerformanceMarker('dataProcessing');
    }
}

\begin{lstlisting}
comprehensiveJavaScript

###  - 

\important{JavaScript}

#### 1. 

\end{lstlisting}

javascript //

// \faTimes  let d = new Date(); let wl = 15.2; let temp = 25.3; let
calc = (a, b) =\textgreater{} a + b;

// \faCheck  let currentTimestamp = new Date(); let waterLevelInMeters =
15.2; let temperatureInCelsius = 25.3; let calculateAverageValue =
(firstValue, secondValue) =\textgreater{} firstValue + secondValue;

// class WaterLevelMonitor \{ constructor(stationId, alertThresholds) \{
// this.stationId = stationId; this.alertThresholds = alertThresholds;
this.currentWaterLevel = 0; this.lastUpdateTimestamp = null;

\begin{verbatim}
    // 
    this.MAX_WATER_LEVEL = 50;
    this.MIN_WATER_LEVEL = 0;
    this.DEFAULT_UPDATE_INTERVAL = 5000; // 5
    
    // 
    this._internalBuffer = [];
    this._lastValidReading = null;
}

// 
updateWaterLevel(newLevel) {
    if (!this._isValidWaterLevel(newLevel)) {
        throw new Error(}: ${newLevel}\code{);
    }
    
    this.currentWaterLevel = newLevel;
    this.lastUpdateTimestamp = new Date();
    
    this._checkAlertConditions(newLevel);
    this._logWaterLevelChange(newLevel);
}

// ishascan
isAlertTriggered() {
    return this.currentWaterLevel > this.alertThresholds.warning;
}

hasValidData() {
    return this.lastUpdateTimestamp !== null && 
           this._isValidWaterLevel(this.currentWaterLevel);
}

canTriggerAlert() {
    return this.hasValidData() && this.isAlertTriggered();
}

// get
getFormattedWaterLevel() {
    return }${this.currentWaterLevel.toFixed(2)} \code{;
}

getCurrentStatus() {
    if (!this.hasValidData()) {
        return {
            status: 'NO_DATA',
            message: '',
            level: this.currentWaterLevel
        };
    }
    
    if (this.currentWaterLevel > this.alertThresholds.critical) {
        return {
            status: 'CRITICAL',
            message: '',
            level: this.currentWaterLevel
        };
    }
    
    if (this.currentWaterLevel > this.alertThresholds.warning) {
        return {
            status: 'WARNING', 
            message: '',
            level: this.currentWaterLevel
        };
    }
    
    return {
        status: 'NORMAL',
        message: '',
        level: this.currentWaterLevel
    };
}

// 
_isValidWaterLevel(level) {
    return typeof level === 'number' && 
           !isNaN(level) && 
           level >= this.MIN_WATER_LEVEL && 
           level <= this.MAX_WATER_LEVEL;
}

_checkAlertConditions(waterLevel) {
    const previousLevel = this._lastValidReading;
    
    if (previousLevel && Math.abs(waterLevel - previousLevel) > 5) {
        console.warn(}: ${previousLevel}${waterLevel}\code{);
    }
    
    if (waterLevel > this.alertThresholds.warning) {
        this._triggerAlert('WARNING', waterLevel);
    }
    
    if (waterLevel > this.alertThresholds.critical) {
        this._triggerAlert('CRITICAL', waterLevel);
    }
}

_triggerAlert(level, waterLevel) {
    const alertMessage = }${level}${this.stationId}: ${waterLevel}\code{;
    console.log(alertMessage);
    
    // 
    // 
}

_logWaterLevelChange(newLevel) {
    const timestamp = new Date().toISOString();
    console.log(}[${timestamp}] ${this.stationId}: ${newLevel}\code{);
    
    this._lastValidReading = newLevel;
}
\end{verbatim}

}

// // \faCheck  function calculateDailyAverageWaterLevel(dailyReadings)
\{ const validReadings = dailyReadings.filter(reading =\textgreater{}
reading !== null \&\& reading !== undefined);

\begin{verbatim}
if (validReadings.length === 0) {
    return null;
}

const totalWaterLevel = validReadings.reduce((sum, reading) => sum + reading.waterLevel, 0);
return totalWaterLevel / validReadings.length;
\end{verbatim}

}

function formatWaterLevelForDisplay(waterLevel, unit = '`) \{ if
(waterLevel === null \textbar\textbar{} waterLevel === undefined) \{
return''; }

\begin{verbatim}
return }${waterLevel.toFixed(2)} ${unit}\code{;
\end{verbatim}

}

function validateStationConfiguration(stationConfig) \{ const
requiredFields = {[}`stationId', `name', `location',
`alertThresholds'{]};

\begin{verbatim}
for (const field of requiredFields) {
    if (!stationConfig[field]) {
        throw new Error(}: ${field}\code{);
    }
}

return true;
\end{verbatim}

}

// const WATER\_LEVEL\_ALERT\_LEVELS = \{ NORMAL: `normal', WARNING:
`warning', CRITICAL: `critical', EMERGENCY: `emergency' };

const STATION\_TYPES = \{ RIVER\_MAIN\_STREAM: `river\_main\_stream',
RIVER\_TRIBUTARY: `river\_tributary', RESERVOIR\_INLET:
`reservoir\_inlet', RESERVOIR\_OUTLET: `reservoir\_outlet', GROUNDWATER:
`groundwater' };

const DEFAULT\_CONFIG = \{ updateInterval: 5000, // maxDataAge: 300000,
// 5 alertDelayTime: 30000, // 30 maxRetryAttempts: 3, //
retryIntervalMs: 2000 // };

\begin{lstlisting}
#### 2. 

\end{lstlisting}

javascript //

// 1. ES6 // water-station.js - export class WaterStation \{
constructor(config) \{ this.config =
this.\_validateAndNormalizeConfig(config); this.sensors = new Map();
this.dataBuffer = {[}{]}; this.alertHandlers = new Set(); this.status =
`inactive'; }

\begin{verbatim}
// API
async initialize() {
    try {
        await this._setupSensors();
        await this._establishConnections();
        this.status = 'active';
        console.log(} ${this.config.id} \code{);
    } catch (error) {
        console.error(} ${this.config.id} :\code{, error.message);
        this.status = 'error';
        throw error;
    }
}

addSensor(sensorType, sensorConfig) {
    const sensor = SensorFactory.create(sensorType, sensorConfig);
    this.sensors.set(sensor.id, sensor);
    return sensor.id;
}

async collectData() {
    const collectedData = {};
    
    for (const [sensorId, sensor] of this.sensors) {
        try {
            const reading = await sensor.read();
            collectedData[sensorId] = reading;
        } catch (error) {
            console.error(} ${sensorId} :\code{, error.message);
            collectedData[sensorId] = null;
        }
    }
    
    return this._processCollectedData(collectedData);
}

// 
_validateAndNormalizeConfig(config) {
    const required = ['id', 'name', 'location'];
    for (const field of required) {
        if (!config[field]) {
            throw new Error(}: ${field}\code{);
        }
    }
    
    return {
        ...DEFAULT_STATION_CONFIG,
        ...config,
        location: this._normalizeLocation(config.location)
    };
}

_normalizeLocation(location) {
    return {
        latitude: parseFloat(location.latitude),
        longitude: parseFloat(location.longitude), 
        altitude: parseFloat(location.altitude || 0)
    };
}

async _setupSensors() {
    // 
}

async _establishConnections() {
    // 
}

_processCollectedData(rawData) {
    // 
    return {
        stationId: this.config.id,
        timestamp: new Date().toISOString(),
        data: rawData,
        quality: this._assessDataQuality(rawData)
    };
}

_assessDataQuality(data) {
    // 
    return { score: 100, level: 'excellent' };
}
\end{verbatim}

}

// sensor-factory.js - class BaseSensor \{ constructor(id, config) \{
this.id = id; this.config = config; this.calibrationData = null;
this.lastReading = null; }

\begin{verbatim}
async read() {
    throw new Error('read');
}

async calibrate(calibrationData) {
    this.calibrationData = calibrationData;
}

_applyCalibration(rawValue) {
    if (!this.calibrationData) return rawValue;
    
    // 
    return rawValue \highlight{ this.calibrationData.factor + this.calibrationData.offset;
}
\end{verbatim}

}

class WaterLevelSensor extends BaseSensor \{ async read() \{ try \{ //
const rawValue = Math.random() } 30 + 10; // 10-40 const
calibratedValue = this.\_applyCalibration(rawValue);

\begin{verbatim}
        this.lastReading = {
            value: calibratedValue,
            unit: 'meters',
            timestamp: new Date(),
            quality: this._assessReadingQuality(calibratedValue)
        };
        
        return this.lastReading;
    } catch (error) {
        throw new Error(}: ${error.message}\code{);
    }
}

_assessReadingQuality(value) {
    // 
    if (value < 0 || value > 50) return 'invalid';
    if (this.lastReading) {
        const change = Math.abs(value - this.lastReading.value);
        if (change > 5) return 'suspect'; // 
    }
    return 'good';
}
\end{verbatim}

}

class TemperatureSensor extends BaseSensor \{ async read() \{ try \{
const rawValue = Math.random() \highlight{ 30 + 5; // 5-35
            const calibratedValue = this._applyCalibration(rawValue);
            
            this.lastReading = {
                value: calibratedValue,
                unit: 'celsius',
                timestamp: new Date(),
                quality: this._assessReadingQuality(calibratedValue)
            };
            
            return this.lastReading;
        } catch (error) \{ throw new Error(}:
\$\{error.message}\code{);
        } }

\begin{verbatim}
_assessReadingQuality(value) {
    if (value < -20 || value > 60) return 'invalid';
    return 'good';
}
\end{verbatim}

}

export const SensorFactory = \{ create(type, config) \{ const sensorId =
}\({type}_\)\{Date.now()}\_\$\{Math.random().toString(36).substr(2,
9)}\code{;
        
        switch (type) {
            case 'water_level':
                return new WaterLevelSensor(sensorId, config);
            case 'temperature':
                return new TemperatureSensor(sensorId, config);
            default:
                throw new Error(}: ${type}\code{);
        }
    },

\begin{verbatim}
getSupportedTypes() {
    return ['water_level', 'temperature', 'ph', 'turbidity', 'flow_rate'];
}
\end{verbatim}

};

// data-validator.js - export class DataValidator \{ constructor() \{
this.rules = new Map(); this.\_initializeDefaultRules(); }

\begin{verbatim}
addRule(fieldName, validator) {
    if (!this.rules.has(fieldName)) {
        this.rules.set(fieldName, []);
    }
    this.rules.get(fieldName).push(validator);
}

validate(data) {
    const errors = [];
    const warnings = [];
    
    for (const [fieldName, validators] of this.rules) {
        if (data[fieldName] === undefined || data[fieldName] === null) {
            continue; // 
        }
        
        for (const validator of validators) {
            const result = validator(data[fieldName], data);
            
            if (!result.valid) {
                if (result.severity === 'error') {
                    errors.push({
                        field: fieldName,
                        message: result.message,
                        value: data[fieldName]
                    });
                } else if (result.severity === 'warning') {
                    warnings.push({
                        field: fieldName,
                        message: result.message,
                        value: data[fieldName]
                    });
                }
            }
        }
    }
    
    return {
        valid: errors.length === 0,
        errors,
        warnings
    };
}

_initializeDefaultRules() {
    // 
    this.addRule('waterLevel', (value) => {
        if (typeof value !== 'number' || isNaN(value)) {
            return { valid: false, severity: 'error', message: '' };
        }
        
        if (value < 0) {
            return { valid: false, severity: 'error', message: '' };
        }
        
        if (value > 100) {
            return { valid: false, severity: 'warning', message: '' };
        }
        
        return { valid: true };
    });
    
    // 
    this.addRule('temperature', (value) => {
        if (typeof value !== 'number' || isNaN(value)) {
            return { valid: false, severity: 'error', message: '' };
        }
        
        if (value < -50 || value > 80) {
            return { valid: false, severity: 'error', message: '(-50°C80°C)' };
        }
        
        return { valid: true };
    });
    
    // 
    this.addRule('timestamp', (value) => {
        const timestamp = new Date(value);
        
        if (isNaN(timestamp.getTime())) {
            return { valid: false, severity: 'error', message: '' };
        }
        
        const now = new Date();
        const maxAge = 24 } 60 \highlight{ 60 } 1000; // 24
        
        if (now - timestamp > maxAge) {
            return { valid: false, severity: 'warning', message: '' };
        }
        
        if (timestamp > now) {
            return { valid: false, severity: 'warning', message: '' };
        }
        
        return { valid: true };
    });
}
\end{verbatim}

}

// import \{ WaterStation, SensorFactory, DataValidator } from
`./water-monitoring';

const stationConfig = \{ id: `WS001', name: '\,', location: \{ latitude:
30.7, longitude: 111.3, altitude: 45.5 } };

const waterStation = new WaterStation(stationConfig); const validator =
new DataValidator();

// waterStation.initialize().then(() =\textgreater{} \{ // const
waterLevelSensorId = waterStation.addSensor(`water\_level', \{ range:
{[}0, 50{]}, accuracy: 0.01 });

\begin{verbatim}
const temperatureSensorId = waterStation.addSensor('temperature', {
    range: [-10, 50],
    accuracy: 0.1 
});

// 
setInterval(async () => {
    try {
        const data = await waterStation.collectData();
        const validationResult = validator.validate(data.data);
        
        if (validationResult.valid) {
            console.log(':', data);
        } else {
            console.error(':', validationResult.errors);
            console.warn(':', validationResult.warnings);
        }
    } catch (error) {
        console.error(':', error.message);
    }
}, 10000); // 10
\end{verbatim}

});

\begin{lstlisting}
#### 3. 

\end{lstlisting}

javascript //

// 1. class EfficientDataProcessor \{ constructor() \{ // GC
this.dataPointPool = {[}{]}; this.processingResults = new Map();

\begin{verbatim}
    // WeakMap
    this.stationMetadata = new WeakMap();
    
    // 
    this.validationPatterns = {
        stationId: /^WS\d{3}$/,
        timestamp: /^\d{4}-\d{2}-\d{2}T\d{2}:\d{2}:\d{2}\.\d{3}Z$/
    };
}

// 
createDataPoint(stationId, value, timestamp) {
    let dataPoint = this.dataPointPool.pop();
    
    if (!dataPoint) {
        dataPoint = {
            stationId: null,
            value: null,
            timestamp: null,
            processed: false
        };
    }
    
    // 
    dataPoint.stationId = stationId;
    dataPoint.value = value;
    dataPoint.timestamp = timestamp;
    dataPoint.processed = false;
    
    return dataPoint;
}

// 
recycleDataPoint(dataPoint) {
    // 
    dataPoint.stationId = null;
    dataPoint.value = null;
    dataPoint.timestamp = null;
    dataPoint.processed = false;
    
    this.dataPointPool.push(dataPoint);
}

// 
processBatchData(dataPoints) {
    const batchSize = 1000;
    const results = [];
    
    for (let i = 0; i < dataPoints.length; i += batchSize) {
        const batch = dataPoints.slice(i, i + batchSize);
        const batchResults = this._processBatch(batch);
        results.push(...batchResults);
    }
    
    return results;
}

_processBatch(batch) {
    const results = new Array(batch.length); // 
    
    for (let i = 0; i < batch.length; i++) {
        const dataPoint = batch[i];
        
        // 
        results[i] = {
            stationId: dataPoint.stationId,
            processedValue: dataPoint.value \highlight{ 1.1 + 0.5, // 
            timestamp: dataPoint.timestamp,
            quality: dataPoint.value > 0 && dataPoint.value < 100 ? 'good' : 'poor'
        };
    }
    
    return results;
}
\end{verbatim}

}

// 2. class OptimizedStationManager \{ constructor() \{ // MapO(1)
this.stationById = new Map(); this.stationsByLocation = new Map();

\begin{verbatim}
    // Set
    this.activeStations = new Set();
    
    // 
    this.locationIndex = new Map(); // 
    this.typeIndex = new Map();     // 
}

addStation(station) {
    // 
    this.stationById.set(station.id, station);
    
    // 
    const locationKey = }${station.location.latitude.toFixed(3)},${station.location.longitude.toFixed(3)}\code{;
    if (!this.locationIndex.has(locationKey)) {
        this.locationIndex.set(locationKey, new Set());
    }
    this.locationIndex.get(locationKey).add(station.id);
    
    // 
    if (!this.typeIndex.has(station.type)) {
        this.typeIndex.set(station.type, new Set());
    }
    this.typeIndex.get(station.type).add(station.id);
    
    // 
    this.activeStations.add(station.id);
}

// 
findNearbyStations(latitude, longitude, radiusKm) {
    const nearby = [];
    const targetLat = parseFloat(latitude);
    const targetLng = parseFloat(longitude);
    
    // 
    for (const [locationKey, stationIds] of this.locationIndex) {
        const [lat, lng] = locationKey.split(',').map(parseFloat);
        
        // 
        const roughDistance = Math.abs(lat - targetLat) + Math.abs(lng - targetLng);
        if (roughDistance <= radiusKm } 0.02) { // 1=111km
            
            for (const stationId of stationIds) {
                const station = this.stationById.get(stationId);
                const exactDistance = this._calculateDistance(
                    targetLat, targetLng, 
                    station.location.latitude, 
                    station.location.longitude
                );
                
                if (exactDistance <= radiusKm) {
                    nearby.push({
                        station,
                        distance: exactDistance
                    });
                }
            }
        }
    }
    
    // 
    return nearby.sort((a, b) => a.distance - b.distance);
}

_calculateDistance(lat1, lng1, lat2, lng2) {
    // Haversine
    const R = 6371; // (km)
    const dLat = this._toRadians(lat2 - lat1);
    const dLng = this._toRadians(lng2 - lng1);
    
    const a = Math.sin(dLat/2) \highlight{ Math.sin(dLat/2) +
             Math.cos(this._toRadians(lat1)) } Math.cos(this._toRadians(lat2)) \highlight{
             Math.sin(dLng/2) } Math.sin(dLng/2);
    
    const c = 2 \highlight{ Math.atan2(Math.sqrt(a), Math.sqrt(1-a));
    return R } c;
}

_toRadians(degrees) {
    return degrees \highlight{ (Math.PI / 180);
}
\end{verbatim}

}

// 3. class AsyncDataManager \{ constructor(concurrencyLimit = 10) \{
this.concurrencyLimit = concurrencyLimit; this.requestQueue = {[}{]};
this.activeRequests = 0; }

\begin{verbatim}
// 
async fetchStationData(stationId) {
    return new Promise((resolve, reject) => {
        this.requestQueue.push({ stationId, resolve, reject });
        this._processQueue();
    });
}

async _processQueue() {
    if (this.activeRequests >= this.concurrencyLimit || this.requestQueue.length === 0) {
        return;
    }
    
    const request = this.requestQueue.shift();
    this.activeRequests++;
    
    try {
        const data = await this._doFetch(request.stationId);
        request.resolve(data);
    } catch (error) {
        request.reject(error);
    } finally {
        this.activeRequests--;
        // 
        setImmediate(() => this._processQueue());
    }
}

async _doFetch(stationId) {
    // 
    const response = await fetch(}/api/stations/${stationId}/data\code{);
    if (!response.ok) {
        throw new Error(}HTTP ${response.status}: ${response.statusText}\code{);
    }
    return await response.json();
}

// 
async fetchMultipleStationsData(stationIds) {
    // 
    const batchSize = 50;
    const batches = [];
    
    for (let i = 0; i < stationIds.length; i += batchSize) {
        batches.push(stationIds.slice(i, i + batchSize));
    }
    
    // 
    const allResults = [];
    
    for (const batch of batches) {
        const batchPromises = batch.map(stationId => 
            this.fetchStationData(stationId).catch(error => ({ error, stationId }))
        );
        
        const batchResults = await Promise.all(batchPromises);
        allResults.push(...batchResults);
    }
    
    return allResults;
}
\end{verbatim}

}

// 4. class MemoryEfficientCache \{ constructor(maxSize = 1000, ttlMs =
300000) \{ // 5TTL this.cache = new Map(); this.accessTimes = new Map();
this.maxSize = maxSize; this.ttlMs = ttlMs;

\begin{verbatim}
    // 
    this.cleanupInterval = setInterval(() => {
        this._cleanup();
    }, 60000); // 
}

set(key, value) {
    const now = Date.now();
    
    // 
    if (this.cache.size >= this.maxSize && !this.cache.has(key)) {
        this._evictLeastRecentlyUsed();
    }
    
    this.cache.set(key, {
        value,
        timestamp: now,
        accessCount: 1
    });
    
    this.accessTimes.set(key, now);
}

get(key) {
    const item = this.cache.get(key);
    
    if (!item) {
        return null;
    }
    
    const now = Date.now();
    
    // 
    if (now - item.timestamp > this.ttlMs) {
        this.cache.delete(key);
        this.accessTimes.delete(key);
        return null;
    }
    
    // 
    item.accessCount++;
    this.accessTimes.set(key, now);
    
    return item.value;
}

_evictLeastRecentlyUsed() {
    let leastRecentKey = null;
    let leastRecentTime = Date.now();
    
    for (const [key, time] of this.accessTimes) {
        if (time < leastRecentTime) {
            leastRecentTime = time;
            leastRecentKey = key;
        }
    }
    
    if (leastRecentKey) {
        this.cache.delete(leastRecentKey);
        this.accessTimes.delete(leastRecentKey);
    }
}

_cleanup() {
    const now = Date.now();
    const expiredKeys = [];
    
    for (const [key, item] of this.cache) {
        if (now - item.timestamp > this.ttlMs) {
            expiredKeys.push(key);
        }
    }
    
    for (const key of expiredKeys) {
        this.cache.delete(key);
        this.accessTimes.delete(key);
    }
    
    console.log(} ${expiredKeys.length} \code{);
}

destroy() {
    if (this.cleanupInterval) {
        clearInterval(this.cleanupInterval);
    }
    
    this.cache.clear();
    this.accessTimes.clear();
}
\end{verbatim}

}

// const processor = new EfficientDataProcessor(); const stationManager
= new OptimizedStationManager(); const dataManager = new
AsyncDataManager(5); // 5 const cache = new MemoryEfficientCache(500,
180000); // 5003

// console.time('\,');

const testData = Array.from(\{ length: 10000 }, (\_, i) =\textgreater{}
(\{ stationId: }WS\$\{String(i \% 100).padStart(3, `0')}\code{,
    value: Math.random() } 100, timestamp: new Date().toISOString()
}));

const results = processor.processBatchData(testData);
console.timeEnd('\,');

console.log(} \$\{results.length} \code{);
\begin{lstlisting}
comprehensiveJavaScript

!!! info "JavaScript"
    
    \important{}
    - - 
    - \important{}
    - - 
    - \important{}
    - API
    - WebSocket
    - ## 4.4.1 ES6+

ECMAScript 2015ES6JavaScript\important{JavaScript}

### 

JavaScript}var\code{ES6}let\code{}const\code{\important{}

\end{lstlisting}javascript
// var
for (var i = 0; i < monitoringStations.length; i++) {
    // var
    setTimeout(function() {
        console.log(':', i); // 
    // ...  ...
    
    return processedData;
}
\begin{lstlisting}
### 

\important{}Template LiteralsES6HTML

\end{lstlisting}javascript
// 
function generateAlertMessage(station, waterLevel, threshold) {
    // 
    const alertMessage = }
  
    // ...  ...
    
    return templates[language] || templates.zh;
}
\begin{lstlisting}
### 

\important{}ES6API

API

\end{lstlisting}javascript
// 
function extractStationData(monitoringResponse) {
    //  - 
    const {
        stationInfo: {
    // ...  ...
        ? validValues.reduce((sum, val) => sum + val, 0) / validValues.length
        : null;
}
\begin{lstlisting}
### 

\important{}ES6\code{this}JavaScript

\code{map}\code{filter}\code{reduce}

\end{lstlisting}javascript
// 
class WaterMonitoringDataProcessor {
    constructor(rawData) {
        this.rawData = rawData;
    }
    // ...  ...
        }
    }
}
\begin{lstlisting}
## 4.4.2 Promise

JavaScript\important{Promise}ES6

Promise

### Promise

\important{Promise}pendingfulfilledrejectedPromise

Promise\important{Promise}\code{.then()}Promise

\end{lstlisting}javascript
// Promise
class WaterMonitoringAPI {
    constructor(baseURL = '/api/monitoring') {
        this.baseURL = baseURL;
        this.cache = new Map(); // 
    // ...  ...
        };
    }
}

\begin{lstlisting}
### async/await

\important{async/await}ES2017Promise\code{async}Promise\code{await}asyncPromise

async/awaitPromise

\end{lstlisting}

javascript // async/await class ModernWaterMonitoringService \{
constructor() \{ this.apiClient = new WaterMonitoringAPI();
this.eventEmitter = new EventEmitter(); // \ldots{} \ldots{}
showErrorNotification('\,'); } }

\begin{lstlisting}
### 

\important{}try-catchasync/await

\end{lstlisting}

javascript // class ErrorHandler \{ constructor() \{ this.errorTypes =
\{ NETWORK\_ERROR: `network', // \ldots{} \ldots{} return new
Promise(resolve =\textgreater{} setTimeout(resolve, ms)); } }

\begin{lstlisting}
### Promise

PromiseJavaScriptPromise\code{Promise.all()}\code{Promise.allSettled()}\code{Promise.race()}\important{Promise}\important{}

\end{lstlisting}

javascript // Promise class AdvancedMonitoringOperations \{
constructor() \{ this.concurrencyLimit = 5; // this.batchSize = 10; //
// \ldots{} \ldots{} }); } }

\begin{lstlisting}
## 4.4.3 DOM

DOMWebHTMLJavaScriptDOM APIDOM\important{DOM}

DOMAPI\code{querySelector}\code{classList}\code{dataset}DOMDOM

### DOM

\code{document.getElementById()}\code{getElementsByClassName()}\important{DOMAPI}\code{querySelector()}\code{querySelectorAll()}CSS

API

\end{lstlisting}

javascript // DOM class MonitoringDOMManager \{ constructor() \{
this.container = document.querySelector(`.monitoring-dashboard');
this.cache = new Map(); // DOM // \ldots{} \ldots{} return new
Date(timestamp).toLocaleString(`zh-CN'); } }

\begin{lstlisting}
### 

\important{}JavaScript

\important{}\important{}

\end{lstlisting}

javascript // class MonitoringEventSystem \{ constructor(container) \{
this.container = container; this.eventHandlers = new Map(); // \ldots{}
\ldots{} console.log('\,'); // });

\begin{lstlisting}
## 4.4.4 

\important{}JavaScriptJavaScriptES6 Modules

\important{}JavaScript

### ES6

ES6\code{import}\code{export}

\important{}\important{}

\end{lstlisting}

javascript // utils/dataValidator.js - export class DataValidator \{
static validateWaterLevel(level) \{ if (typeof level !== `number'
\textbar\textbar{} isNaN(level)) \{ throw new Error('`); // \ldots{}
\ldots{} MISSING\_FIELD:'`, INVALID\_FORMAT:'' };

\begin{lstlisting}
\important{API}HTTPAPI

\end{lstlisting}

javascript // services/apiService.js - API import \{ DataValidator,
ERROR\_MESSAGES } from `../utils/dataValidator.js';

export class APIService \{ constructor(baseURL = `/api', options = \{})
\{ // \ldots{} \ldots{}

// API export default new APIService();

\begin{lstlisting}
\important{}

\end{lstlisting}

javascript // components/chartRenderer.js - import apiService from
`../services/apiService.js';

export class ChartRenderer \{ constructor() \{ // \ldots{} \ldots{}

// export default new ChartRenderer();

\begin{lstlisting}
### 

JavaScript\important{}\important{}\important{}\important{}

\important{}JavaScript

\end{lstlisting}

javascript // debug/debugUtils.js - export class DebugUtils \{
constructor() \{ this.isDebugMode = this.checkDebugMode();
this.performanceMarks = new Map(); // \ldots{} \ldots{}

\begin{verbatim}
return descriptor;
\end{verbatim}

}

\begin{lstlisting}
\important{}

\end{lstlisting}

javascript // main.js - import apiService from
`./services/apiService.js'; import chartRenderer from
`./components/chartRenderer.js'; import debugUtils, \{
performanceMonitor } from `./debug/debugUtils.js'; import \{
DataValidator } from `./utils/dataValidator.js'; // \ldots{} \ldots{}

// export default WaterMonitoringApp;

\begin{lstlisting}
## 

JavaScriptJavaScript

### 

||  ||
|----------|----------|----------------------|
|\important{ES6+}| let/const ||
|\important{}| Promiseasync/await ||
|\important{DOM}| API ||
|\important{}| ES6 ||

### 

1. \important{}ES6+
2. \important{}Promiseasync/await
3. \important{}DOM
4. \important{}

### 

JavaScriptVue.jsVue.jsVue.js

JavaScriptVue.js
\end{lstlisting}

\texttt{\textbackslash{}code\{\ }}

\hypertarget{vue}{%
\section{4.5 Vue}\label{vue}}

Vue.jsJavaScriptVueAPIVue.js

jQuery DOMVue.jsMVVMModel-View-ViewModelDOM

Vue.jsVue.jsMVVMVue.js

!!! info ``Vue.js''

\begin{verbatim}
Vue.js
\end{verbatim}

\hypertarget{section-29}{%
\subsection{4.5.1}\label{section-29}}

\hypertarget{section-30}{%
\subsubsection{}\label{section-30}}

WebHTMLCSSJavaScript

\important{Web}

\begin{itemize}
\item
  \important{DOM}DOM

  WebDOM

  ``\,``DOM

  \textbackslash begin\{lstlisting}{[}style=javascript{]} // function
  deleteUser(userId) \{ // 1. const userRow =
  document.getElementById(`user-' + userId);
  userRow.parentNode.removeChild(userRow);

  // 2. const totalCount = document.getElementById(`total-users'); const
  currentCount = parseInt(totalCount.textContent) - 1;
  totalCount.textContent = currentCount;

  // 3. const currentPage = Math.ceil(currentCount / pageSize);
  document.getElementById(`current-page').textContent = currentPage;

  // 4. if (currentCount \% pageSize === 0 \&\& currentPage
  \textgreater{} 1) \{ loadPage(currentPage - 1); }

  // 5. if (user.role === `admin') \{ const adminCount =
  document.getElementById(`admin-count'); adminCount.textContent =
  parseInt(adminCount.textContent) - 1; }

  // 6. showToast('\,');

  // 7. if (userId === currentUserId) \{ logout(); } }
  ``\textbackslash code\{

  \important{}

  \important{}DOM

  \important{}``\,``AB
\item
  \important{}

  Web

  }`\code{javascript
  // cart.js - 

  // 
  var cartItems = [];
  var totalPrice = 0;
  var discountRate = 0;
  var shippingFee = 10;

  // 
  function loadCartData() {
    $.ajax({
      url: '/api/cart',
      success: function(data) {
        cartItems = data.items;
        updateCartDisplay();
        calculateTotal();
        updateShippingInfo();
        checkCouponValidity();
      }
    });
  }

  // UI
  function updateCartDisplay() {
    var html = '';
    for (var i = 0; i < cartItems.length; i++) {
      html += '<div class="cart-item" id="item-' + cartItems[i].id + '">';
      html += '<span>' + cartItems[i].name + '</span>';
      html += '<input type="number" value="' + cartItems[i].quantity + '" onchange="updateQuantity(' + cartItems[i].id + ', this.value)">';
      html += '<button onclick="removeItem(' + cartItems[i].id + ')"></button>';
      html += '</div>';
    }
    document.getElementById('cart-list').innerHTML = html;
  }

  // 
  function calculateTotal() {
    totalPrice = 0;
    for (var i = 0; i < cartItems.length; i++) {
      totalPrice += cartItems[i].price \highlight{ cartItems[i].quantity;
    }
    totalPrice = totalPrice } (1 - discountRate) + shippingFee;
    document.getElementById('total-price').textContent = '' + totalPrice.toFixed(2);
  }

  // 
  function updateQuantity(itemId, newQuantity) {
    // 
    for (var i = 0; i < cartItems.length; i++) {
      if (cartItems[i].id === itemId) {
        cartItems[i].quantity = parseInt(newQuantity);
        break;
      }
    }

    // UI
    calculateTotal();
    updateInventoryWarning(itemId);
    saveToLocalStorage();

    // 
    $.post('/api/cart/update', {
      itemId: itemId,
      quantity: newQuantity
    });
  }
  }`\code{

  \important{}API

  \important{}}cartItems\code{}totalPrice\textbackslash code\{

  \important{}

  \important{}DOMAJAX

  \important{}
\item
  \important{}

  Web

  }`\code{javascript
  // 
  function updateUserOnlineStatus(userId, isOnline) {
    // 1. 
    const friendItem = document.querySelector(}#friend-${userId} .status-icon\code{);
    if (friendItem) {
      friendItem.className = isOnline ? 'status-online' : 'status-offline';
    }

    // 2. 
    const chatTitle = document.querySelector(}#chat-${userId} .user-status\code{);
    if (chatTitle) {
      chatTitle.textContent = isOnline ? '' : '';
    }

    // 3. 
    const groupMembers = document.querySelectorAll(}.group-member[data-user="${userId}"]\code{);
    groupMembers.forEach(member => {
      member.setAttribute('data-status', isOnline ? 'online' : 'offline');
    });

    // 4. 
    const onlineCount = document.querySelector('#online-count');
    if (onlineCount) {
      const current = parseInt(onlineCount.textContent);
      onlineCount.textContent = isOnline ? current + 1 : current - 1;
    }

    // 5. 
    const profileStatus = document.querySelector(}#profile-${userId} .status\code{);
    if (profileStatus) {
      profileStatus.textContent = isOnline ? '' : '';
    }
  }
  }`\textbackslash code\{

  \important{}

  \important{}``\,``\,``\,''

  \important{}

  \important{}

  \important{}DOM
\item
  \important{}

  Web

  }`\code{javascript
  // 
  function createEmployeeTable(employees) {
    var html = '<table class="employee-table">';
    html += '<thead><tr>';
    html += '<th onclick="sortEmployees(\'name\')"> <span id="name-sort">↕</span></th>';
    html += '<th onclick="sortEmployees(\'department\')"> <span id="dept-sort">↕</span></th>';
    html += '<th onclick="sortEmployees(\'salary\')"> <span id="salary-sort">↕</span></th>';
    html += '<th></th>';
    html += '</tr></thead><tbody>';

    for (var i = 0; i < employees.length; i++) {
      html += '<tr>';
      html += '<td>' + employees[i].name + '</td>';
      html += '<td>' + employees[i].department + '</td>';
      html += '<td>' + employees[i].salary + '</td>';
      html += '<td><button onclick="editEmployee(' + employees[i].id + ')"></button>';
      html += '<button onclick="deleteEmployee(' + employees[i].id + ')"></button></td>';
      html += '</tr>';
    }
    html += '</tbody></table>';

    // 
    html += '<div class="pagination">';
    for (var page = 1; page <= Math.ceil(employees.length / 10); page++) {
      html += '<button onclick="loadEmployeePage(' + page + ')">' + page + '</button>';
    }
    html += '</div>';

    document.getElementById('employee-container').innerHTML = html;
  }

  //  - 
  function createProjectTable(projects) {
    var html = '<table class="project-table">';
    html += '<thead><tr>';
    html += '<th onclick="sortProjects(\'name\')"> <span id="name-sort">↕</span></th>';
    html += '<th onclick="sortProjects(\'status\')"> <span id="status-sort">↕</span></th>';
    html += '<th onclick="sortProjects(\'deadline\')"> <span id="deadline-sort">↕</span></th>';
    html += '<th></th>';
    html += '</tr></thead><tbody>';

    for (var i = 0; i < projects.length; i++) {
      html += '<tr>';
      html += '<td>' + projects[i].name + '</td>';
      html += '<td>' + projects[i].status + '</td>';
      html += '<td>' + projects[i].deadline + '</td>';
      html += '<td><button onclick="editProject(' + projects[i].id + ')"></button>';
      html += '<button onclick="deleteProject(' + projects[i].id + ')"></button></td>';
      html += '</tr>';
    }
    html += '</tbody></table>';

    // 
    html += '<div class="pagination">';
    for (var page = 1; page <= Math.ceil(projects.length / 10); page++) {
      html += '<button onclick="loadProjectPage(' + page + ')">' + page + '</button>';
    }
    html += '</div>';

    document.getElementById('project-container').innerHTML = html;
  }
  }`\textbackslash code\{

  \important{}HTML

  \important{}

  \important{}Bug

  \important{}

  \important{}

  \important{}

  Web
\end{itemize}

WebDOM

\textbackslash end\{lstlisting}javascript // jQuery function
updateWaterLevel(stationId, newLevel) \{ // \$(`\#chart-' +
stationId).updateChart(newLevel); // \$(`\#table-row-' + stationId + '
.water-level').text(newLevel + `m'); // if (newLevel \textgreater{} 5.0)
\{ \$(`\#warning-' + stationId).addClass(`alert-danger').text('`); } //
\$('\#total-stations').text(calculateActiveStations()); }

\begin{lstlisting}
DOM

### 

\important{}\important{}\important{}

\end{lstlisting}

vue

\begin{verbatim}
<!--  -->
<water-level-chart :data="stationData.level" :station-id="stationId" />
<data-table :rows="tableData" />
<alert-panel :status="warningStatus" :level="stationData.level" />
<statistics-summary :total="totalActiveStations" />
\end{verbatim}

\begin{lstlisting}
### 

ReactAngularVue.js

|| React |Angular| Vue.js |
|---------|-------|---------|---------|
|\important{}|  ||  |
|\important{}|  ||  |
|\important{}|  ||  |
|\important{}|  ||  |
|\important{}|  ||  |
|\important{}|  ||  |
|\important{}|  ||  |

\important{}

\important{1. }
- - HTMLCSSJavaScript
- \important{2. }
- - 
- \important{3. }
- - 
- ### Vue.js

\important{Vue.js}

\important{1. }

Vue.js\important{}

\end{lstlisting}

html

\hypertarget{water-monitor}{}
\{\{ stationName }}

: \{\{ waterLevel }}

: \{\{ flowRate }}/

\textless button @click=``refreshData''\textgreater{}

\begin{lstlisting}
\important{2. }

Vue.jsAPI

\important{3. }

Vue

- \important{Vue Router}: 
- \important{Vuex/Pinia}: 
- \important{Element Plus}: UI
- \important{ECharts}: 
- \important{Vue CLI/Vite}: 

\important{4. }

Vue.js(.vue)

\end{lstlisting}

vue

\begin{verbatim}
<div class="gauge-container">
  <svg class="gauge-svg" viewBox="0 0 200 200">
    <circle 
      cx="100" 
      cy="100" 
      r="80" 
      fill="none" 
      stroke="#e0e6ed" 
      stroke-width="8"
    />
    <circle 
      cx="100" 
      cy="100" 
      r="80" 
      fill="none" 
      :stroke="gaugeColor" 
      stroke-width="8"
      :stroke-dasharray="circumference"
      :stroke-dashoffset="dashOffset"
      transform="rotate(-90 100 100)"
    />
  </svg>
  <div class="gauge-text">
    <div class="water-level">{{ level.toFixed(1) }}m</div>
    <div class="gauge-label"></div>
  </div>
</div>
\end{verbatim}

\begin{lstlisting}
!!! tip ""
    
    Vue.js
    
    1. \important{}
    2. \important{}
    3. \important{}
    4. \important{}
    5. \important{}

Vue.jsMVVM

## 4.5.2 Vue.jsMVVM

### Vue.js

\important{Vue.js}"view"\important{JavaScript}2014Vue.jsAPIVue.js

\important{Vue.js}

1. \important{}
2. \important{}
3. \important{}
4. \important{}

Vue.js

\end{lstlisting}

html \textless!DOCTYPE html\textgreater{}

\begin{verbatim}
    <div class="monitoring-dashboard">
        <h1>{{ systemTitle }}</h1>
        
        <!--  -->
        <div class="station-card" 
             :class="getStationStatus(station)" 
             v-for="station in monitoringStations" 
             :key="station.id">
            <h3>{{ station.name }}</h3>
            
            <div class="data-row">
                <span>:</span>
                <strong>{{ station.waterLevel.toFixed(2) }} </strong>
            </div>
            
            <div class="data-row">
                <span>:</span>
                <strong>{{ station.flowRate.toFixed(1) }} /</strong>
            </div>
            
            <div class="data-row">
                <span>:</span>
                <span :style="{ color: station.waterLevel > 5 ? 'red' : 'green' }">
                    {{ station.waterLevel > 5 ? '' : '' }}
                </span>
            </div>
            
            <div class="data-row">
                <span>:</span>
                <span>{{ formatTime(station.lastUpdate) }}</span>
            </div>
            
            <button class="btn" @click="refreshStationData(station.id)">
                
            </button>
        </div>
        
        <!--  -->
        <div class="station-card">
            <h3></h3>
            <div class="data-row">
                <span>:</span>
                <strong>{{ totalStations }}</strong>
            </div>
            <div class="data-row">
                <span>:</span>
                <strong>{{ normalStations }}</strong>
            </div>
            <div class="data-row">
                <span>:</span>
                <strong>{{ warningStations }}</strong>
            </div>
        </div>
    </div>
</div>

<script>
    const { createApp, ref, computed } = Vue
    
    createApp({
        setup() {
            // 
            const systemTitle = ref('')
            const monitoringStations = ref([
                {
                    id: 1,
                    name: '',
                    waterLevel: 4.25,
                    flowRate: 1580.5,
                    lastUpdate: new Date()
                },
                {
                    id: 2,
                    name: '',
                    waterLevel: 6.80,
                    flowRate: 2450.3,
                    lastUpdate: new Date()
                },
                {
                    id: 3,
                    name: '',
                    waterLevel: 3.15,
                    flowRate: 890.7,
                    lastUpdate: new Date()
                }
            ])
            
            //  - 
            const totalStations = computed(() => {
                return monitoringStations.value.length
            })
            
            const warningStations = computed(() => {
                return monitoringStations.value.filter(station => station.waterLevel > 5).length
            })
            
            const normalStations = computed(() => {
                return totalStations.value - warningStations.value
            })
            
            // 
            const refreshStationData = (stationId) => {
                const station = monitoringStations.value.find(s => s.id === stationId)
                if (station) {
                    // 
                    station.waterLevel = (Math.random() } 4 + 2).toFixed(2) \highlight{ 1
                    station.flowRate = (Math.random() } 1000 + 500).toFixed(1) \highlight{ 1
                    station.lastUpdate = new Date()
                }
            }
            
            const getStationStatus = (station) => {
                return station.waterLevel > 5 ? 'warning' : 'normal'
            }
            
            const formatTime = (date) => {
                return date.toLocaleTimeString('zh-CN')
            }
            
            return {
                systemTitle,
                monitoringStations,
                totalStations,
                normalStations,
                warningStations,
                refreshStationData,
                getStationStatus,
                formatTime
            }
        }
    }).mount('#water-monitoring-app')
</script>
\end{verbatim}

\begin{lstlisting}
Vue.js

1. \important{}}monitoringStations\code{
2. \important{}}totalStations\code{}warningStations\code{
3. \important{}}v-for\code{
4. \important{}}@click\code{

### Vue.js

\important{1. }

Vue.js\important{}Vue

\end{lstlisting}

vue

\begin{verbatim}
<h2></h2>
<div v-if="isLoading" class="loading">...</div>
<div v-else>
  <div class="quality-item" v-for="item in waterQualityData" :key="item.parameter">
    <span class="parameter">{{ item.parameter }}:</span>
    <span 
      class="value" 
      :class="{ 'exceeded': item.value > item.standard }"
    >
      {{ item.value }} {{ item.unit }}
    </span>
    <span class="standard">(: {{ item.standard }}{{ item.unit }})</span>
  </div>
</div>
\end{verbatim}

\begin{lstlisting}
\important{2. }

Vue.js\important{}DOM

\end{lstlisting}

vue

\begin{verbatim}
<h2>{{ reservoirName }}</h2>

<!--  -->
<div class="data-grid">
  <div class="data-card">
    <h3></h3>
    <div class="value">{{ currentLevel.toFixed(2) }} </div>
    <div class="trend" :class="levelTrend">
      {{ levelTrend === 'rising' ? '↗ ' : levelTrend === 'falling' ? '↘ ' : '→ ' }}
    </div>
  </div>
  
  <div class="data-card">
    <h3></h3>
    <div class="value">{{ currentVolume }} </div>
    <div class="percentage">{{ volumePercentage }}% </div>
  </div>
  
  <div class="data-card">
    <h3></h3>
    <div class="value">{{ inflowRate }} /</div>
  </div>
  
  <div class="data-card">
    <h3></h3>
    <div class="value">{{ outflowRate }} /</div>
  </div>
</div>

<!--  -->
<div class="control-panel">
  <h3></h3>
  <div class="controls">
    <button @click="simulateRainfall" :disabled="isSimulating"></button>
    <button @click="adjustOutflow" :disabled="isSimulating"></button>
    <button @click="resetData"></button>
  </div>
  <div v-if="isSimulating" class="simulation-info">
    ... {{ simulationProgress }}%
  </div>
</div>

<!--  -->
<div v-if="warnings.length > 0" class="warnings">
  <h3></h3>
  <div class="warning-item" v-for="warning in warnings" :key="warning.id">
    <span class="warning-level" :class="warning.level">{{ warning.level }}</span>
    <span class="warning-message">{{ warning.message }}</span>
    <span class="warning-time">{{ formatTime(warning.time) }}</span>
  </div>
</div>
\end{verbatim}

\begin{lstlisting}
Vue.js

- \important{}}currentLevel\code{
- \important{}}currentVolume\code{}volumePercentage\code{}currentLevel\code{
- \important{}}watch\code{

### MVVM

\important{MVVMModel-View-ViewModel}Vue.js

||  |Vue|  |
|------|------|--------------|-------------|
|\important{Model}|  |API|  |
|\important{View}|  |Template|  |
|\important{ViewModel}| ModelView |Vue|  |

MVVM

\end{lstlisting}

vue

\begin{verbatim}
<header class="manager-header">
  <h1></h1>
  <div class="header-actions">
    <button @click="refreshAllData" :disabled="isLoading" class="btn btn-primary">
      {{ isLoading ? '...' : '' }}
    </button>
    <button @click="showAddDialog = true" class="btn btn-success"></button>
  </div>
</header>

<!--  -->
<div class="filter-section">
  <input 
    v-model="searchKeyword" 
    placeholder="..."
    class="search-input"
  />
  <select v-model="filterStatus" class="filter-select">
    <option value=""></option>
    <option value="normal"></option>
    <option value="warning"></option>
    <option value="offline"></option>
  </select>
</div>

<!--  -->
<div class="stations-grid">
  <div 
    v-for="station in filteredStations" 
    :key="station.id"
    class="station-card"
    :class="getStationStatusClass(station)"
  >
    <div class="station-header">
      <h3>{{ station.name }}</h3>
      <span class="status-badge" :class="station.status">
        {{ getStatusText(station.status) }}
      </span>
    </div>
    
    <div class="station-data">
      <div class="data-item">
        <span class="label">:</span>
        <span class="value">{{ station.longitude }}°</span>
      </div>
      <div class="data-item">
        <span class="label">:</span>
        <span class="value">{{ station.latitude }}°</span>
      </div>
      <div class="data-item">
        <span class="label">:</span>
        <span class="value">{{ station.currentData.waterLevel }} m</span>
      </div>
      <div class="data-item">
        <span class="label">:</span>
        <span class="value">{{ station.currentData.flowRate }} m³/s</span>
      </div>
      <div class="data-item">
        <span class="label">:</span>
        <span class="value">{{ formatTime(station.lastUpdate) }}</span>
      </div>
    </div>
    
    <div class="station-actions">
      <button @click="viewDetails(station)" class="btn btn-info btn-sm"></button>
      <button @click="editStation(station)" class="btn btn-warning btn-sm"></button>
      <button @click="deleteStation(station.id)" class="btn btn-danger btn-sm"></button>
    </div>
  </div>
</div>

<!-- / -->
<div v-if="showAddDialog || editingStation" class="modal-overlay" @click.self="closeDialog">
  <div class="modal-content">
    <div class="modal-header">
      <h2>{{ editingStation ? '' : '' }}</h2>
      <button @click="closeDialog" class="close-btn">&times;</button>
    </div>
    
    <form @submit.prevent="saveStation" class="station-form">
      <div class="form-group">
        <label>:</label>
        <input 
          v-model="stationForm.name" 
          type="text" 
          required 
          class="form-input"
        />
      </div>
      
      <div class="form-row">
        <div class="form-group">
          <label>:</label>
          <input 
            v-model.number="stationForm.longitude" 
            type="number" 
            step="0.000001"
            required 
            class="form-input"
          />
        </div>
        <div class="form-group">
          <label>:</label>
          <input 
            v-model.number="stationForm.latitude" 
            type="number" 
            step="0.000001"
            required 
            class="form-input"
          />
        </div>
      </div>
      
      <div class="form-group">
        <label>:</label>
        <textarea 
          v-model="stationForm.description" 
          class="form-textarea"
          rows="3"
        ></textarea>
      </div>
      
      <div class="form-actions">
        <button type="submit" class="btn btn-primary" :disabled="isSubmitting">
          {{ isSubmitting ? '...' : '' }}
        </button>
        <button type="button" @click="closeDialog" class="btn btn-secondary"></button>
      </div>
    </form>
  </div>
</div>
\end{verbatim}

\begin{lstlisting}
Vue.jsMVVM

\important{Model}
- }stations\code{: 
- }stationForm\code{: 
- }searchKeyword\code{, }filterStatus\code{: 

\important{View}
- Template
- (}v-for\code{, }v-if\code{, }v-model\code{)
- (}@click\code{, }@submit\code{)

\important{ViewModel}
- (}filteredStations\code{, }stationStats\code{)
- (}watch\code{)
- (}refreshAllData\code{, }saveStation\code{)

### 

Vue.js\important{}MVVM(Model)(View)

\important{}
1. \important{} → \important{} → \important{}

\important{}
1. \important{} → \important{}
2. \important{} → \important{} → \important{}

\end{lstlisting}

vue

\begin{verbatim}
<h2></h2>

<!--  -->
<div class="input-section">
  <div class="input-group">
    <label> ():</label>
    <input 
      v-model.number="waterLevel" 
      type="number" 
      step="0.1"
      min="0"
      max="20"
      class="level-input"
    />
  </div>
  
  <div class="input-group">
    <label>:</label>
    <select v-model="selectedStation" class="station-select">
      <option value=""></option>
      <option v-for="station in availableStations" :key="station.id" :value="station">
        {{ station.name }}
      </option>
    </select>
  </div>
  
  <div class="input-group">
    <label>:</label>
    <textarea 
      v-model="notes" 
      placeholder="..."
      class="notes-textarea"
    ></textarea>
  </div>
</div>

<!--  -->
<div class="preview-section">
  <h3></h3>
  <div class="preview-card">
    <div class="preview-item">
      <strong>:</strong> 
      <span :class="getLevelClass(waterLevel)">{{ waterLevel || 0 }} </span>
    </div>
    <div class="preview-item">
      <strong>:</strong> 
      <span>{{ selectedStation ? selectedStation.name : '' }}</span>
    </div>
    <div class="preview-item" v-if="selectedStation">
      <strong>:</strong> 
      <span>{{ selectedStation.longitude }}°E, {{ selectedStation.latitude }}°N</span>
    </div>
    <div class="preview-item">
      <strong>:</strong> 
      <span :class="getWarningClass(waterLevel)">{{ getWarningText(waterLevel) }}</span>
    </div>
    <div class="preview-item" v-if="notes">
      <strong>:</strong> 
      <span>{{ notes }}</span>
    </div>
    <div class="preview-item">
      <strong>:</strong> 
      <span>{{ currentTime }}</span>
    </div>
  </div>
</div>

<!--  -->
<div class="sync-demo">
  <h3></h3>
  <p></p>
  <div class="sync-display">
    <div class="gauge-wrapper">
      <div class="level-gauge">
        <div 
          class="gauge-fill" 
          :style="{ height: (waterLevel / 20 \highlight{ 100) + '%' }"
        ></div>
        <div class="gauge-text">{{ waterLevel || 0 }}m</div>
      </div>
      <div class="gauge-label"></div>
    </div>
    
    <div class="chart-wrapper">
      <div class="mini-chart">
        <div 
          v-for="(point, index) in chartData" 
          :key="index"
          class="chart-bar"
          :style="{ 
            height: (point / 20 } 100) + '%',
            background: point > 15 ? '#ff4757' : point > 10 ? '#ffa502' : '#2ed573'
          }"
        ></div>
      </div>
      <div class="chart-label"></div>
    </div>
  </div>
</div>
\end{verbatim}

\begin{lstlisting}
Vue.js

1. \important{}
2. \important{}
3. \important{}
4. \important{}

### DOM

\important{DOMVirtual DOM}Vue.jsDOMJavaScriptVueDOMDOM

\important{DOM}

\end{lstlisting}

javascript // 1. DOM const initialVNode = \{ tag: `div', props: \{
class: `water-station' }, children: {[} \{ tag: `h3', props: \{},
children: `A' }, \{ tag: `span', props: \{}, children: `: 4.5m' } {]}
}

// 2. DOM const updatedVNode = \{ tag: `div', props: \{ class:
`water-station' }, children: {[} \{ tag: `h3', props: \{}, children:
`A' // }, \{ tag: `span', props: \{}, children: `: 5.2m' // } {]} }

// 3. Vuediff //

\begin{lstlisting}
\important{DOM}

1. \important{}
2. \important{}DOM
3. \important{}DOM

Vue.js

## 4.5.3 Vue

### Vue

Vue 3Vue 2}createApp\code{

\important{Vue 3}

1. \important{Vue}
2. \important{}
3. \important{}
4. \important{DOM}

Vue

\end{lstlisting}

html \textless!DOCTYPE html\textgreater{}

\begin{verbatim}
    <div class="monitoring-console">
        <!--  -->
        <header class="console-header">
            <h1>{{ systemInfo.name }}</h1>
            <p>{{ systemInfo.description }}</p>
            <p>: {{ formatDate(systemInfo.startTime) }} | : {{ uptime }}</p>
        </header>
        
        <!--  -->
        <div class="status-grid">
            <div class="status-card" :class="getCardClass('stations')">
                <div class="metric-label"></div>
                <div class="metric-value">{{ metrics.onlineStations }} / {{ metrics.totalStations }}</div>
                <div class="metric-label">
                    : {{ stationOnlineRate }}%
                </div>
            </div>
            
            <div class="status-card" :class="getCardClass('dataFlow')">
                <div class="metric-label"></div>
                <div class="metric-value">{{ metrics.dataFlow }} MB/h</div>
                <div class="metric-label">
                    {{ metrics.dataFlow > 100 ? '' : '' }}
                </div>
            </div>
            
            <div class="status-card" :class="getCardClass('alerts')">
                <div class="metric-label"></div>
                <div class="metric-value">{{ metrics.activeAlerts }}</div>
                <div class="metric-label">
                    : {{ formatTime(metrics.lastAlertTime) }}
                </div>
            </div>
            
            <div class="status-card" :class="getCardClass('storage')">
                <div class="metric-label"></div>
                <div class="metric-value">{{ metrics.storageUsage }}%</div>
                <div class="metric-label">
                    : {{ 100 - metrics.storageUsage }}%
                </div>
            </div>
        </div>
        
        <!--  -->
        <div class="control-panel">
            <h3></h3>
            <div class="btn-group">
                <button 
                    class="btn btn-primary" 
                    @click="refreshSystemData"
                    :disabled="isLoading"
                >
                    {{ isLoading ? '...' : '' }}
                </button>
                
                <button 
                    class="btn btn-success" 
                    @click="startDataCollection"
                    :disabled="dataCollectionActive"
                >
                    {{ dataCollectionActive ? '...' : '' }}
                </button>
                
                <button 
                    class="btn btn-warning" 
                    @click="pauseDataCollection"
                    :disabled="!dataCollectionActive"
                >
                    
                </button>
                
                <button 
                    class="btn btn-danger" 
                    @click="clearAllAlerts"
                    :disabled="metrics.activeAlerts === 0"
                >
                    
                </button>
            </div>
        </div>
        
        <!--  -->
        <div class="log-section">
            <h4 style="margin-top: 0;"> ({{ systemLogs.length }})</h4>
            <div v-for="log in systemLogs" :key="log.id" class="log-entry">
                <span class="log-timestamp">[{{ formatTime(log.timestamp) }}]</span>
                <span :class="'log-' + log.level">{{ log.message }}</span>
            </div>
        </div>
    </div>
</div>

<script>
    // Vue 3API
    const { createApp, ref, reactive, computed, watch, onMounted, onUnmounted } = Vue
    
    // Vue
    const app = createApp({
        // setupComposition API
        setup() {
            // =====  =====
            
            // 
            const systemInfo = reactive({
                name: '',
                description: '',
                startTime: new Date(),
                version: '2.0.1'
            })
            
            // 
            const metrics = reactive({
                totalStations: 156,
                onlineStations: 142,
                dataFlow: 85.6,
                activeAlerts: 3,
                storageUsage: 67,
                lastAlertTime: new Date()
            })
            
            // 
            const isLoading = ref(false)
            const dataCollectionActive = ref(true)
            const uptime = ref('00:00:00')
            
            // 
            const systemLogs = ref([
                {
                    id: 1,
                    timestamp: new Date(),
                    level: 'info',
                    message: ''
                },
                {
                    id: 2,
                    timestamp: new Date(Date.now() - 30000),
                    level: 'warning',
                    message: ' #045 '
                },
                {
                    id: 3,
                    timestamp: new Date(Date.now() - 60000),
                    level: 'success',
                    message: ''
                }
            ])
            
            // =====  =====
            
            // 
            const stationOnlineRate = computed(() => {
                if (metrics.totalStations === 0) return 0
                return Math.round((metrics.onlineStations / metrics.totalStations) \highlight{ 100)
            })
            
            // 
            const systemStatus = computed(() => {
                const onlineRate = stationOnlineRate.value
                const alertCount = metrics.activeAlerts
                
                if (onlineRate < 80 || alertCount > 5) return 'error'
                if (onlineRate < 95 || alertCount > 2) return 'warning'
                return 'success'
            })
            
            // =====  =====
            
            // 
            watch(() => metrics.activeAlerts, (newCount, oldCount) => {
                if (newCount > oldCount) {
                    addLog('warning', } ${newCount} \code{)
                } else if (newCount < oldCount) {
                    addLog('success', } ${newCount} \code{)
                }
            })
            
            // 
            watch(() => metrics.onlineStations, (newCount) => {
                const rate = Math.round((newCount / metrics.totalStations) } 100)
                if (rate < 80) {
                    addLog('error', } ${rate}%\code{)
                } else if (rate >= 95) {
                    addLog('success', } ${rate}%\code{)
                }
            })
            
            // =====  =====
            
            // 
            const refreshSystemData = async () => {
                isLoading.value = true
                addLog('info', '...')
                
                // API
                await new Promise(resolve => setTimeout(resolve, 1500))
                
                try {
                    // 
                    metrics.onlineStations = Math.floor(Math.random() \highlight{ 20) + 140
                    metrics.dataFlow = Math.round((Math.random() } 50 + 60) \highlight{ 10) / 10
                    metrics.activeAlerts = Math.floor(Math.random() } 5)
                    metrics.storageUsage = Math.floor(Math.random() \highlight{ 30) + 50
                    metrics.lastAlertTime = new Date()
                    
                    addLog('success', '')
                } catch (error) {
                    addLog('error', }: ${error.message}\code{)
                } finally {
                    isLoading.value = false
                }
            }
            
            // 
            const startDataCollection = () => {
                dataCollectionActive.value = true
                addLog('success', '')
                
                // 
                const collectionInterval = setInterval(() => {
                    if (!dataCollectionActive.value) {
                        clearInterval(collectionInterval)
                        return
                    }
                    
                    // 
                    metrics.dataFlow = Math.round((metrics.dataFlow + (Math.random() - 0.5) } 10) \highlight{ 10) / 10
                    metrics.dataFlow = Math.max(0, Math.min(200, metrics.dataFlow))
                }, 3000)
            }
            
            // 
            const pauseDataCollection = () => {
                dataCollectionActive.value = false
                addLog('warning', '')
            }
            
            // 
            const clearAllAlerts = () => {
                const clearedCount = metrics.activeAlerts
                metrics.activeAlerts = 0
                addLog('info', } ${clearedCount} \code{)
            }
            
            // 
            const addLog = (level, message) => {
                const newLog = {
                    id: Date.now(),
                    timestamp: new Date(),
                    level,
                    message
                }
                
                systemLogs.value.unshift(newLog)
                
                // 50
                if (systemLogs.value.length > 50) {
                    systemLogs.value = systemLogs.value.slice(0, 50)
                }
            }
            
            // 
            const getCardClass = (type) => {
                switch (type) {
                    case 'stations':
                        return stationOnlineRate.value < 80 ? 'error' : 
                               stationOnlineRate.value < 95 ? 'warning' : 'success'
                    case 'dataFlow':
                        return metrics.dataFlow > 150 ? 'error' : 
                               metrics.dataFlow > 100 ? 'warning' : 'success'
                    case 'alerts':
                        return metrics.activeAlerts > 5 ? 'error' : 
                               metrics.activeAlerts > 2 ? 'warning' : 'success'
                    case 'storage':
                        return metrics.storageUsage > 90 ? 'error' : 
                               metrics.storageUsage > 80 ? 'warning' : 'success'
                    default:
                        return ''
                }
            }
            
            // 
            const formatDate = (date) => {
                return date.toLocaleDateString('zh-CN') + ' ' + date.toLocaleTimeString('zh-CN')
            }
            
            // 
            const formatTime = (date) => {
                return date.toLocaleTimeString('zh-CN')
            }
            
            // 
            const updateUptime = () => {
                const now = new Date()
                const diff = now - systemInfo.startTime
                const hours = Math.floor(diff / (1000 } 60 \highlight{ 60))
                const minutes = Math.floor((diff % (1000 } 60 \highlight{ 60)) / (1000 } 60))
                const seconds = Math.floor((diff % (1000 \highlight{ 60)) / 1000)
                
                uptime.value = }${hours.toString().padStart(2, '0')}:${minutes.toString().padStart(2, '0')}:${seconds.toString().padStart(2, '0')}\code{
            }
            
            // =====  =====
            
            let uptimeTimer = null
            
            onMounted(() => {
                // 
                addLog('info', '')
                
                // 
                uptimeTimer = setInterval(updateUptime, 1000)
                
                // 
                if (dataCollectionActive.value) {
                    startDataCollection()
                }
            })
            
            onUnmounted(() => {
                // 
                if (uptimeTimer) {
                    clearInterval(uptimeTimer)
                }
                addLog('info', '')
            })
            
            // =====  =====
            return {
                // 
                systemInfo,
                metrics,
                isLoading,
                dataCollectionActive,
                uptime,
                systemLogs,
                
                // 
                stationOnlineRate,
                systemStatus,
                
                // 
                refreshSystemData,
                startDataCollection,
                pauseDataCollection,
                clearAllAlerts,
                getCardClass,
                formatDate,
                formatTime
            }
        }
    })
    
    // 
    app.config.globalProperties.$version = '2.0.1'
    app.config.errorHandler = (err, instance, info) => {
        console.error('Vue:', err)
        console.error(':', instance)
        console.error(':', info)
    }
    
    // DOM
    app.mount('#water-monitoring-console')
</script>
\end{verbatim}

\begin{lstlisting}
Vue 3

\important{1. }
- }createApp()\code{
- }setup()\code{Composition API
- }mount()\code{DOM

\important{2. }
- }ref()\code{
- }reactive()\code{
- \important{3. }
- }computed()\code{
- }watch()\code{

\important{4. }
- }onMounted()\code{
- }onUnmounted()\code{

### 

Vue 3Composition API

\important{}

\end{lstlisting}

vue

\begin{verbatim}
<h2></h2>

<!--  -->
<div class="data-section">
  <h3></h3>
  <div class="sensor-grid">
    <div v-for="sensor in sensorList" :key="sensor.id" class="sensor-card">
      <h4>{{ sensor.name }}</h4>
      <div class="sensor-value">
        <span class="value">{{ sensor.currentValue }}</span>
        <span class="unit">{{ sensor.unit }}</span>
      </div>
      <div class="sensor-status" :class="getSensorStatus(sensor)">
        {{ getSensorStatusText(sensor) }}
      </div>
      <div class="last-update">
        : {{ formatTime(sensor.lastUpdate) }}
      </div>
    </div>
  </div>
</div>

<!--  -->
<div class="config-section">
  <h3></h3>
  <form @submit.prevent="saveConfiguration">
    <div class="config-group">
      <label> ():</label>
      <input 
        v-model.number="config.samplingInterval" 
        type="number" 
        min="1" 
        max="3600"
        class="config-input"
      />
    </div>
    
    <div class="config-group">
      <label> ():</label>
      <input 
        v-model.number="config.dataRetentionDays" 
        type="number" 
        min="1" 
        max="3650"
        class="config-input"
      />
    </div>
    
    <div class="config-group">
      <label>:</label>
      <div class="threshold-settings">
        <div v-for="(threshold, key) in config.alertThresholds" :key="key">
          <span>{{ getThresholdLabel(key) }}:</span>
          <input 
            v-model.number="threshold.min" 
            type="number" 
            step="0.1"
            placeholder=""
            class="threshold-input"
          />
          <span> ~ </span>
          <input 
            v-model.number="threshold.max" 
            type="number" 
            step="0.1"
            placeholder=""
            class="threshold-input"
          />
          <span>{{ getThresholdUnit(key) }}</span>
        </div>
      </div>
    </div>
    
    <div class="config-group">
      <label>:</label>
      <div class="feature-toggles">
        <label class="toggle-item">
          <input 
            v-model="config.features.autoAlert" 
            type="checkbox"
          />
          
        </label>
        <label class="toggle-item">
          <input 
            v-model="config.features.dataBackup" 
            type="checkbox"
          />
          
        </label>
        <label class="toggle-item">
          <input 
            v-model="config.features.remoteMonitoring" 
            type="checkbox"
          />
          
        </label>
      </div>
    </div>
    
    <button type="submit" class="btn btn-primary" :disabled="isSaving">
      {{ isSaving ? '...' : '' }}
    </button>
  </form>
</div>

<!--  -->
<div class="stats-section">
  <h3></h3>
  <div class="stats-grid">
    <div class="stat-card">
      <div class="stat-label"></div>
      <div class="stat-value">{{ totalSensors }}</div>
    </div>
    <div class="stat-card">
      <div class="stat-label"></div>
      <div class="stat-value">{{ onlineSensors }}</div>
    </div>
    <div class="stat-card">
      <div class="stat-label"></div>
      <div class="stat-value">{{ abnormalSensors }}</div>
    </div>
    <div class="stat-card">
      <div class="stat-label"></div>
      <div class="stat-value">{{ dataIntegrityRate }}%</div>
    </div>
  </div>
</div>
\end{verbatim}

\begin{lstlisting}
Vue 3

\important{1. }
- }ref()\code{: 
- }reactive()\code{: 

\important{2. }
- }watch()\code{: 
- (}{ deep: true }\code{): 

\important{3. }
- - 

### 

Vue.jsHTML

\important{Vue}

1. \important{} - 
2. \important{} - HTML
3. \important{} - 
4. \important{} - 
5. \important{} - 

\end{lstlisting}

vue

\begin{verbatim}
<!--  -  -->
<header class="dashboard-header">
  <h1>{{ dashboardTitle }}</h1>
  <div class="header-info">
    <span class="location">{{ currentLocation }}</span>
    <span class="datetime">{{ formatDateTime(currentTime) }}</span>
    <span class="weather" :class="weatherClass">
      {{ weatherInfo.description }} {{ weatherInfo.temperature }}°C
    </span>
  </div>
</header>

<!--  -  -->
<div class="system-status">
  <div 
    class="status-indicator" 
    :class="systemStatusClass"
    :title="systemStatusTooltip"
  >
    <!-- v-if/v-else-if/v-else  -->
    <span v-if="systemStatus === 'healthy'" class="status-icon"></span>
    <span v-else-if="systemStatus === 'warning'" class="status-icon"></span>
    <span v-else-if="systemStatus === 'error'" class="status-icon"></span>
    <span v-else class="status-icon">?</span>
    
    <span class="status-text">{{ systemStatusText }}</span>
  </div>
  
  <!-- v-show  -->
  <div v-show="showDetailedStatus" class="detailed-status">
    <p>: {{ onlineStations }} / {{ totalStations }}</p>
    <p>: {{ dataIntegrityRate }}%</p>
    <p>: {{ formatTime(lastUpdateTime) }}</p>
  </div>
</div>

<!--  -  -->
<div class="monitoring-grid">
  <h2></h2>
  
  <!-- v-for  -->
  <div class="grid-container">
    <div 
      v-for="station in filteredStations" 
      :key="station.id"
      class="station-card"
      :class="getStationCardClass(station)"
      @click="selectStation(station)"
    >
      <!--  -->
      <div class="station-header">
        <h3>{{ station.name }}</h3>
        <span 
          class="station-id"
          :style="{ color: station.status === 'online' ? '#52c41a' : '#ff4757' }"
        >
          #{{ station.id }}
        </span>
      </div>
      
      <!--  -->
      <div class="station-data">
        <div 
          v-for="(value, key) in station.measurements" 
          :key="key"
          class="data-item"
        >
          <span class="data-label">{{ getMeasurementLabel(key) }}:</span>
          <span 
            class="data-value" 
            :class="getValueStatusClass(key, value)"
            :title="}: ${getValueRange(key)}\code{"
          >
            {{ formatValue(key, value) }}
          </span>
        </div>
      </div>
      
      <!--  -->
      <div v-if="station.alerts && station.alerts.length > 0" class="station-alerts">
        <div 
          v-for="alert in station.alerts" 
          :key="alert.id"
          class="alert-item"
          :class="}alert-${alert.level}\code{"
        >
          <span class="alert-icon">{{ getAlertIcon(alert.level) }}</span>
          <span class="alert-message">{{ alert.message }}</span>
        </div>
      </div>
      
      <!--  -->
      <div class="station-actions">
        <button 
          class="btn btn-sm"
          :class="{ 'btn-primary': !station.isMonitoring, 'btn-warning': station.isMonitoring }"
          @click.stop="toggleMonitoring(station)"
          :disabled="station.status === 'offline'"
        >
          {{ station.isMonitoring ? '' : '' }}
        </button>
        
        <button 
          class="btn btn-sm btn-info"
          @click.stop="viewHistory(station)"
        >
          
        </button>
      </div>
    </div>
  </div>
  
  <!--  -->
  <div v-if="filteredStations.length === 0" class="empty-state">
    <div class="empty-icon"></div>
    <h3></h3>
    <p v-if="searchKeyword">
      "{{ searchKeyword }}"
    </p>
    <p v-else>
      
    </p>
    <button class="btn btn-primary" @click="refreshData"></button>
  </div>
</div>

<!--  -->
<div class="filter-controls">
  <h3></h3>
  
  <!-- v-model  -->
  <div class="filter-group">
    <label>:</label>
    <input 
      v-model="searchKeyword" 
      type="text"
      placeholder="ID..."
      class="search-input"
      @input="handleSearchInput"
    />
  </div>
  
  <div class="filter-group">
    <label>:</label>
    <select v-model="selectedStatus" class="status-filter">
      <option value=""></option>
      <option value="online"></option>
      <option value="offline"></option>
      <option value="warning"></option>
    </select>
  </div>
  
  <div class="filter-group">
    <label>:</label>
    <div class="region-checkboxes">
      <label 
        v-for="region in availableRegions" 
        :key="region"
        class="checkbox-label"
      >
        <input 
          type="checkbox" 
          :value="region"
          v-model="selectedRegions"
          @change="handleRegionChange"
        />
        {{ region }}
      </label>
    </div>
  </div>
  
  <div class="filter-group">
    <label>:</label>
    <div class="data-type-radios">
      <label 
        v-for="dataType in dataTypes" 
        :key="dataType.value"
        class="radio-label"
      >
        <input 
          type="radio" 
          :value="dataType.value"
          v-model="selectedDataType"
        />
        {{ dataType.label }}
      </label>
    </div>
  </div>
</div>

<!--  -  -->
<div class="data-table-section" v-if="showDataTable">
  <h3></h3>
  <div class="table-wrapper">
    <table class="data-table">
      <thead>
        <tr>
          <th 
            v-for="column in tableColumns" 
            :key="column.key"
            :class="{ 'sortable': column.sortable }"
            @click="column.sortable && sortBy(column.key)"
          >
            {{ column.title }}
            <span 
              v-if="column.sortable"
              class="sort-indicator"
              :class="getSortClass(column.key)"
            >
              {{ getSortIcon(column.key) }}
            </span>
          </th>
        </tr>
      </thead>
      <tbody>
        <!-- v-for -->
        <template v-for="station in paginatedTableData" :key="station.id">
          <tr 
            class="station-row"
            :class="{ 'selected': selectedTableRows.includes(station.id) }"
            @click="toggleTableRowSelection(station.id)"
          >
            <td>{{ station.name }}</td>
            <td>{{ station.id }}</td>
            <td>
              <span 
                class="status-badge"
                :class="}status-${station.status}\code{"
              >
                {{ station.status }}
              </span>
            </td>
            <td>{{ station.location }}</td>
            <td>{{ formatTime(station.lastUpdate) }}</td>
            <td>
              <div class="table-actions">
                <button 
                  class="btn-icon" 
                  @click.stop="editStation(station)"
                  title=""
                >
                  
                </button>
                <button 
                  class="btn-icon" 
                  @click.stop="deleteStation(station.id)"
                  title=""
                >
                  
                </button>
              </div>
            </td>
          </tr>
          
          <!--  -  -->
          <tr 
            v-if="expandedRows.includes(station.id)"
            class="expanded-row"
          >
            <td :colspan="tableColumns.length">
              <div class="expanded-content">
                <div class="measurement-details">
                  <h4></h4>
                  <div class="measurement-grid">
                    <div 
                      v-for="(value, key) in station.measurements"
                      :key="key"
                      class="measurement-item"
                    >
                      <span class="measurement-name">{{ getMeasurementLabel(key) }}</span>
                      <span class="measurement-value">{{ formatValue(key, value) }}</span>
                      <div 
                        class="measurement-trend"
                        :class="getTrendClass(station.trends[key])"
                      >
                        {{ getTrendIcon(station.trends[key]) }} 
                        {{ station.trends[key] }}%
                      </div>
                    </div>
                  </div>
                </div>
              </div>
            </td>
          </tr>
        </template>
      </tbody>
    </table>
    
    <!--  -  -->
    <div class="pagination" v-if="totalPages > 1">
      <button 
        class="page-btn"
        :disabled="currentPage === 1"
        @click="changePage(currentPage - 1)"
      >
        
      </button>
      
      <span 
        v-for="page in visiblePageNumbers" 
        :key="page"
        class="page-number"
        :class="{ 'active': page === currentPage }"
        @click="changePage(page)"
      >
        {{ page }}
      </span>
      
      <button 
        class="page-btn"
        :disabled="currentPage === totalPages"
        @click="changePage(currentPage + 1)"
      >
        
      </button>
    </div>
  </div>
</div>

<!--  -  -->
<div class="charts-section">
  <h3></h3>
  <div class="chart-container">
    <!--  -->
    <div 
      v-for="chartData in chartDataSets" 
      :key="chartData.id"
      class="chart-item"
      :style="getChartStyle(chartData)"
    >
      <h4>{{ chartData.title }}</h4>
      <div class="chart-visualization">
        <!--  -->
        <div class="chart-bars">
          <div 
            v-for="(dataPoint, index) in chartData.data.slice(-10)"
            :key="index"
            class="chart-bar"
            :style="{ 
              height: (dataPoint.value / chartData.maxValue \highlight{ 100) + '%',
              backgroundColor: getBarColor(dataPoint.value, chartData)
            }"
            :title="}${dataPoint.label}: ${dataPoint.value}${chartData.unit}\code{"
          ></div>
        </div>
        
        <div class="chart-info">
          <span class="current-value">
            : {{ chartData.currentValue }}{{ chartData.unit }}
          </span>
          <span 
            class="trend-indicator"
            :class="chartData.trend"
          >
            {{ getTrendText(chartData.trend) }}
          </span>
        </div>
      </div>
    </div>
  </div>
</div>
\end{verbatim}

\begin{lstlisting}
Vue.js

\important{1. }
- }{{ dashboardTitle }}\code{ - 
- }{{ formatDateTime(currentTime) }}\code{ - 

\important{2. }
- }:class="systemStatusClass"\code{ - class
- }:style="getChartStyle(chartData)"\code{ - 
- }:disabled="station.status === 'offline'"\code{ - 

\important{3. }
- }v-if/v-else-if/v-else\code{ - 
- }v-show\code{ - DOM

\important{4. }
- }v-for="station in filteredStations"\code{ - 
- }v-for="(value, key) in station.measurements"\code{ - 
- \important{5. }
- }@click="selectStation(station)"\code{ - 
- }@click.stop\code{ - 
- }@submit.prevent\code{ - 

\important{6. }
- }v-model="searchKeyword"\code{ - 
- }v-model="selectedRegions"\code{ - 
- }v-model.number\code{ - 

## 4.5.4 

### Vue

\important{}

\important{Vue}

1. \important{}
2. \important{}
3. \important{}
4. \important{}

## 4.5.5 

### Vue Router

Vue RouterVue.js

#### 

URL

\end{lstlisting}

javascript // router/index.js - import \{ createRouter, createWebHistory
} from `vue-router'

const routes = {[} // \{ path: `/login', name: `Login', component:
Login, meta: \{ requiresAuth: false } // }, \{ path: `/', component:
Layout, // children: {[} \{ path: `/dashboard', name: `Dashboard',
component: Dashboard, meta: \{ title: '`, requiresAuth: true } }, \{
path:'/stations', name: `StationManagement', component:
StationManagement, meta: \{ title: '`, requiresAuth: true } }, // - \{
path:'/stations/:stationId', name: `StationDetail', component:
StationDetail, meta: \{ title: '\,', requiresAuth: true } } {]} }{]}

// const router = createRouter(\{ history: createWebHistory(), // HTML5
History routes })

\begin{lstlisting}
\important{}
- \important{}}children\code{
- \important{}}:stationId\code{}/stations/WS001\code{
- \important{}}meta\code{

#### 

\end{lstlisting}

javascript // - router.beforeEach((to, from, next) =\textgreater{} \{ //
document.title = to.meta.title \textbar\textbar{} '\,'

// if (to.meta.requiresAuth \&\& !userStore.isAuthenticated) \{
next(`/login') // } else \{ next() // } })

\begin{lstlisting}
#### 

\end{lstlisting}

vue

\begin{verbatim}
<!--  -->
<router-link to="/dashboard"></router-link>

<!--  -->
<router-link 
  :to="{ name: 'StationDetail', params: { stationId: station.id } }"
>
  
</router-link>

<!--  -->
<button @click="navigateToAnalysis"></button>
\end{verbatim}

\begin{lstlisting}
\important{}
- }router.push()\code{
- }router.replace()\code{  
- }router.go(n)\code{n

### Pinia

PiniaVue 3

#### 

- - 
- #### Store

\end{lstlisting}

javascript // stores/user.js - import \{ defineStore } from `pinia'
import \{ ref, computed } from `vue'

export const useUserStore = defineStore(`user', () =\textgreater{} \{ //
===== ===== const userInfo = ref(null) const token =
ref(localStorage.getItem(`token') \textbar\textbar{} '\,') const
permissions = ref({[}{]})

// ===== ===== const isAuthenticated = computed(() =\textgreater{} \{
return !!token.value \&\& !!userInfo.value })

const userName = computed(() =\textgreater{} \{ return
userInfo.value?.name \textbar\textbar{} '\,' })

// ===== ===== const login = async (credentials) =\textgreater{} \{ try
\{ const response = await authAPI.login(credentials) if
(response.success) \{ token.value = response.data.token userInfo.value =
response.data.user permissions.value = response.data.permissions

\begin{verbatim}
    // 
    localStorage.setItem('token', token.value)
    return { success: true }
  }
} catch (error) {
  return { success: false, message: '' }
}
\end{verbatim}

}

const logout = () =\textgreater{} \{ token.value = '\,' userInfo.value =
null permissions.value = {[}{]} localStorage.removeItem(`token') }

const hasPermission = (permission) =\textgreater{} \{ return
permissions.value.includes(permission) }

// return \{ userInfo, token, permissions, isAuthenticated, userName,
login, logout, hasPermission } })

\begin{lstlisting}
\important{Store}
1. \important{State}}ref()\code{}reactive()\code{
2. \important{Getters}}computed()\code{
3. \important{Actions}

#### Store

\end{lstlisting}

vue

\begin{verbatim}
<div v-if="userStore.isAuthenticated">
  {{ userStore.userName }}
  <button @click="handleLogout"></button>
</div>
<div v-else>
  <button @click="showLogin = true"></button>
</div>
\end{verbatim}

\begin{lstlisting}

#### Store

\end{lstlisting}

javascript // stores/station.js - import \{ defineStore } from `pinia'
import \{ ref, computed } from `vue' import \{ stationAPI } from
`@/api/station'

export const useStationStore = defineStore(`station', () =\textgreater{}
\{ // ===== ===== const stations = ref({[}{]}) const selectedStation =
ref(null) const loading = ref(false)

// ===== ===== const onlineStations = computed(() =\textgreater{} \{
return stations.value.filter(station =\textgreater{} station.status ===
`online') })

const statusStatistics = computed(() =\textgreater{} (\{ total:
stations.value.length, online: stations.value.filter(s =\textgreater{}
s.status === `online').length, offline: stations.value.filter(s
=\textgreater{} s.status === `offline').length, warning:
stations.value.filter(s =\textgreater{} s.status === `warning').length
}))

// ===== ===== const loadStations = async () =\textgreater{} \{
loading.value = true try \{ const response = await
stationAPI.getStations() if (response.success) \{ stations.value =
response.data.stations } } catch (error) \{ console.error(`:', error)
} finally \{ loading.value = false } }

const updateStationData = (stationId, newData) =\textgreater{} \{ const
index = stations.value.findIndex(s =\textgreater{} s.id === stationId)
if (index \textgreater{} -1) \{ stations.value{[}index{]} = \{
\ldots stations.value{[}index{]}, \ldots newData } } }

const selectStation = (stationId) =\textgreater{} \{
selectedStation.value = stations.value.find(s =\textgreater{} s.id ===
stationId) }

return \{ // stations, selectedStation, loading, // onlineStations,
statusStatistics, // loadStations, updateStationData, selectStation }
})

\begin{lstlisting}
\important{Store}

\end{lstlisting}

vue

\begin{verbatim}
<div class="statistics">
  <div class="stat-item">
    <h3></h3>
    <div class="value">{{ stationStore.statusStatistics.total }}</div>
  </div>
  <div class="stat-item">
    <h3></h3>
    <div class="value">{{ stationStore.statusStatistics.online }}</div>
  </div>
</div>

<div class="station-list">
  <div 
    v-for="station in stationStore.stations" 
    :key="station.id"
    @click="stationStore.selectStation(station.id)"
  >
    {{ station.name }} - {{ station.status }}
  </div>
</div>
\end{verbatim}

\begin{lstlisting}
### Vue.js

#### 

\important{1. }

\end{lstlisting}

vue

\begin{verbatim}
<div class="current-level">{{ level }} </div>
<div class="status" :class="statusClass">{{ statusText }}</div>
\end{verbatim}

\begin{lstlisting}
\important{2. }

- \important{}
- \important{}  
- \important{}

\important{3. }

TypeScriptPropTypes

\end{lstlisting}

javascript // const props = defineProps(\{ stationData: \{ type: Object,
required: true, validator: (value) =\textgreater{} value.id \&\&
value.name }, editable: \{ type: Boolean, default: false } })

// const emit = defineEmits({[}`update', `delete', `select'{]})

\begin{lstlisting}
#### 

\important{}

- \important{}}ref()\code{}reactive()\code{
- \important{}Pinia Store
- \important{}Props/Events

\end{lstlisting}

javascript // const componentState = ref(\{}) // UI const businessStore
= useStationStore() // const globalStore = useAppStore() //

\begin{lstlisting}
#### 

\important{1. }

\end{lstlisting}

javascript // //

\{\{ stations.filter(s =\textgreater{} s.status === `online').length
}}

// const onlineStationsCount = computed(() =\textgreater{}
stations.value.filter(s =\textgreater{} s.status === `online').length )

\begin{lstlisting}
\important{2. }

}v-for\code{key

\end{lstlisting}

vue \textless station-card v-for=``station in stations''
:key=``station.id''\\
:station=``station'' /\textgreater{}

\begin{lstlisting}
\important{3. }

\end{lstlisting}

javascript const routes = {[} \{ path: `/analysis', name:
`DataAnalysis', // component: () =\textgreater{}
import(`@/views/DataAnalysis.vue') }{]}

\begin{lstlisting}
### 

Vue.js

\important{1. }
- Vue.js
- MVVM
- DOM

\important{2. }
- Vue
- - Vue RouterPinia

\important{3. }
- Vue.js
- - 

\important{4. }
- Vue.js
- - 

!!! tip ""
    
    Vue.js
    
    1. \important{}
    2. \important{}
    3. \important{}
    4. \important{}Vue.js
\end{lstlisting}

} station.id.toLowerCase().includes(keyword) \textbar\textbar{}
station.location?.toLowerCase().includes(keyword) ) }

\begin{verbatim}
return filtered
\end{verbatim}

}

// const subscribeRealTimeData = (stationId, callback) =\textgreater{}
\{ if (subscriptions.value.has(stationId)) \{ //
subscriptions.value.get(stationId).callbacks.push(callback) } else \{
// const websocket = new
WebSocket(\code{${import.meta.env.VITE_WS_URL}/stations/${stationId}/realtime})

\begin{verbatim}
  websocket.onmessage = (event) => {
    const data = JSON.parse(event.data)
    
    // 
    const index = stations.value.findIndex(s => s.id === stationId)
    if (index > -1) {
      stations.value[index].measurements = data.measurements
      stations.value[index].lastUpdate = data.timestamp
    }
    
    // 
    const subscription = subscriptions.value.get(stationId)
    if (subscription) {
      subscription.callbacks.forEach(cb => cb(data))
    }
  }
  
  websocket.onerror = (error) => {
    console.error(\code{WebSocket (${stationId}):}, error)
  }
  
  websocket.onclose = () => {
    console.log(\code{WebSocket (${stationId})})
    subscriptions.value.delete(stationId)
  }
  
  subscriptions.value.set(stationId, {
    websocket,
    callbacks: [callback]
  })
}
\end{verbatim}

}

// const unsubscribeRealTimeData = (stationId, callback = null)
=\textgreater{} \{ const subscription =
subscriptions.value.get(stationId) if (!subscription) return

\begin{verbatim}
if (callback) {
  // 
  const index = subscription.callbacks.indexOf(callback)
  if (index > -1) {
    subscription.callbacks.splice(index, 1)
  }
  
  // 
  if (subscription.callbacks.length === 0) {
    subscription.websocket.close()
    subscriptions.value.delete(stationId)
  }
} else {
  // 
  subscription.websocket.close()
  subscriptions.value.delete(stationId)
}
\end{verbatim}

}

// const updateFilters = (newFilters) =\textgreater{} \{ filters.value =
\{ \ldots filters.value, \ldots newFilters } }

// const resetState = () =\textgreater{} \{ stations.value = {[}{]}
selectedStation.value = null loading.value = false error.value = null
pagination.value = \{ current: 1, pageSize: 12, total: 0 }
filters.value = \{ status: {[}{]}, region: '`, stationType:'`,
keyword:'' }

\begin{verbatim}
// WebSocket
subscriptions.value.forEach((subscription) => {
  subscription.websocket.close()
})
subscriptions.value.clear()
\end{verbatim}

}

// store return \{ // stations: readonly(stations), selectedStation:
readonly(selectedStation), loading: readonly(loading), error:
readonly(error), pagination: readonly(pagination), filters,
lastUpdateTime: readonly(lastUpdateTime),

\begin{verbatim}
// 
totalCount,
onlineStations,
offlineStations,
warningStations,
errorStations,
statusStatistics,
regions,
stationTypes,

// 
loadStations,
getStationById,
createStation,
updateStation,
deleteStation,
refreshStationData,
getFilteredStations,
subscribeRealTimeData,
unsubscribeRealTimeData,
updateFilters,
resetState
\end{verbatim}

} }) ```

Vue.jsVue.jsVue.js

\hypertarget{section-31}{%
\section{4.6}\label{section-31}}

Vue CLI

!!! info ``\,''

\hypertarget{section-32}{%
\subsection{4.6.1}\label{section-32}}

\hypertarget{section-33}{%
\subsubsection{}\label{section-33}}

\important{}JavaScriptCSS\code{<script>}

\important{}JavaScriptES6+CSS

\important{}

\important{}``\,''

\important{}

\important{}

\important{}ES6CommonJSAMD

\important{}

\important{}

\hypertarget{section-34}{%
\subsubsection{}\label{section-34}}

\important{1. Scaffolding}

\begin{itemize}
\item
  \important{}
\item
  \important{}
\item
  \important{}
\item
  \important{}
\end{itemize}

\begin{lstlisting}
water-monitoring-frontend/
 public/              # 
    index.html      # 
    favicon.ico     # 
 src/                # 
    components/     # 
       DataChart/  # 
       WaterGauge/ # 
    views/          # 
       Dashboard/  # 
       Monitor/    # 
    router/         # 
    store/          # 
    api/            # API
    utils/          # 
    assets/         # 
 tests/              # 
 package.json        # 
 vue.config.js       # Vue
\end{lstlisting}

\important{2. Module System}

\begin{itemize}
\item
  \important{ES6 Modules}\code{import}\code{export}
\item
  \important{CommonJS}Node.js
\item
  \important{AMD/UMD}
\end{itemize}

\important{}

\important{}import \important{} \important{}

\begin{lstlisting}[style=javascript]
// api/waterLevel.js - API
export const getWaterLevelData = async (stationId) => {
  const response = await fetch(\code{/api/stations/${stationId}/water-level})
  return response.json()
}

export const getWaterLevelHistory = async (stationId, dateRange) => {
  const response = await fetch(\code{/api/stations/${stationId}/history}, {
    method: 'POST',
    body: JSON.stringify(dateRange)
  })
  return response.json()
}

// components/WaterLevelChart.vue - 
import { getWaterLevelData, getWaterLevelHistory } from '@/api/waterLevel'

export default {
  name: 'WaterLevelChart',
  async created() {
    this.currentData = await getWaterLevelData(this.stationId)
    this.historyData = await getWaterLevelHistory(this.stationId, this.dateRange)
  }
}
\end{lstlisting}

\important{3. Build Tools}

\begin{itemize}
\item
  \important{}ES6+ES5
\item
  \important{}HTTP
\item
  \important{}
\item
  \important{}CSS
\end{itemize}

\important{4. Package Manager}

npmyarnpnpm

\begin{itemize}
\item
  \important{}
\item
  \important{}
\item
  \important{}
\item
  \important{}
\end{itemize}

\important{5. }

\begin{itemize}
\item
  \important{}ESLintJavaScript
\item
  \important{}Prettier
\item
  \important{}TypeScript
\item
  \important{}
\end{itemize}

\hypertarget{section-35}{%
\subsubsection{}\label{section-35}}

Web

\important{}

\important{}

\important{}

\important{}

\important{}

\begin{itemize}
\item
  \begin{itemize}
  \tightlist
  \item
  \end{itemize}
\item
  \begin{itemize}
  \tightlist
  \item
  \end{itemize}
\end{itemize}

\hypertarget{vue-cli}{%
\subsection{4.6.2 Vue CLI}\label{vue-cli}}

\hypertarget{vue-cli-1}{%
\subsubsection{Vue CLI}\label{vue-cli-1}}

Vue CLICommand Line InterfaceVue.jsVueVue CLI

Vue CLI\important{}\important{}WebpackBabelESLint

\important{Vue CLI}

\important{}Vue

\important{}TypeScriptPWA

\important{}Vue UI

\important{}Webpack

\hypertarget{vue-cli-2}{%
\subsubsection{Vue CLI}\label{vue-cli-2}}

\important{Vue CLI}

Vue CLINode.jsVue CLINode.js 8.9

\begin{lstlisting}[style=terminal]
# Vue CLI
npm install -g @vue/cli

# 
vue --version

# 5.0.8
\end{lstlisting}

yarn

\begin{lstlisting}[style=terminal]
# yarn
yarn global add @vue/cli

# 
vue --version
\end{lstlisting}

\important{Vue}

Vue CLI

\begin{lstlisting}[style=terminal]
# 
vue create water-monitoring-dashboard

# 
cd water-monitoring-dashboard

# 
npm run serve
\end{lstlisting}

Vue CLI

\begin{lstlisting}[style=terminal]
Vue CLI v5.0.8
? Please pick a preset: (Use arrow keys)
 Default ([Vue 3] babel, eslint) 
  Default ([Vue 2] babel, eslint) 
  Manually select features
\end{lstlisting}

\important{}

\begin{itemize}
\item
  \important{Default ([Vue 3] babel, eslint)}Vue 3BabelESLint
\item
  \important{Default ([Vue 2] babel, eslint)}Vue 2
\item
  \important{Manually select features}
\end{itemize}

``Manually select features''

\begin{lstlisting}[style=terminal]
? Check the features needed for your project: (Press <space> to select, <a> to toggle all, <i> to invert selection)
 Babel
  TypeScript
  Progressive Web App (PWA) Support
  Router
  Vuex
  CSS Pre-processors
  Linter / Formatter
  Unit Testing
  E2E Testing
\end{lstlisting}

\important{}

\begin{itemize}
\item
  \important{Babel}JavaScriptES6+
\item
  \important{TypeScript}
\item
  \important{PWA Support}Web
\item
  \important{Router}
\item
  \important{Vuex}
\item
  \important{CSS Pre-processors}CSSSassLessStylus
\item
  \important{Linter / Formatter}
\item
  \important{Unit Testing}
\item
  \important{E2E Testing}
\end{itemize}

\begin{lstlisting}[style=terminal]
# Vue
? Choose a version of Vue.js that you want to start the project with 
 3.x

# TypeScript
? Use class-style component syntax? No
? Use Babel alongside TypeScript? Yes

# 
? Use history mode for router? Yes

# CSS
? Pick a CSS pre-processor:
 Sass/SCSS (with dart-sass)

# 
? Pick a linter / formatter config:
 ESLint + Prettier

# 
? Pick additional lint features:
 Lint on save
  Lint and fix on commit

# 
? Pick a unit testing solution:
 Jest

# E2E
? Pick an E2E testing solution:
 Cypress

# 
? Where do you prefer placing config for Babel, ESLint, etc.?
 In dedicated config files
\end{lstlisting}

\hypertarget{section-36}{%
\subsubsection{}\label{section-36}}

Vue CLI

\begin{lstlisting}
water-monitoring-dashboard/
 public/                     # 
    index.html             # HTML
    favicon.ico            # 
    manifest.json          # PWA
 src/                       # 
    assets/                # 
       images/           # 
       styles/           # 
       fonts/            # 
    components/            # 
       charts/           # 
          LineChart.vue # 
          BarChart.vue  # 
          PieChart.vue  # 
       common/           # 
          Header.vue    # 
          Sidebar.vue   # 
          Footer.vue    # 
       monitoring/       # 
           StationCard.vue    # 
           DataTable.vue      # 
           AlertPanel.vue     # 
    views/                 # 
       Dashboard.vue     # 
       Monitoring.vue    # 
       Analysis.vue      # 
       Settings.vue      # 
    router/               # 
       index.js         # 
    store/               # Vuex
       index.js        # Store
       modules/        # Store
          user.js    # 
          monitoring.js # 
          settings.js   # 
       getters.js      # 
    api/                # API
       http.js        # HTTP
       monitoring.js  # API
       user.js        # API
       common.js      # API
    utils/              # 
       date.js        # 
       format.js      # 
       validate.js    # 
       constants.js   # 
    plugins/            # Vue
       element.js     # Element UI
       charts.js      # 
    directives/         # 
       loading.js     # 
       permission.js  # 
    filters/            # 
       date.js        # 
       number.js      # 
    App.vue            # 
    main.js            # 
 tests/                  # 
    unit/              # 
       components/    # 
       utils/         # 
       setup.js       # 
    e2e/               # 
        specs/         # 
        support/       # 
 docs/                   # 
 .env                    # 
 .env.development        # 
 .env.production         # 
 .gitignore             # Git
 .eslintrc.js           # ESLint
 babel.config.js        # Babel
 jest.config.js         # Jest
 package.json           # 
 README.md              # 
 vue.config.js          # Vue CLI
\end{lstlisting}

\important{}

\important{src/components/}Vue

\important{src/views/}

\important{src/api/}APIHTTP

\important{src/utils/}

\important{vue.config.js}Vue CLIWebpack

\hypertarget{vue-cli-3}{%
\subsubsection{Vue CLI}\label{vue-cli-3}}

Vue CLI

\important{}

\begin{lstlisting}[style=terminal]
# 
npm run serve

# 
npm run serve -- --port 8888

# 
npm run serve -- --host 0.0.0.0

# 
npm run serve -- --open
\end{lstlisting}

\begin{itemize}
\item
  \important{Hot Reload}
\item
  \important{Hot Module Replacement}
\item
  \important{}
\item
  \important{HTTPS}HTTPS
\end{itemize}

\important{}

\begin{lstlisting}[style=terminal]
# 
npm run build

# 
npm run build -- --analyze

# 
npm run build -- --mode staging

# 
npm run build -- --report
\end{lstlisting}

\important{}

\begin{lstlisting}[style=terminal]
# 
npm run lint

# 
npm run lint -- --fix

# 
npm run lint src/components/Chart.vue
\end{lstlisting}

\important{}

\begin{lstlisting}[style=terminal]
# 
npm run test:unit

# 
npm run test:unit -- --coverage

# E2E
npm run test:e2e

# 
npm run test:unit -- --watch
\end{lstlisting}

\hypertarget{vue-cli-4}{%
\subsubsection{Vue CLI}\label{vue-cli-4}}

Vue CLI

\important{}

\begin{lstlisting}[style=terminal]
# Element Plus UI
vue add element-plus

# PWA
vue add pwa

# Vuetify UI
vue add vuetify

# TypeScript
vue add typescript
\end{lstlisting}

\important{}

\begin{lstlisting}[style=terminal]
# UI
vue add element-plus      # Element Plus UI
vue add ant-design-vue    # Ant Design Vue

# 
vue add echarts          # Apache ECharts
npm install vue-chartjs  # Chart.jsVue

# 
npm install vue-amap     # Vue
npm install vue-baidu-map # Vue

# 
vue add axios            # HTTP
vue add dayjs            # 
npm install lodash       # 

# 
vue add storybook        # 
npm install @vue/devtools # Vue
\end{lstlisting}

\hypertarget{section-37}{%
\subsubsection{}\label{section-37}}

Vue CLI

\important{.env}

\begin{lstlisting}[style=terminal]
# .env - 
VUE_APP_TITLE=
VUE_APP_VERSION=1.0.0

# .env.development - 
NODE_ENV=development
VUE_APP_API_BASE_URL=http://localhost:3000/api
VUE_APP_WS_URL=ws://localhost:3000/ws
VUE_APP_MAP_KEY=development_map_api_key
VUE_APP_DEBUG=true

# .env.production - 
NODE_ENV=production
VUE_APP_API_BASE_URL=https://api.water-monitoring.com
VUE_APP_WS_URL=wss://api.water-monitoring.com/ws
VUE_APP_MAP_KEY=production_map_api_key
VUE_APP_DEBUG=false

# .env.staging - 
NODE_ENV=staging
VUE_APP_API_BASE_URL=https://staging-api.water-monitoring.com
VUE_APP_WS_URL=wss://staging-api.water-monitoring.com/ws
VUE_APP_MAP_KEY=staging_map_api_key
VUE_APP_DEBUG=true
\end{lstlisting}

\important{vue.config.js}

\begin{lstlisting}[style=javascript]
const { defineConfig } = require('@vue/cli-service')
const path = require('path')

module.exports = defineConfig({
  // 
  publicPath: process.env.NODE_ENV === 'production' ? '/water-monitoring/' : '/',
  outputDir: 'dist',
  assetsDir: 'static',
  
  // 
  devServer: {
    port: 8080,
    open: true,
    hot: true,
    //  - 
    proxy: {
      '/api': {
        target: 'http://localhost:3000',
        changeOrigin: true,
        pathRewrite: {
          '^/api': '/api'
        }
      },
      '/ws': {
        target: 'ws://localhost:3000',
        ws: true,
        changeOrigin: true
      }
    }
  },
  
  // 
  configureWebpack: {
    resolve: {
      alias: {
        '@': path.resolve(__dirname, 'src'),
        '@components': path.resolve(__dirname, 'src/components'),
        '@views': path.resolve(__dirname, 'src/views'),
        '@utils': path.resolve(__dirname, 'src/utils'),
        '@api': path.resolve(__dirname, 'src/api'),
        '@assets': path.resolve(__dirname, 'src/assets')
      }
    }
  },
  
  // CSS
  css: {
    loaderOptions: {
      scss: {
        // Sass
        additionalData: \code{
          @import "@/assets/styles/variables.scss";
          @import "@/assets/styles/mixins.scss";
        }
      }
    }
  },
  
  // 
  chainWebpack: config => {
    // console
    if (process.env.NODE_ENV === 'production') {
      config.optimization.minimizer('terser').tap(args => {
        args[0].terserOptions.compress.drop_console = true
        return args
      })
    }
    
    // 
    if (process.env.NODE_ENV === 'production' && process.env.ANALYZE) {
      config
        .plugin('webpack-bundle-analyzer')
        .use(require('webpack-bundle-analyzer').BundleAnalyzerPlugin)
    }
  },
  
  // PWA
  pwa: {
    name: '',
    themeColor: '#1890ff',
    msTileColor: '#1890ff',
    workboxOptions: {
      skipWaiting: true,
      clientsClaim: true
    }
  },
  
  // 
  pluginOptions: {
    // Element Plus
    'element-plus': {
      useSource: true
    }
  }
})
\end{lstlisting}

\hypertarget{webpackvite}{%
\subsection{4.6.3 WebpackVite}\label{webpackvite}}

\hypertarget{section-38}{%
\subsubsection{}\label{section-38}}

WebHTMLCSSJavaScript\code{<script>}\code{<link>}HTML

\important{}JavaScript

\important{}HTTP

\important{}JavaScriptES6+CSS

\important{}

\important{}ES6CommonJSAMD\code{import}\code{export}

\important{}BabelES6+ES5TypeScriptJSX

\important{}HTTP

\important{}Source Map

\hypertarget{webpack}{%
\subsubsection{Webpack}\label{webpack}}

WebpackLoaderPlugin

\important{Webpack}

\important{Entry}WebpackJavaScriptWebpack

\begin{lstlisting}[style=javascript]
// webpack.config.js
module.exports = {
  entry: {
    app: './src/main.js',           // 
    vendor: ['vue', 'vue-router']   // 
  }
}
\end{lstlisting}

\important{Output}

\begin{lstlisting}[style=javascript]
module.exports = {
  output: {
    path: path.resolve(__dirname, 'dist'),
    filename: '[name].[contenthash].js',  // 
    chunkFilename: '[name].[contenthash].js',
    publicPath: '/static/'  // 
  }
}
\end{lstlisting}

\important{Loader}WebpackJavaScriptLoaderCSS

\begin{lstlisting}[style=javascript]
module.exports = {
  module: {
    rules: [
      // Vue
      {
        test: /\.vue$/,
        use: 'vue-loader'
      },
      // JavaScript/TypeScript
      {
        test: /\.(js|ts)$/,
        exclude: /node_modules/,
        use: {
          loader: 'babel-loader',
          options: {
            presets: ['@babel/preset-env', '@babel/preset-typescript']
          }
        }
      },
      // CSS
      {
        test: /\.css$/,
        use: [
          'vue-style-loader',
          'css-loader',
          'postcss-loader'
        ]
      },
      // Sass/SCSS
      {
        test: /\.scss$/,
        use: [
          'vue-style-loader',
          'css-loader',
          'sass-loader'
        ]
      },
      // 
      {
        test: /\.(png|jpg|gif|svg)$/,
        type: 'asset',
        parser: {
          dataUrlCondition: {
            maxSize: 8 \highlight{ 1024 // 8KBbase64
          }
        }
      },
      // 
      {
        test: /\.(woff|woff2|eot|ttf|otf)$/,
        type: 'asset/resource'
      }
    ]
  }
}
\end{lstlisting}

\important{Plugin}HTML

\begin{lstlisting}[style=javascript]
const HtmlWebpackPlugin = require('html-webpack-plugin')
const MiniCssExtractPlugin = require('mini-css-extract-plugin')
const { CleanWebpackPlugin } = require('clean-webpack-plugin')

module.exports = {
  plugins: [
    // 
    new CleanWebpackPlugin(),
    
    // HTML
    new HtmlWebpackPlugin({
      template: 'public/index.html',
      title: '',
      minify: {
        removeComments: true,
        collapseWhitespace: true
      }
    }),
    
    // CSS
    new MiniCssExtractPlugin({
      filename: 'css/[name].[contenthash].css',
      chunkFilename: 'css/[name].[contenthash].css'
    }),
    
    // 
    new webpack.DefinePlugin({
      'process.env.NODE_ENV': JSON.stringify(process.env.NODE_ENV),
      'process.env.VUE_APP_API_URL': JSON.stringify(process.env.VUE_APP_API_URL)
    })
  ]
}
\end{lstlisting}

\important{Webpack}

\important{}Webpack

\important{}Tree ShakingHot Module Replacement

\important{}

\important{}

\important{Webpack}

\important{}Webpack

\important{}Webpack

\important{}

\hypertarget{vite}{%
\subsubsection{Vite}\label{vite}}

ViteVue.jsES

\important{Vite}

Vite\important{}

\important{}ES

\important{}RollupRollupES

\important{Vite}

\important{}1-2

\begin{lstlisting}[style=terminal]
# Webpack30
npm run serve
# Starting development server...
# webpack compiled successfully in 45.67s

# Vite1-2
npm run dev
# Local: http://localhost:3000/
# ready in 524ms
\end{lstlisting}

\important{}100ms

\important{ES}ES

\begin{lstlisting}[style=javascript]
// Vite
import { createApp } from 'vue'
import App from './App.vue'
import router from './router'

// 
\end{lstlisting}

\important{TypeScript}TypeScript

\important{CSS}SassLessStylus

\important{Vite}

\begin{lstlisting}[style=javascript]
// vite.config.js
import { defineConfig } from 'vite'
import vue from '@vitejs/plugin-vue'
import { resolve } from 'path'

export default defineConfig({
  // 
  plugins: [vue()],
  
  // 
  resolve: {
    alias: {
      '@': resolve(__dirname, 'src'),
      '@components': resolve(__dirname, 'src/components'),
      '@utils': resolve(__dirname, 'src/utils')
    }
  },
  
  // 
  server: {
    port: 3000,
    open: true,
    // 
    proxy: {
      '/api': {
        target: 'http://localhost:8080',
        changeOrigin: true,
        rewrite: (path) => path.replace(/^\/api/, '')
      }
    }
  },
  
  // 
  build: {
    outDir: 'dist',
    assetsDir: 'assets',
    sourcemap: false,
    // 
    rollupOptions: {
      output: {
        chunkFileNames: 'js/[name]-[hash].js',
        entryFileNames: 'js/[name]-[hash].js',
        assetFileNames: '[ext]/[name]-[hash].[ext]'
      }
    }
  },
  
  // CSS
  css: {
    preprocessorOptions: {
      scss: {
        additionalData: \code{@import "@/assets/styles/variables.scss";}
      }
    }
  },
  
  // 
  define: {
    __VUE_OPTIONS_API__: true,
    __VUE_PROD_DEVTOOLS__: false
  }
})
\end{lstlisting}

\hypertarget{webpack-vs-vite}{%
\subsubsection{Webpack vs Vite}\label{webpack-vs-vite}}

\begin{longtable}[]{@{}lll@{}}
\toprule\noalign{}
& Webpack & Vite \\
\midrule\noalign{}
\endhead
\bottomrule\noalign{}
\endlastfoot
\important{} & & 1-2 \\
\important{} & & \textless100ms \\
\important{} & & \\
\important{} & & \\
\important{} & & \\
\important{} & Webpack & Rollup \\
\important{} & & \\
\important{} & Source Map & ES \\
\important{} & & \\
\end{longtable}

\important{}

\important{Webpack}

\begin{itemize}
\item
  \begin{itemize}
  \tightlist
  \item
  \end{itemize}
\item
  Webpack
\item
  Webpack
\end{itemize}

\important{Vite}

\begin{itemize}
\item
  \begin{itemize}
  \tightlist
  \item
    Vue 3React
  \end{itemize}
\item
  \begin{itemize}
  \tightlist
  \item
  \end{itemize}
\end{itemize}

\important{}

Vite

\important{}Vite

\important{}Vue 3TypeScriptVite

\important{}Vite

\important{}ViteVite

\hypertarget{section-39}{%
\subsection{4.6.4}\label{section-39}}

\hypertarget{section-40}{%
\subsubsection{}\label{section-40}}

Web

\important{}

\important{}\code{file://}HTMLAJAXWeb

\important{}APIJSONPCORS

\important{}

\important{}Source Map

\important{}HTTPHTTPS

\important{}

\important{}API

\important{Source Map}Source Map

\important{}

\hypertarget{section-41}{%
\subsubsection{}\label{section-41}}

Vue CLI

\important{}

\begin{lstlisting}[style=javascript]
// vue.config.js
module.exports = {
  devServer: {
    // 
    host: '0.0.0.0',        // 
    port: 8080,             // 
    https: false,           // HTTPS
    open: true,             // 
    
    // 
    hot: true,              // 
    liveReload: true,       // 
    
    // 
    client: {
      overlay: {            // 
        errors: true,       // 
        warnings: false     // 
      },
      progress: true        // 
    }
  }
}
\end{lstlisting}

\important{}

API

\begin{lstlisting}[style=javascript]
module.exports = {
  devServer: {
    proxy: {
      // API
      '/api': {
        target: 'http://localhost:3000',    // API
        changeOrigin: true,                 // 
        pathRewrite: {
          '^/api': '/api'                   // 
        },
        logLevel: 'debug'                   // 
      },
      
      // WebSocket
      '/socket.io': {
        target: 'http://localhost:3000',
        ws: true,                           // WebSocket
        changeOrigin: true
      },
      
      // API
      '/api/monitoring': {
        target: 'http://monitoring-service:8081',
        changeOrigin: true,
        pathRewrite: {
          '^/api/monitoring': ''
        },
        // 
        onProxyReq: (proxyReq, req, res) => {
          console.log(':', req.method, req.url)
          // 
          proxyReq.setHeader('Authorization', 'Bearer ' + getToken())
        },
        // 
        onProxyRes: (proxyRes, req, res) => {
          console.log(':', proxyRes.statusCode, req.url)
        }
      },
      
      // 
      '/api/data': {
        target: 'http://localhost:3000',
        changeOrigin: true,
        bypass: function (req, res, proxyOptions) {
          // mock
          if (process.env.USE_MOCK === 'true') {
            return '/mock' + req.url
          }
        }
      }
    }
  }
}
\end{lstlisting}

\important{HTTPS}

HTTPSAPIHTTPS

\begin{lstlisting}[style=javascript]
const fs = require('fs')
const path = require('path')

module.exports = {
  devServer: {
    https: {
      key: fs.readFileSync(path.resolve(__dirname, 'certs/server.key')),
      cert: fs.readFileSync(path.resolve(__dirname, 'certs/server.crt'))
    },
    // 
    // https: true
  }
}
\end{lstlisting}

\hypertarget{section-42}{%
\subsubsection{}\label{section-42}}

Hot Reload

\important{}

\important{}

\important{}

\important{}

\important{WebSocket}

\begin{lstlisting}[style=javascript]
// 
// 1. 
const chokidar = require('chokidar')
const watcher = chokidar.watch('./src', {
  ignored: /node_modules/,
  persistent: true
})

watcher.on('change', (path) => {
  console.log(':', path)
  // 2. 
  recompile(path)
})

// 3. 
function onCompileComplete(updatedModules) {
  // WebSocket
  webSocketServer.clients.forEach(client => {
    client.send(JSON.stringify({
      type: 'hot-update',
      modules: updatedModules
    }))
  })
}
\end{lstlisting}

\important{Vue}

Vue

\important{}

\begin{lstlisting}
<!--  -->
<template>
  <div>
    <input v-model="message" />
    <p>{{ message }}</p>
  </div>
</template>

<script>
export default {
  data() {
    return {
      message: ''
    }
  }
}
</script>

<!-- input -->
<template>
  <div class="container">
    <h3></h3>
    <input v-model="message" />
    <p class="message">{{ message }}</p>
  </div>
</template>
\end{lstlisting}

\important{}

\begin{itemize}
\item
  \important{}
\item
  \important{}CSS
\item
  \important{}
\end{itemize}

\important{}

\begin{lstlisting}[style=javascript]
// vue.config.js
module.exports = {
  devServer: {
    hot: true,
    hotOnly: true,  // 
  },
  
  chainWebpack: config => {
    // 
    if (process.env.NODE_ENV === 'development') {
      // 
      config.cache(true)
      
      // 
      config.watchOptions({
        ignored: /node_modules/,
        aggregateTimeout: 300,
        poll: 1000
      })
    }
  }
}
\end{lstlisting}

\hypertarget{section-43}{%
\subsubsection{}\label{section-43}}

\important{}

Vue CLI

\begin{lstlisting}[style=terminal]
.env                # 
.env.local          # git
.env.development    # 
.env.development.local  # 
.env.production     # 
.env.production.local   # 
\end{lstlisting}

\important{}

\begin{lstlisting}[style=terminal]
# .env - 
VUE_APP_TITLE=
VUE_APP_VERSION=2.1.0
VUE_APP_BUILD_TIME=2024-01-15

# .env.development - 
NODE_ENV=development
VUE_APP_BASE_API=http://localhost:3000/api
VUE_APP_WS_URL=ws://localhost:3000/ws

# 
VUE_APP_MAP_TYPE=amap
VUE_APP_AMAP_KEY=dev_amap_key_123456
VUE_APP_BAIDU_KEY=dev_baidu_key_123456

# 
VUE_APP_DATA_UPDATE_INTERVAL=5000
VUE_APP_MAX_CHART_POINTS=1000
VUE_APP_ENABLE_MOCK_DATA=true

# 
VUE_APP_ENABLE_PWA=false
VUE_APP_ENABLE_DEBUG=true
VUE_APP_ENABLE_PERFORMANCE_MONITOR=true

# .env.production - 
NODE_ENV=production
VUE_APP_BASE_API=https://api.water-monitoring.com
VUE_APP_WS_URL=wss://api.water-monitoring.com/ws

# 
VUE_APP_MAP_TYPE=amap
VUE_APP_AMAP_KEY=prod_amap_key_789012
VUE_APP_BAIDU_KEY=prod_baidu_key_789012

# 
VUE_APP_DATA_UPDATE_INTERVAL=3000
VUE_APP_MAX_CHART_POINTS=5000
VUE_APP_ENABLE_MOCK_DATA=false

# 
VUE_APP_ENABLE_PWA=true
VUE_APP_ENABLE_DEBUG=false
VUE_APP_ENABLE_PERFORMANCE_MONITOR=false
\end{lstlisting}

\important{}

\begin{lstlisting}[style=javascript]
// src/config/index.js - 
export const config = {
  // API
  api: {
    baseURL: process.env.VUE_APP_BASE_API,
    timeout: 10000,
    withCredentials: true
  },
  
  // WebSocket
  websocket: {
    url: process.env.VUE_APP_WS_URL,
    reconnectInterval: 5000,
    maxReconnectAttempts: 10
  },
  
  // 
  map: {
    type: process.env.VUE_APP_MAP_TYPE,
    keys: {
      amap: process.env.VUE_APP_AMAP_KEY,
      baidu: process.env.VUE_APP_BAIDU_KEY
    },
    defaultCenter: [116.397428, 39.90923],
    defaultZoom: 11
  },
  
  // 
  monitoring: {
    updateInterval: parseInt(process.env.VUE_APP_DATA_UPDATE_INTERVAL),
    maxChartPoints: parseInt(process.env.VUE_APP_MAX_CHART_POINTS),
    enableMockData: process.env.VUE_APP_ENABLE_MOCK_DATA === 'true'
  },
  
  // 
  features: {
    pwa: process.env.VUE_APP_ENABLE_PWA === 'true',
    debug: process.env.VUE_APP_ENABLE_DEBUG === 'true',
    performanceMonitor: process.env.VUE_APP_ENABLE_PERFORMANCE_MONITOR === 'true'
  },
  
  // 
  app: {
    title: process.env.VUE_APP_TITLE,
    version: process.env.VUE_APP_VERSION,
    buildTime: process.env.VUE_APP_BUILD_TIME
  }
}

// src/utils/api.js - API
import axios from 'axios'
import { config } from '@/config'

const apiClient = axios.create({
  baseURL: config.api.baseURL,
  timeout: config.api.timeout,
  withCredentials: config.api.withCredentials
})

// 
apiClient.interceptors.request.use(
  (config) => {
    if (process.env.NODE_ENV === 'development') {
      console.log('API:', config.method?.toUpperCase(), config.url)
    }
    return config
  },
  (error) => {
    return Promise.reject(error)
  }
)

// src/components/Map/index.vue - 
<template>
  <div class="map-container">
    <component 
      :is="mapComponent" 
      :api-key="mapApiKey"
      :center="mapCenter"
      :zoom="mapZoom"
      @ready="onMapReady"
    />
  </div>
</template>

<script>
import { config } from '@/config'
import AmapComponent from './AMap.vue'
import BaiduMapComponent from './BaiduMap.vue'

export default {
  name: 'MapContainer',
  
  computed: {
    mapComponent() {
      return config.map.type === 'amap' ? AmapComponent : BaiduMapComponent
    },
    
    mapApiKey() {
      return config.map.keys[config.map.type]
    },
    
    mapCenter() {
      return this.center || config.map.defaultCenter
    },
    
    mapZoom() {
      return this.zoom || config.map.defaultZoom
    }
  },
  
  methods: {
    onMapReady(map) {
      if (config.features.debug) {
        console.log(':', map)
      }
      this.$emit('map-ready', map)
    }
  }
}
</script>
\end{lstlisting}

\important{}

\begin{lstlisting}[style=javascript]
// src/utils/env.js - 
class EnvironmentManager {
  constructor() {
    this.currentEnv = this.detectEnvironment()
    this.config = this.loadConfig()
  }
  
  detectEnvironment() {
    const hostname = window.location.hostname
    
    if (hostname === 'localhost' || hostname === '127.0.0.1') {
      return 'development'
    } else if (hostname.includes('staging') || hostname.includes('test')) {
      return 'staging'
    } else {
      return 'production'
    }
  }
  
  loadConfig() {
    const configs = {
      development: {
        apiUrl: 'http://localhost:3000/api',
        wsUrl: 'ws://localhost:3000/ws',
        debugMode: true
      },
      staging: {
        apiUrl: 'https://staging-api.water-monitoring.com',
        wsUrl: 'wss://staging-api.water-monitoring.com/ws',
        debugMode: true
      },
      production: {
        apiUrl: 'https://api.water-monitoring.com',
        wsUrl: 'wss://api.water-monitoring.com/ws',
        debugMode: false
      }
    }
    
    return {
      ...configs[this.currentEnv],
      environment: this.currentEnv
    }
  }
  
  getConfig(key) {
    return key ? this.config[key] : this.config
  }
  
  switchEnvironment(env) {
    this.currentEnv = env
    this.config = this.loadConfig()
    // 
    window.dispatchEvent(new CustomEvent('env-config-changed', {
      detail: this.config
    }))
  }
}

export const envManager = new EnvironmentManager()
\end{lstlisting}

\hypertarget{section-44}{%
\subsection{4.6.5}\label{section-44}}

\hypertarget{section-45}{%
\subsubsection{}\label{section-45}}

\important{}

\important{}

\important{}

\important{}

\begin{lstlisting}[style=javascript]
// 
function getStationData(stationId,dateRange){
    if(!stationId)return null;
    const data=fetchData( stationId , dateRange );
    return{
        id:stationId,
        data:data,
        timestamp:new Date( ).toISOString()
    };
}

// 
function getStationData(stationId, dateRange) {
  if (!stationId) {
    return null;
  }
  
  const data = fetchData(stationId, dateRange);
  
  return {
    id: stationId,
    data: data,
    timestamp: new Date().toISOString()
  };
}
\end{lstlisting}

\important{}

\important{}

\begin{lstlisting}[style=javascript]
// 
const d = new Date();
const u = users.filter(x => x.active);
const calc = (a, b) => a } b \highlight{ 0.1;

// 
const currentDate = new Date();
const activeUsers = users.filter(user => user.isActive);
const calculateDiscountPrice = (originalPrice, discountRate) => 
  originalPrice } discountRate \highlight{ 0.1;
\end{lstlisting}

\important{}

\begin{lstlisting}[style=javascript]
// 
function processWaterData(stationId, data, options) {
  if (!stationId || !data) return null;
  
  let result = {};
  if (options.includeHistory) {
    // 50
  }
  if (options.calculateAverage) {
    // 30
  }
  if (options.detectAnomalies) {
    // 40
  }
  // ... 
  
  return result;
}

// 
function processWaterData(stationId, data, options) {
  if (!stationId || !data) {
    return null;
  }
  
  const result = {
    stationId,
    rawData: data,
    processedAt: new Date()
  };
  
  if (options.includeHistory) {
    result.history = processHistoryData(data, options.historyDays);
  }
  
  if (options.calculateAverage) {
    result.averages = calculateDataAverages(data, options.averagePeriod);
  }
  
  if (options.detectAnomalies) {
    result.anomalies = detectDataAnomalies(data, options.anomalyThreshold);
  }
  
  return result;
}
\end{lstlisting}

\important{}

\begin{lstlisting}[style=javascript]
// ESLint
function updateStationStatus(station) {
  // 
  if (stationId) {  //  station.id
    // 
    station.status.active = true;  // statusnull
    
    // 
    updateDatabase(station);  // 
  }
}

// 
function updateStationStatus(station) {
  if (!station || !station.id) {
    throw new Error('Invalid station object');
  }
  
  if (!station.status) {
    station.status = {};
  }
  
  station.status.active = true;
  
  try {
    return updateDatabase(station);
  } catch (error) {
    console.error('Failed to update station:', error);
    throw error;
  }
}
\end{lstlisting}

\hypertarget{eslint}{%
\subsubsection{ESLint}\label{eslint}}

ESLintJavaScriptTypeScript

\important{ESLint}

ESLintASTESLint

\important{}JavaScript

\important{}

\important{}

\important{}

\important{ESLint}

\begin{lstlisting}[style=javascript]
// .eslintrc.js
module.exports = {
  // 
  env: {
    node: true,
    browser: true,
    es2021: true
  },
  
  // 
  extends: [
    'plugin:vue/vue3-essential',      // Vue 3
    'eslint:recommended',             // ESLint
    '@vue/typescript/recommended',    // TypeScript
    '@vue/prettier',                  // Prettier
    '@vue/prettier/@typescript-eslint' // TypeScript Prettier
  ],
  
  // 
  parserOptions: {
    ecmaVersion: 2021,
    parser: '@typescript-eslint/parser',
    sourceType: 'module'
  },
  
  // 
  plugins: [
    'vue',
    '@typescript-eslint',
    'import'
  ],
  
  // 
  rules: {
    // Vue
    'vue/multi-word-component-names': 'off',
    'vue/component-definition-name-casing': ['error', 'PascalCase'],
    'vue/component-name-in-template-casing': ['error', 'kebab-case'],
    'vue/prop-name-casing': ['error', 'camelCase'],
    'vue/attribute-hyphenation': ['error', 'always'],
    
    // JavaScript
    'no-console': process.env.NODE_ENV === 'production' ? 'warn' : 'off',
    'no-debugger': process.env.NODE_ENV === 'production' ? 'warn' : 'off',
    'no-unused-vars': 'error',
    'no-undef': 'error',
    'no-duplicate-keys': 'error',
    'no-unreachable': 'error',
    
    // 
    'indent': ['error', 2],
    'quotes': ['error', 'single'],
    'semi': ['error', 'always'],
    'comma-dangle': ['error', 'never'],
    'object-curly-spacing': ['error', 'always'],
    'array-bracket-spacing': ['error', 'never'],
    
    // 
    'func-names': ['error', 'as-needed'],
    'camelcase': ['error', { properties: 'never' }],
    'max-len': ['error', { code: 100, ignoreUrls: true }],
    'max-params': ['error', 4],
    'complexity': ['error', 10],
    
    // Import
    'import/order': ['error', {
      groups: [
        'builtin',
        'external', 
        'internal',
        'parent',
        'sibling',
        'index'
      ],
      'newlines-between': 'always'
    }],
    'import/no-unresolved': 'error',
    'import/no-duplicates': 'error',
    
    // TypeScript
    '@typescript-eslint/no-explicit-any': 'warn',
    '@typescript-eslint/explicit-function-return-type': 'off',
    '@typescript-eslint/no-unused-vars': 'error',
    '@typescript-eslint/prefer-const': 'error'
  },
  
  // 
  overrides: [
    {
      files: ['}.vue'],
      rules: {
        'max-len': 'off'  // Vue
      }
    },
    {
      files: ['tests/**/*'],
      env: {
        jest: true
      },
      rules: {
        'no-console': 'off'  // console
      }
    }
  ],
  
  // 
  globals: {
    defineProps: 'readonly',
    defineEmits: 'readonly',
    defineExpose: 'readonly',
    withDefaults: 'readonly'
  }
};
\end{lstlisting}

\important{ESLint}

\important{}

\begin{lstlisting}[style=javascript]
// no-undef - 
function calculateAverage(data) {
  // totalValue
  return totalValue / data.length;  // ESLint
}

// 
function calculateAverage(data) {
  const totalValue = data.reduce((sum, item) => sum + item.value, 0);
  return totalValue / data.length;
}

// no-unreachable - 
function processData(data) {
  if (!data) {
    return null;
    console.log('This will never execute');  // ESLint
  }
  return data;
}
\end{lstlisting}

\important{}

\begin{lstlisting}[style=javascript]
// complexity - 
function processStationData(station, options) {
  // ESLint
  if (station.type === 'water-level') {
    if (options.includeHistory) {
      if (options.period === 'daily') {
        if (options.format === 'chart') {
          // 
        }
      }
    }
  }
}

// 
function processStationData(station, options) {
  const processor = getDataProcessor(station.type);
  return processor.process(station, options);
}

// max-params - 
// 
function createChart(type, data, width, height, colors, options, callbacks) {
  // ...
}

// 
function createChart(type, data, config) {
  const { width, height, colors, options, callbacks } = config;
  // ...
}
\end{lstlisting}

\important{}

\begin{lstlisting}[style=javascript]
// object-curly-spacing - 
const station = {id: 1, name: 'Station A'};  // 
const station = { id: 1, name: 'Station A' };  // 

// quotes - 
const message = "Hello World";  // single
const message = 'Hello World';  // 

// comma-dangle - 
const config = {
  api: 'http://localhost:3000',
  timeout: 5000,  // 
};
\end{lstlisting}

\hypertarget{prettier}{%
\subsubsection{Prettier}\label{prettier}}

PrettierESLintESLintPrettier

\important{Prettier}

Prettier''opinionated''

\important{Prettier}

\begin{lstlisting}[style=javascript]
// .prettierrc.js
module.exports = {
  // 
  printWidth: 80,                    // 
  tabWidth: 2,                       // 
  useTabs: false,                    // tab
  semi: true,                        // 
  singleQuote: true,                 // 
  quoteProps: 'as-needed',           // 
  
  // 
  trailingComma: 'es5',              // 
  bracketSpacing: true,              // 
  bracketSameLine: false,            // 
  
  // 
  arrowParens: 'avoid',              // 
  
  // Vue
  vueIndentScriptAndStyle: false,    // Vuescriptstyle
  
  // HTML
  htmlWhitespaceSensitivity: 'css',  // HTML
  
  // 
  endOfLine: 'lf',                   // LF
  
  // 
  embeddedLanguageFormatting: 'auto'
};
\end{lstlisting}

\important{PrettierESLint}

ESLintPrettier

\begin{lstlisting}[style=terminal]
# 
npm install --save-dev prettier eslint-config-prettier eslint-plugin-prettier

# Vue
npm install --save-dev @vue/eslint-config-prettier
\end{lstlisting}

\begin{lstlisting}[style=javascript]
// .eslintrc.js - 
module.exports = {
  extends: [
    'eslint:recommended',
    'plugin:vue/vue3-recommended',
    '@vue/prettier',                 // 
    '@vue/prettier/@typescript-eslint'
  ],
  
  plugins: ['prettier'],
  
  rules: {
    'prettier/prettier': 'error',    // PrettierESLint
  }
};
\end{lstlisting}

\important{VS Code}

\begin{lstlisting}[style=json]
// .vscode/settings.json
{
  "editor.formatOnSave": true,
  "editor.defaultFormatter": "esbenp.prettier-vscode",
  "editor.codeActionsOnSave": {
    "source.fixAll.eslint": true
  },
  "[vue]": {
    "editor.defaultFormatter": "esbenp.prettier-vscode"
  },
  "[javascript]": {
    "editor.defaultFormatter": "esbenp.prettier-vscode"
  },
  "[typescript]": {
    "editor.defaultFormatter": "esbenp.prettier-vscode"
  }
}
\end{lstlisting}

\hypertarget{git-hooks}{%
\subsubsection{Git Hooks}\label{git-hooks}}

Git HooksGitGitGit Hooks

\important{Husky}

HuskyGit HooksGit Hooks

\begin{lstlisting}[style=terminal]
# Husky
npm install --save-dev husky

# Git Hooks
npx husky install

# package.json
npm set-script prepare "husky install"
\end{lstlisting}

\begin{lstlisting}[style=javascript]
// package.json
{
  "scripts": {
    "prepare": "husky install",
    "lint": "eslint --ext .js,.vue,.ts src",
    "lint:fix": "eslint --ext .js,.vue,.ts src --fix",
    "format": "prettier --write src/**/\highlight{.{js,vue,ts}"
  },
  "lint-staged": {
    "}.{js,vue,ts}": [
      "eslint --fix",
      "prettier --write"
    ],
    "\highlight{.{css,scss,vue}": [
      "prettier --write"
    ]
  }
}
\end{lstlisting}

\important{Pre-commit Hook}

\begin{lstlisting}[style=terminal]
# pre-commit
npx husky add .husky/pre-commit "npx lint-staged"
\end{lstlisting}

\begin{lstlisting}[style=terminal]
#!/bin/sh
# .husky/pre-commit

. "$(dirname "$0")/_/husky.sh"

echo "\faSearch Running pre-commit checks..."

# lint-staged
npx lint-staged

# TypeScript
echo "\faSearch Checking TypeScript types..."
npx vue-tsc --noEmit

# 
echo "🧪 Running unit tests..."
npm run test:unit

echo "\faCheck Pre-commit checks passed!"
\end{lstlisting}

\important{Commit Message}

commitlint

\begin{lstlisting}[style=terminal]
# commitlint
npm install --save-dev @commitlint/cli @commitlint/config-conventional

# commit-msg
npx husky add .husky/commit-msg "npx --no-install commitlint --edit $1"
\end{lstlisting}

\begin{lstlisting}[style=javascript]
// commitlint.config.js
module.exports = {
  extends: ['@commitlint/config-conventional'],
  rules: {
    'type-enum': [
      2,
      'always',
      [
        'feat',     // 
        'fix',      // bug
        'docs',     // 
        'style',    // 
        'refactor', // 
        'test',     // 
        'chore',    // 
        'perf',     // 
        'ci',       // CI/CD
        'revert'    // 
      ]
    ],
    'type-case': [2, 'always', 'lower-case'],
    'type-empty': [2, 'never'],
    'scope-empty': [2, 'never'],
    'scope-enum': [
      2,
      'always',
      [
        'components',  // 
        'views',       // 
        'api',         // API
        'utils',       // 
        'config',      // 
        'monitoring',  // 
        'charts',      // 
        'map',         // 
        'auth',        // 
        'build',       // 
        'docs'         // 
      ]
    ],
    'subject-case': [2, 'never', ['sentence-case', 'start-case', 'pascal-case', 'upper-case']],
    'subject-empty': [2, 'never'],
    'subject-full-stop': [2, 'never', '.'],
    'header-max-length': [2, 'always', 72]
  }
};
\end{lstlisting}

\important{}

\begin{lstlisting}[style=terminal]
# 
feat(monitoring): add real-time water level chart component

# bug  
fix(api): resolve station data fetching timeout issue

# 
docs(readme): update installation and setup instructions

# 
refactor(components): extract common chart logic to base class

# 
perf(charts): implement virtual scrolling for large datasets
\end{lstlisting}

\hypertarget{typescript}{%
\subsubsection{TypeScript}\label{typescript}}

TypeScriptJavaScript

\important{TypeScript}

\begin{lstlisting}[style=json]
// tsconfig.json
{
  "compilerOptions": {
    "target": "ES2020",
    "lib": [
      "ES2020",
      "DOM",
      "DOM.Iterable",
      "ES6"
    ],
    "allowJs": true,
    "skipLibCheck": true,
    "esModuleInterop": true,
    "allowSyntheticDefaultImports": true,
    "strict": true,
    "forceConsistentCasingInFileNames": true,
    "module": "ESNext",
    "moduleResolution": "Node",
    "resolveJsonModule": true,
    "isolatedModules": true,
    "noEmit": true,
    "jsx": "preserve",
    "baseUrl": ".",
    "paths": {
      "@/}": [
        "src/*"
      ]
    }
  },
  "include": [
    "src/**/*.ts",
    "src/**/*.tsx", 
    "src/**/*.vue",
    "tests/**/*.ts",
    "tests/**/\highlight{.tsx"
  ],
  "exclude": [
    "node_modules"
  ]
}
\end{lstlisting}

\important{VueTypeScript}

\begin{lstlisting}
<!-- components/WaterLevelChart.vue -->
<template>
  <div class="water-level-chart">
    <div class="chart-header">
      <h3>{{ title }}</h3>
      <div class="chart-controls">
        <select v-model="selectedPeriod" @change="onPeriodChange">
          <option v-for="period in periods" :key="period.value" :value="period.value">
            {{ period.label }}
          </option>
        </select>
      </div>
    </div>
    <div ref="chartContainer" class="chart-container"></div>
  </div>
</template>

<script setup lang="ts">
import { ref, onMounted, watch, computed } from 'vue';
import } as echarts from 'echarts';

// 
interface WaterLevelData {
  timestamp: string;
  value: number;
  stationId: string;
  stationName: string;
}

interface ChartPeriod {
  value: string;
  label: string;
  days: number;
}

interface Props {
  stationId: string;
  title?: string;
  height?: number;
  autoRefresh?: boolean;
  refreshInterval?: number;
}

interface Emits {
  (e: 'data-loaded', data: WaterLevelData[]): void;
  (e: 'period-changed', period: string): void;
  (e: 'chart-ready', chart: echarts.ECharts): void;
}

// Props
const props = withDefaults(defineProps<Props>(), {
  title: '',
  height: 400,
  autoRefresh: true,
  refreshInterval: 30000
});

// Emits
const emit = defineEmits<Emits>();

// 
const chartContainer = ref<HTMLDivElement>();
const chart = ref<echarts.ECharts>();
const chartData = ref<WaterLevelData[]>([]);
const selectedPeriod = ref<string>('7d');
const loading = ref<boolean>(false);
const error = ref<string | null>(null);

// 
const periods: ChartPeriod[] = [
  { value: '1d', label: '1', days: 1 },
  { value: '7d', label: '7', days: 7 },
  { value: '30d', label: '30', days: 30 },
  { value: '90d', label: '90', days: 90 }
];

// 
const currentPeriod = computed((): ChartPeriod => {
  return periods.find(p => p.value === selectedPeriod.value) || periods[1];
});

const chartOptions = computed((): echarts.EChartOption => {
  const data = chartData.value.map(item => [
    new Date(item.timestamp).getTime(),
    item.value
  ]);

  return {
    title: {
      text: props.title,
      left: 'center'
    },
    tooltip: {
      trigger: 'axis',
      formatter: (params: any) => {
        const point = params[0];
        return \code{
          <div>
            <div>: ${new Date(point.data[0]).toLocaleString()}</div>
            <div>: ${point.data[1]} </div>
          </div>
        };
      }
    },
    xAxis: {
      type: 'time',
      name: ''
    },
    yAxis: {
      type: 'value',
      name: ' ()'
    },
    series: [{
      name: '',
      type: 'line',
      data: data,
      smooth: true,
      symbol: 'circle',
      symbolSize: 4,
      lineStyle: {
        color: '#1890ff',
        width: 2
      },
      areaStyle: {
        color: new echarts.graphic.LinearGradient(0, 0, 0, 1, [
          { offset: 0, color: 'rgba(24, 144, 255, 0.3)' },
          { offset: 1, color: 'rgba(24, 144, 255, 0.05)' }
        ])
      }
    }]
  };
});

// 
const fetchData = async (): Promise<void> => {
  if (!props.stationId) {
    error.value = 'Station ID is required';
    return;
  }

  loading.value = true;
  error.value = null;

  try {
    const response = await fetch(\code{/api/stations/${props.stationId}/water-level}, {
      method: 'POST',
      headers: {
        'Content-Type': 'application/json'
      },
      body: JSON.stringify({
        period: currentPeriod.value.days,
        limit: 1000
      })
    });

    if (!response.ok) {
      throw new Error(\code{HTTP error! status: ${response.status}});
    }

    const result = await response.json();
    chartData.value = result.data as WaterLevelData[];
    emit('data-loaded', chartData.value);
  } catch (err) {
    error.value = err instanceof Error ? err.message : 'Unknown error occurred';
    console.error('Failed to fetch water level data:', err);
  } finally {
    loading.value = false;
  }
};

const initChart = (): void => {
  if (!chartContainer.value) {
    console.error('Chart container not found');
    return;
  }

  chart.value = echarts.init(chartContainer.value);
  chart.value.setOption(chartOptions.value);
  emit('chart-ready', chart.value);

  // 
  window.addEventListener('resize', handleResize);
};

const updateChart = (): void => {
  if (chart.value) {
    chart.value.setOption(chartOptions.value);
  }
};

const handleResize = (): void => {
  if (chart.value) {
    chart.value.resize();
  }
};

const onPeriodChange = (): void => {
  emit('period-changed', selectedPeriod.value);
  fetchData();
};

// 
onMounted(async () => {
  await fetchData();
  initChart();

  // 
  if (props.autoRefresh && props.refreshInterval > 0) {
    setInterval(fetchData, props.refreshInterval);
  }
});

// 
watch(chartData, () => {
  updateChart();
}, { deep: true });

// Props
watch(() => props.stationId, () => {
  fetchData();
});

// 
import { onUnmounted } from 'vue';

onUnmounted(() => {
  window.removeEventListener('resize', handleResize);
  if (chart.value) {
    chart.value.dispose();
  }
});
</script>

<style scoped lang="scss">
.water-level-chart {
  background: white;
  border-radius: 8px;
  padding: 20px;
  box-shadow: 0 2px 8px rgba(0, 0, 0, 0.1);

  .chart-header {
    display: flex;
    justify-content: space-between;
    align-items: center;
    margin-bottom: 20px;

    h3 {
      margin: 0;
      color: #2c3e50;
    }

    .chart-controls {
      select {
        padding: 8px 12px;
        border: 1px solid #d9d9d9;
        border-radius: 4px;
        background: white;
        font-size: 14px;

        &:focus {
          outline: none;
          border-color: #1890ff;
          box-shadow: 0 0 0 2px rgba(24, 144, 255, 0.2);
        }
      }
    }
  }

  .chart-container {
    height: v-bind('props.height + "px"');
    width: 100%;
  }
}
</style>
\end{lstlisting}

TypeScript

\begin{itemize}
\item
  \important{}
\item
  \important{}IDE
\item
  \important{}
\item
  \important{}
\end{itemize}

\hypertarget{section-46}{%
\subsection{4.6.6}\label{section-46}}

\hypertarget{section-47}{%
\subsubsection{}\label{section-47}}

\begin{longtable}[]{@{}ll@{}}
\toprule\noalign{}
\i & mportant\{1. } \\
\midrule\noalign{}
\endhead
\bottomrule\noalign{}
\endlastfoot
- & \important{2. Vue CLI} \\
- & Vue CLI \\
- & - Vue CLI \\
- & \important{3. } \\
- & WebpackVite \\
\end{longtable}

\begin{longtable}[]{@{}ll@{}}
\toprule\noalign{}
\i & mportant\{4. } \\
\midrule\noalign{}
\endhead
\bottomrule\noalign{}
\endlastfoot
- & \important{5. } \\
- & ESLintPrettierTypeScript \\
- & Git Hooks \\
- & \#\#\# \\
\end{longtable}

\important{}

\important{}ESLintTypeScript

\important{}GIS

\important{}

\hypertarget{section-48}{%
\subsubsection{}\label{section-48}}

\important{}Vite

\important{}

\important{}

\important{}

\hypertarget{section-49}{%
\subsubsection{}\label{section-49}}

\important{1. }

\important{2. }

\important{3. }

\important{4. }

{[}\{``content'':
``\u9605\u8bfbchapter04.md\u5927\u7eb2\uff0c\u4e86\u89e34.6\u8282\u7684\u8981\u6c42\u548c\u7ed3\u6784'',
``id'': ``read-outline'', ``status'': ``completed''}, \{``content'':
``\u5206\u67904.6\u8282\u9700\u8981\u64b0\u5199\u7684\u5185\u5bb9\u548c\u6df1\u5ea6\u8981\u6c42'',
``id'': ``analyze-requirements'', ``status'': ``completed''},
\{``content'':
``\u64b0\u51994.6\u8282\u5185\u5bb9\uff0c\u5305\u542b\u8be6\u7ec6\u7406\u8bba\u9610\u8ff0\u548c\u5b9e\u8df5\u6848\u4f8b'',
``id'': ``write-content'', ``status'': ``completed''}, \{``content'':
``\u68c0\u67e5\u5185\u5bb9\u4e0e\u524d\u9762\u7ae0\u8282\u7684\u4e00\u81f4\u6027\u548c\u8fde\u8d2f\u6027'',
``id'': ``review-consistency'', ``status'': ``in\_progress''}{]}

\hypertarget{section-50}{%
\section{4.7}\label{section-50}}

FTPCI/CDCDN

!!! info ``\,''

\begin{verbatim}
\important{}\important{}\important{}
\end{verbatim}

\hypertarget{section-51}{%
\subsection{4.7.1}\label{section-51}}

\hypertarget{section-52}{%
\subsubsection{}\label{section-52}}

Tree Shaking50-80\%60-90\%

\hypertarget{section-53}{%
\subsubsection{}\label{section-53}}

\hypertarget{section-54}{%
\paragraph{}\label{section-54}}

\important{Code Splitting}

\begin{itemize}
\item
  \important{}
\item
  \important{}
\item
  \important{}ECharts
\item
  \important{}
\end{itemize}

\important{Tree Shaking}Element Plus

\hypertarget{section-55}{%
\paragraph{}\label{section-55}}

\important{}JavaScriptCSSHTMLTerser

\important{}

\begin{itemize}
\item
  WebPAVIF
\item
  \begin{itemize}
  \tightlist
  \item
  \end{itemize}
\end{itemize}

\hypertarget{section-56}{%
\paragraph{}\label{section-56}}

\begin{lstlisting}[style=javascript]
// vue.config.js - 
module.exports = {
  // 
  configureWebpack: config => {
    if (process.env.NODE_ENV === 'production') {
      // Gzip
      config.plugins.push(new CompressionWebpackPlugin({
        algorithm: 'gzip',
        test: /\.(js|css|html|svg)$/
      }))
      
      // 
      config.optimization.splitChunks = {
        chunks: 'all',
        cacheGroups: {
          vendor: {
            name: 'chunk-vendors',
            test: /[\\/]node_modules[\\/]/,
            priority: 10
          },
          // 
          waterComponents: {
            name: 'chunk-water',
            test: /[\\/]src[\\/]components[\\/]water/,
            priority: 20
          }
        }
      }
    }
  }
}
\end{lstlisting}

\hypertarget{section-57}{%
\subsubsection{}\label{section-57}}

\hypertarget{tree-shaking}{%
\paragraph{Tree Shaking}\label{tree-shaking}}

\important{Tree Shaking}ES6''\,``------

Tree ShakingUITree Shaking

\begin{itemize}
\tightlist
\item
  \important{UI}Element PlusAnt Design
\item
  \important{}lodash
\item
  \important{}ECharts
\end{itemize}

\hypertarget{section-58}{%
\paragraph{}\label{section-58}}

\begin{lstlisting}[style=javascript]
// 
const config = {
  development: {
    apiUrl: 'http://localhost:3000/api',
    mapServiceKey: 'dev-key-123'
  },
  production: {
    apiUrl: 'https://api.water-monitor.gov.cn',
    mapServiceKey: 'prod-key-456'
  }
}
\end{lstlisting}

\hypertarget{section-59}{%
\paragraph{}\label{section-59}}

\important{}

\begin{itemize}
\item
  \important{Bundle Analyzer}
\item
  \important{}
\item
  \important{}
\end{itemize}

\important{}

\hypertarget{cdn}{%
\subsection{4.7.2 CDN}\label{cdn}}

CDN

\hypertarget{cdn-1}{%
\subsubsection{CDN}\label{cdn-1}}

\important{CDNContent Delivery Network}

\hypertarget{cdn-2}{%
\paragraph{CDN}\label{cdn-2}}

\begin{enumerate}
\def\labelenumi{\arabic{enumi}.}
\item
  \important{}
\item
  \important{}
\item
  \important{}
\item
  \important{}
\end{enumerate}

\hypertarget{section-60}{%
\paragraph{}\label{section-60}}

\begin{itemize}
\item
  \important{}
\item
  \important{}
\item
  \important{}
\item
  \important{}
\end{itemize}

\hypertarget{section-61}{%
\subsubsection{}\label{section-61}}

\begin{longtable}[]{@{}llll@{}}
\toprule\noalign{}
\endhead
\bottomrule\noalign{}
\endlastfoot
\important{HTML} & & 1 & \\
\important{JavaScript/CSS} & & 1 & hash \\
\important{} & & 6 & \\
\important{} & & 1+ & \\
\end{longtable}

\hypertarget{nginx}{%
\subsubsection{Nginx}\label{nginx}}

\important{}

\begin{lstlisting}
server {
    listen 443 ssl http2;
    server_name watermonitor.gov.cn;
    
    # SSL
    ssl_certificate /path/to/cert.pem;
    ssl_certificate_key /path/to/key.pem;
    
    # 
    root /var/www/water-monitor/dist;
    index index.html;
    
    # Gzip
    gzip on;
    gzip_types text/css application/javascript image/svg+xml;
    
    # 
    location ~* \.(js|css|png|jpg|svg)$ {
        expires 1y;
        add_header Cache-Control "public, immutable";
    }
    
    # SPA
    location / {
        try_files $uri $uri/ /index.html;
    }
}
\end{lstlisting}

\hypertarget{cdn-3}{%
\subsubsection{CDN}\label{cdn-3}}

\hypertarget{cdn-4}{%
\paragraph{CDN}\label{cdn-4}}

CDN

\important{1. }

\begin{itemize}
\item
  \begin{itemize}
  \tightlist
  \item
    URL
  \end{itemize}
\item
  \important{2. }
\item
  \begin{itemize}
  \tightlist
  \item
    IP
  \end{itemize}
\item
  HTTPS
\end{itemize}

\important{3. }

\begin{itemize}
\item
  \begin{itemize}
  \tightlist
  \item
    HTTP/2
  \end{itemize}
\item
  CDN
\end{itemize}

\hypertarget{cdn-5}{%
\paragraph{CDN}\label{cdn-5}}

CDN

\begin{lstlisting}[style=javascript]
// CDN
class CDNManager {
  constructor() {
    this.providers = [
      { name: 'primary', baseUrl: 'https://cdn1.water-monitor.com' },
      { name: 'backup', baseUrl: 'https://cdn2.water-monitor.com' }
    ]
    this.currentProvider = this.providers[0]
  }
  
  // 
  async checkHealth() {
    try {
      const response = await fetch(this.currentProvider.baseUrl + '/health')
      return response.ok
    } catch {
      return false
    }
  }
}
\end{lstlisting}

\important{} CDNCDN60-80\%

\hypertarget{section-62}{%
\subsection{4.7.3}\label{section-62}}

\hypertarget{section-63}{%
\subsubsection{}\label{section-63}}

\hypertarget{section-64}{%
\paragraph{}\label{section-64}}

\important{1. }

\begin{itemize}
\item
  \important{}
\item
  \important{}JavaScript
\item
  \important{}
\end{itemize}

\important{2. }

\begin{itemize}
\item
  \important{}
\item
  \important{}
\item
  \important{}
\end{itemize}

\hypertarget{core-web-vitals}{%
\paragraph{Core Web Vitals}\label{core-web-vitals}}

GoogleCore Web Vitals

\begin{itemize}
\item
  \important{LCP (Largest Contentful Paint)}
\item
  \important{FID (First Input Delay)}
\item
  \important{CLS (Cumulative Layout Shift)}
\end{itemize}

\hypertarget{section-65}{%
\subsubsection{}\label{section-65}}

\hypertarget{section-66}{%
\paragraph{}\label{section-66}}

\begin{itemize}
\item
  \important{WebSocket}
\item
  \important{}
\item
  \important{}
\item
  \important{}GIS
\end{itemize}

\hypertarget{section-67}{%
\paragraph{}\label{section-67}}

\begin{lstlisting}[style=javascript]
// 
class UserBehaviorMonitor {
  trackTaskCompletion(taskName) {
    const startTime = performance.now()
    
    return {
      complete: () => {
        const duration = performance.now() - startTime
        this.reportTaskCompletion(taskName, duration, true)
      },
      fail: (reason) => {
        this.reportTaskCompletion(taskName, 0, false, reason)
      }
    }
  }
}
\end{lstlisting}

\hypertarget{section-68}{%
\paragraph{}\label{section-68}}

\begin{itemize}
\item
  \important{}
\item
  \important{}
\item
  \important{}
\end{itemize}

\hypertarget{section-69}{%
\subsubsection{}\label{section-69}}

\hypertarget{section-70}{%
\paragraph{}\label{section-70}}

\important{}\important{}

\important{}

\important{1. }
\important{}\important{}\important{}

\important{2. }
\important{}

\important{3. }
\important{}\important{}

\hypertarget{section-71}{%
\paragraph{}\label{section-71}}

\important{}\important{}

\important{}

\important{}

\begin{lstlisting}[style=javascript]
// 
class SmartAlertSystem {
  checkPerformanceAnomaly(currentMetrics, historicalData) {
    const threshold = this.calculateDynamicThreshold(historicalData)
    const anomalyScore = this.detectAnomaly(currentMetrics, threshold)
    
    if (anomalyScore > 0.8) {
      this.triggerAlert('performance', currentMetrics)
    }
  }
}
\end{lstlisting}

\important{} \important{}

\hypertarget{section-72}{%
\subsection{4.7.4}\label{section-72}}

\hypertarget{section-73}{%
\subsubsection{}\label{section-73}}

\hypertarget{section-74}{%
\paragraph{}\label{section-74}}

\important{}\important{}\important{}Web

\important{1. }
\important{}\important{}

\important{2. }
\important{}\important{}

\important{3. }
\important{}\important{}

\hypertarget{section-75}{%
\paragraph{}\label{section-75}}

\important{}

\important{}Git

\important{}CI/CD

\important{}

\hypertarget{section-76}{%
\subsubsection{}\label{section-76}}

\hypertarget{section-77}{%
\paragraph{}\label{section-77}}

\important{Semantic Versioning}\important{MAJOR.MINOR.PATCH}

\begin{itemize}
\item
  \important{MAJOR}API
\item
  \important{MINOR}
\item
  \important{PATCH}
\end{itemize}

\important{}

\hypertarget{section-78}{%
\paragraph{}\label{section-78}}

\important{Patch}\important{}24-48

\important{Minor}\important{}

\important{Major}\important{}1-2

\begin{lstlisting}[style=javascript]
// 
const versionConfig = {
  current: "2.1.3",
  nextPatch: "2.1.4",    // 
  nextMinor: "2.2.0",    //   
  nextMajor: "3.0.0"     // 
}
\end{lstlisting}

\hypertarget{git}{%
\subsubsection{Git}\label{git}}

\hypertarget{section-79}{%
\paragraph{}\label{section-79}}

\important{Git Flow}\important{}\important{}

\important{}

\important{1. }

maindevelopfeature

\important{2. }

\important{3. }

\hypertarget{section-80}{%
\paragraph{}\label{section-80}}

\important{}

\important{}

\important{}``\,''

\important{}

\hypertarget{section-81}{%
\subsubsection{}\label{section-81}}

\hypertarget{cicd}{%
\paragraph{CI/CD}\label{cicd}}

\important{CI}\important{}

\important{CD}\important{}

\hypertarget{section-82}{%
\paragraph{}\label{section-82}}

CI/CD

\important{1. }

CI/CD

\important{2. }

\important{3. }

\begin{lstlisting}
# CI
name: CI/CD
on: [push, pull_request]
jobs:
  build-and-test:
    runs-on: ubuntu-latest
    steps:
      - name: 
        run: npm run lint
      - name:   
        run: npm test
      - name: 
        run: npm run build
\end{lstlisting}

\hypertarget{section-83}{%
\subsubsection{}\label{section-83}}

\hypertarget{section-84}{%
\paragraph{}\label{section-84}}

\important{}\important{}

\important{}

\important{}

\important{}

\important{}

\hypertarget{section-85}{%
\paragraph{}\label{section-85}}

\important{}

\important{}

\important{}

\important{}

\important{}

\hypertarget{section-86}{%
\subsection{4.7.5}\label{section-86}}

\hypertarget{section-87}{%
\subsubsection{}\label{section-87}}

\important{1. }

Tree Shaking50-80\%60-90\%

\begin{itemize}
\item
  \begin{itemize}
  \tightlist
  \item
  \end{itemize}
\item
  \begin{itemize}
  \tightlist
  \item
    Source Map
  \end{itemize}
\end{itemize}

\important{2. CDN}

CDNCDN60-80\%

\begin{itemize}
\item
  \begin{itemize}
  \tightlist
  \item
    Nginx
  \end{itemize}
\item
  CDN
\item
  CDN
\end{itemize}

\important{3. }

Core Web Vitals

\begin{itemize}
\item
  \begin{itemize}
  \tightlist
  \item
  \end{itemize}
\item
  \begin{itemize}
  \tightlist
  \item
  \end{itemize}
\end{itemize}

\important{4. }

Git FlowCI/CD

\begin{itemize}
\item
  \begin{itemize}
  \tightlist
  \item
  \end{itemize}
\item
  \begin{itemize}
  \tightlist
  \item
  \end{itemize}
\end{itemize}

\hypertarget{section-88}{%
\subsubsection{}\label{section-88}}

\important{}

\important{}WebSocket

\important{}HTTPS

\important{}CDN

\hypertarget{section-89}{%
\subsubsection{}\label{section-89}}

\important{}5G

\important{}DockerKubernetes

\important{Serverless}

\important{}AI

\hypertarget{section-90}{%
\subsubsection{}\label{section-90}}

\important{1. }

CDN CI/CD

\important{2. }

\begin{itemize}
\item
  \begin{itemize}
  \tightlist
  \item
  \end{itemize}
\item
  \begin{itemize}
  \tightlist
  \item
  \end{itemize}
\end{itemize}

\important{3. }

\begin{itemize}
\item
  \begin{itemize}
  \tightlist
  \item
  \end{itemize}
\item
  \begin{itemize}
  \tightlist
  \item
  \end{itemize}
\end{itemize}

\hypertarget{section-91}{%
\subsubsection{}\label{section-91}}

Web \important{} GitCI/CD

\hypertarget{section-92}{%
\subsection{4.7.5}\label{section-92}}

\hypertarget{section-93}{%
\subsubsection{}\label{section-93}}

\important{1. }

\important{}

\important{}\important{}Tree ShakingGIS

\important{2. CDN}

CDNCDN

\begin{itemize}
\item
  \important{}
\item
  \important{}
\item
  \important{}
\end{itemize}

\important{3. }

\important{}\important{}Core Web Vitals

\important{}\important{}

\important{4. }

\important{}Git FlowCI/CD

\hypertarget{section-94}{%
\subsubsection{}\label{section-94}}

\important{}

\begin{itemize}
\item
  \important{}
\item
  \important{}
\item
  \important{}
\end{itemize}

\important{}

\begin{itemize}
\item
  \important{}
\item
  \important{}
\item
  \important{}
\end{itemize}

\important{}

\begin{itemize}
\item
  \important{CDN}
\item
  \important{}
\item
  \important{}
\end{itemize}

\hypertarget{section-95}{%
\subsubsection{}\label{section-95}}

\important{}

5G

\begin{itemize}
\item
  \begin{itemize}
  \tightlist
  \item
  \end{itemize}
\item
  \important{}
\end{itemize}

DockerKubernetes

\begin{itemize}
\item
  \important{}
\item
  \important{}
\item
  \important{}
\end{itemize}

\important{}

AI

\begin{itemize}
\item
  \important{}
\item
  \important{}
\item
  \important{}
\end{itemize}

\hypertarget{section-96}{%
\subsubsection{}\label{section-96}}

\important{}

\important{}

\begin{itemize}
\item
  \begin{itemize}
  \tightlist
  \item
    CDN
  \end{itemize}
\item
  \important{}
\item
  CI/CD
\item
  \begin{itemize}
  \tightlist
  \item
  \end{itemize}
\end{itemize}

\begin{longtable}[]{@{}ll@{}}
\toprule\noalign{}
\i & mportant\{} \\
\midrule\noalign{}
\endhead
\bottomrule\noalign{}
\endlastfoot
- & \important{} \\
\end{longtable}

\begin{itemize}
\item
  \important{}
\item
  \important{}
\item
  \important{}
\item
  \important{}
\end{itemize}

\important{}

\begin{itemize}
\item
  \important{}
\item
  \important{}
\item
  \important{}
\item
  \important{}
\end{itemize}

\hypertarget{section-97}{%
\subsubsection{}\label{section-97}}

\important{}\important{}

Web\important{}\important{}

\important{}